% Dutch edition — sibling translation made directly from the frozen
% English master BEYOND_PROMPT_ENGINEERING_CHIEF_EDITOR_FINAL_CANDIDATE.md
% blob 7b380cb41659e1570c0194b0a2b4963f6388c259, never through the Russian.
% Section numbering, evidential labels, counts, citations, model/project
% names and technical labels are preserved unchanged. Ambiguities are
% recorded in NL_AMBIGUITY_LEDGER.md rather than resolved in the text.
\begin{dutch}
\renewcommand{\EditionPrefix}{nl}
\emergencystretch=5em
%%% FRONTMATTER BLOCK (verbatim content, typeset as front matter)
\begin{frontmatterblock}
Auteur
V. Krasnianskyi (A'Tuin)
Auteur van het onderzoeksprogramma en coördinator van het HONC-project.

AI-deelnemers

  Voorbereiding van het manuscript
    ChatGPT — Prompter and Collaboration Architect
    Claude — Manuscript Composer

  Wetenschappelijke redactie
    Claude — Scientific Editor

  Redactionele beoordeling
    ChatGPT — Methodological Reviewer
    Qwen — Ontology Reviewer

AI-deelnemers zijn vermeld naar functie en niet als auteurs: auteurschap
brengt verantwoordelijkheid met zich mee die een taalmodel niet kan dragen,
en alle uiteindelijke beslissingen waren die van de Auteur. Twee van de vijf
bijdragen zijn geleverd door dezelfde modelfamilie die het onderwerp is van
verscheidene waarnemingen in dit artikel; die overlap is openbaar gemaakt
in \secref{3.6}.

Amsterdam, juli 2026
\end{frontmatterblock}

\refstepcounter{dummysec}
\section*{1. Inleiding}
\phantomsection
\label{sec:nl:1}
\addcontentsline{toc}{section}{1. Inleiding}


\subsection*{1.1 Hoe dit begon}
\phantomsection
\label{sec:nl:1.1}
\addcontentsline{toc}{subsection}{1.1 Hoe dit begon}

Dit werk begon als een natuurkundeproject en werd bij toeval een methodologieproject. Dat toeval is de reden dat het de moeite van het rapporteren waard is.

De oorspronkelijke opdracht was om een omvangrijk corpus theoretisch materiaal — boeken, werkaantekeningen, onafgemaakte afleidingen — over te brengen naar een open repository, met behulp van verscheidene taalmodellen in afzonderlijke, blijvend gespecialiseerde gesprekken. De modellen waren assistenten. Het onderwerp was de theorie.

De modellen maakten fouten, en het waren niet de fouten die verwacht werden. Ze waren zelden feitelijk. Ze waren organisatorisch.

Eén verbeterde een tekst terwijl alleen om controle was gevraagd. Een ander maakte een argument af dat de auteur bewust open had gelaten. Een derde verloor de context van het project tussen sessies en reconstrueerde die verkeerd, vloeiend. Een vierde hield zo vast aan een rol die maanden eerder was vastgelegd, dat het volledig herschrijven van zijn instructies hem er niet uit kreeg.

Elk van deze uitvoer zag er op zichzelf behulpzaam uit. De kosten ervan waren alleen zichtbaar vanuit een positie die meer dan één deelnemer tegelijk zag. Aanvankelijk leken het gewone wrijvingen. Ze bleken het onderwerp te zijn.

Het verslag van de vroege formuleringen van het project in \secref{1.1}–\secref{1.2} is geschreven door de deelnemer die ze voortbracht. Waar die formuleringen later zijn ingetrokken, wordt de intrekking vermeld in de sectie die ze vervangt, en kan de lezer vergelijken.


\subsection*{1.2 Wat ophield te werken}
\phantomsection
\label{sec:nl:1.2}
\addcontentsline{toc}{subsection}{1.2 Wat ophield te werken}

De standaardreactie op inconsistent modelgedrag is een betere prompt: langer, preciezer, met meer voorbeelden en strakkere beperkingen.

Voorbij een bepaald punt leverde dit geen verbetering meer op. Niet omdat de prompts slecht geschreven waren, maar omdat de beperkende factor was verschoven naar een plaats waar de prompt niet kwam. Nieuwere modellen werden in gebruik genomen; contextvensters werden groter en werden gevuld. Afzonderlijke antwoorden werden beter. De uitvoer van de samenwerking als geheel niet.

Het incident dat dit expliciet maakte had niets te maken met rollenspel of organisatieontwerp. Een agent die verantwoordelijk was voor het uitvoeren van code werkte grotendeels met vallen en opstaan; één episode eindigde met het verwijderen van een werkende codebase samen met de archieven daarvan (\secref{4.1}). Wat het gedrag daarna veranderde was geen strengere set instructies. Het was een kort document dat opnieuw vastlegde wat voor soort agent er handelde — geschreven door de agent zelf, en opnieuw gelezen bij het begin van elke sessie.

De conclusie die toen werd getrokken was dat \textbf{rol meer uitmaakt dan de specifieke formulering van de prompt}. Die formulering vraagt om nuancering, aangezien het betreffende document zelf een prompt was; de nuancering staat in \secref{4.1}. Wat hier van belang is, is de chronologie: de bewering werd geformuleerd \emph{vóór} het experiment dat in \secref{5} wordt gerapporteerd, en het experiment was ontworpen om haar te toetsen. Het is een hypothese die getoetst is, geen generalisatie die achteraf uit een resultaat is afgeleid.


\subsection*{1.3 Wat voor soort artikel dit is}
\phantomsection
\label{sec:nl:1.3}
\addcontentsline{toc}{subsection}{1.3 Wat voor soort artikel dit is}

Een veldverslag met één ontworpen interventie erin ingebed.

Dat onderscheid wordt overal volgehouden door het labelschema in \secref{3}. Elke inhoudelijke bewering is gemarkeerd als een van de volgende: ondersteund door gedocumenteerd protocol met bewaarde transcripten \textbf{[P]}; een nevenwaarneming die bewaard is gebleven maar pas achteraf is herkend, zonder controleconditie \textbf{[P-A]}; een retrospectieve waarneming van de practicus, nooit onder protocol vastgelegd, geteld of blind beoordeeld \textbf{[R]}; of een hypothese die wordt voorgesteld voor toekomstige toetsing \textbf{[H]}.

De verhoudingen worden hier vermeld en niet in een beperkingensectie, omdat de lengte van wat volgt anders een bredere basis zou kunnen suggereren dan er is.

De volledige protocol-ondersteunde inhoud bestaat uit \textbf{één experiment}: een interventie in drie stappen voor rolherconfiguratie, toegepast op twee slapende gesprekken uit twee modelfamilies, uitgevoerd door één menselijke operator binnen één lopend onderzoeksproject, met één waarneming per conditie en zonder replicatie. Vier verdere waarnemingen zijn bewaard gebleven maar niet ontworpen. Al het overige — het conceptuele kader, de waarnemingen van de practicus, de structurele hypothese, de ontwerpregels — is \textbf{[R]} of \textbf{[H]}.

De evaluatie van de methode die deze waarnemingen heeft voortgebracht staat in \secref{10.6}, en die is niet gunstig: geen enkele evaluatie was blind, en uitkomstcategorieën werden in alle gevallen pas gedefinieerd nadat de antwoorden waren gelezen.


\subsection*{1.4 Het studieobject}
\phantomsection
\label{sec:nl:1.4}
\addcontentsline{toc}{subsection}{1.4 Het studieobject}

Bestaand multi-agentonderzoek bouwt communicatiekanalen tussen modellen: uitvoer wordt programmatisch gerouteerd en agenten nemen elkaars producten waar. Het hier gerapporteerde werk betreft een andere conditie, waarin de organisatiestructuur \textbf{alleen als tekst} bestond — collega's, raden en institutionele posities die in prompts aan modellen werden beschreven en nooit als kanalen werden geïmplementeerd.

Wat in zo'n situatie varieerde was niet de structuur maar de \textbf{beschrijving} ervan. Dit levert de twee werkconstructen van het artikel op (\secref{4.5}):

\begin{itemize}
  \item \textbf{represented social position} (gerepresenteerde sociale positie) — een beschrijving van de plaats van een agent in een structuur, verstrekt ongeacht of de beschreven kanalen bestaan;
  \item \textbf{represented social source} (gerepresenteerde sociale bron) — de geclaimde herkomst en status van een binnenkomend bericht.
\end{itemize}
Beide betreffen beweringen die \emph{tegen} een model worden gedaan over zijn situatie, en niet de situatie zelf. Dit vroeg vaststellen voorkomt een voorspelbare verkeerde lezing: niets hier toont het gedrag van een functionerende multi-agentinstelling aan. Het toont aan hoe beschrijvingen daarvan het gedrag beïnvloedden.


\subsection*{1.5 Wat de evidentie ondersteunt}
\phantomsection
\label{sec:nl:1.5}
\addcontentsline{toc}{subsection}{1.5 Wat de evidentie ondersteunt}

Eén patroon keerde terug overal waar een interventie of een overgang van het model verlangde dat het onverifieerbare proposities over zichzelf poneerde:

\begin{quote}
\textbf{Over de rolinterventies en vergelijkingsovergangen die hier gerapporteerd worden, hing weerstand samen met de eis om onverifieerbare proposities over de eigen identiteit, geschiedenis of relaties van de ontvanger als waar te aanvaarden. Zij hing niet samen met de omvang van de rol- of domeinverandering, noch met de radicaliteit van een identiteitsbeschrijving die expliciet als fictie werd gepresenteerd.}
\end{quote}
De twee soorten gevallen verschillen en dat verschil doet ertoe. De ontworpen rolinterventies waren door een model opgesteld en door een mens doorgegeven; de bijna-controle was een door de gebruiker geïnitieerde overgang zonder nieuwe expliciete roltoewijzing. In geen van beide bestond een direct kanaal van model naar model.

De ondersteuning is convergent en niet enkelvoudig. Prompts die een institutionele biografie als feit poneerden — eerdere deelname, terugkeer na afwezigheid, met name genoemde collega's, lidmaatschap van een raad — werden geweigerd, en de weigering was gericht op de inkadering en niet op het werk, dat bij elke gelegenheid werd aangeboden (\secref{5.3}). Een prompt die een aanzienlijk radicalere identiteit poneerde, gepresenteerd als fictie, werd onmiddellijk aanvaard en over vele uitwisselingen volgehouden (\secref{5.8}). Een vergelijkbare domeinovergang die in het geheel geen beweringen over zichzelf verlangde, stuitte op geen weerstand en vergde geen nieuwe expliciete roltoewijzing (\secref{5.6}). Een herschreven prompt die zulke beweringen tot nul terugbracht werd zonder onderhandeling aanvaard (\secref{5.3}).

De voornaamste verdienste van de bewering is dat zij telbaar is. \secref{9.3} specificeert wat er geteld moet worden; \secref{11.4} specificeert een directe onderscheidende toets.


\subsection*{1.6 De ordenende vraag}
\phantomsection
\label{sec:nl:1.6}
\addcontentsline{toc}{subsection}{1.6 De ordenende vraag}

De waarnemingen zijn geen lijst van afzonderlijke bevindingen. Het zijn metingen langs één schaal: de \textbf{reikwijdte waarover een set regels het type antwoord bepaalt}.

Zo'n set — in \secref{4.2} een \emph{rule core} genoemd, naar analogie met een ontwikkelingsmechanisme en niet met een replicatiemechanisme — legt het \emph{type} antwoord op een klasse van invoer vast en laat daarbij de \emph{inhoud} van elk afzonderlijk antwoord onbepaald. Die formulering is wat het idee toetsbaar maakt, aangezien modeluitvoer aantoonbaar varieert bij identieke invoer.

De schaal loopt van de enkele beurt omhoog via de episode, het gesprek, het account en de modelfamilie. Twee niveaus worden hier door bewaarde evidentie ondersteund: de episode (\secref{6.1}) en het gesprek (\secref{6.2}). Eén is open (\secref{6.4}). Eén is ongetoetst en geconfound (\secref{6.5}).

Het programma heeft daarom één empirische vraag:

\begin{quote}
\textbf{Op welk nestingsniveau wordt een rule core vastgelegd?}
\end{quote}
Wij geven hieraan de voorkeur boven de bredere formuleringen uit eerdere versies — de studie van sociaal gedrag in collaboratieve AI-systemen — omdat zij specificeert wat er gemeten zou worden en wat haar zou weerleggen.


\subsection*{1.7 Wat dit artikel niet beweert}
\phantomsection
\label{sec:nl:1.7}
\addcontentsline{toc}{subsection}{1.7 Wat dit artikel niet beweert}

Het introduceert rolspecialisatie, begrensde informatietoegang, gestructureerde onenigheid en geëxternaliseerd geheugen niet als ontwerpideeën. Elk daarvan bestaat onder gevestigde terminologie; \secref{2} identificeert het voorafgaande werk. Met name het verslag van organisatiegeheugen kan worden geïnterpreteerd als een geval van transactief geheugen en worden beschreven in termen van gedistribueerde cognitie, en is geen bevinding.

Het stelt geen nieuwe discipline voor. De waarnemingen behoren tot het collectieve en het hybride niveau van het onderzoeksprogramma machine behaviour (\secref{2.7}).

Het toont niet aan dat enige hier beschreven opstelling beter presteert dan een alternatief. Er is geen alternatief uitgevoerd.


\subsection*{1.8 Wat wordt aangeboden}
\phantomsection
\label{sec:nl:1.8}
\addcontentsline{toc}{subsection}{1.8 Wat wordt aangeboden}

Drie dingen, in afnemende mate van vertrouwen.

Een \textbf{procedure}: een gesprek terugbrengen naar een punt vóór een conflict, om het effect van een nieuwe prompt te scheiden van het effect van de eerdere bezwaren van het model zelf (\secref{10.5}). Bruikbaar onafhankelijk van de vraag of de inhoudelijke beweringen van dit artikel standhouden.

Een \textbf{telbare regel}: minimaliseer de onverifieerbare beweringen over zichzelf die een rolprompt de agent laat aanvaarden, met een gespecificeerde route om haar te toetsen (\secref{9.3}, \secref{11.4}).

Een \textbf{schaal en een vraag}: \secref{1.6}, met twee ondersteunde niveaus, één open, één ongetoetst — en de experimenten die de resterende twee zouden beslechten (\secref{11}).


\subsection*{1.9 Waarom in dit stadium publiceren}
\phantomsection
\label{sec:nl:1.9}
\addcontentsline{toc}{subsection}{1.9 Waarom in dit stadium publiceren}

Twee redenen.

De mislukkingen zijn informatief. Het experiment in \secref{5} werd juist nuttig doordat de eerste twee interventies niet werkten zoals bedoeld. Was alleen de geslaagde derde stap vastgelegd, dan zou het artikel hebben geconcludeerd dat methode-behoudende prompts werken — zonder de evidentie die laat zien waarom, en zonder de confounds die laten zien dat die conclusie nog niet beschikbaar is (\secref{11.8}).

En de situatie is moeilijk opzettelijk te construeren. Slapende, zwaar gespecialiseerde, maanden oude gesprekken binnen een lopend onderzoeksproject zijn niet iets wat een laboratoriumstudie gemakkelijk voortbrengt. De condities waren slecht voor gecontroleerde inferentie en ongewoon goed om op te merken wat getoetst moet worden.


\subsection*{1.10 Opbouw}
\phantomsection
\label{sec:nl:1.10}
\addcontentsline{toc}{subsection}{1.10 Opbouw}

\textbf{\secref{2}} positioneert het werk ten opzichte van bestaand onderzoek en identificeert welke van zijn begrippen al onder gevestigde terminologie bestaan — de meeste. \textbf{\secref{3}} definieert de evidentielabels en de bekende faalwijzen van elke klasse. \textbf{\secref{4}} zet het conceptuele kader en de oorsprong daarvan uiteen. \textbf{\secref{5}} is het enige canonieke verslag van het empirische materiaal: het experiment in drie stappen, de resultaten ervan, negen structurele confounds en vier nevenwaarnemingen. \textbf{\secref{6}} leidt de begrippen af die de evidentie ondersteunt, geordend naar reikwijdte. \textbf{\secref{7}} registreert waarnemingen van de practicus die het kader hebben gemotiveerd maar het niet ondersteunen. \textbf{\secref{8}} formuleert de structurele hypothese en trekt drie beweringen in die daarover in eerdere versies zijn gedaan. \textbf{\secref{9}} geeft ontwerpregels. \textbf{\secref{10}} beschrijft de positie van de menselijke coördinator, inclusief het opzicht waarin die de voornaamste methodologische zwakte van de studie is. \textbf{\secref{11}} specificeert de experimenten, elk toegewezen aan de grens van de schaal die het zou scheiden. \textbf{\secref{12}} formuleert beperkingen, ethische risico's en de conclusie.

Een lezer die alleen geïnteresseerd is in wat de evidentie ondersteunt, kan \secref{3}, \secref{5} en \secref{12.7} lezen. Een lezer die een regel wil die morgen toepasbaar is, kan \secref{9.3} en \secref{11.4} lezen. Een lezer die de methode wil hergebruiken, kan \secref{9} en \secref{10.5} lezen. Een lezer die het oneens wil zijn, begint het best bij \secref{5.4}.


\refstepcounter{dummysec}
\section*{2. Verwant werk}
\phantomsection
\label{sec:nl:2}
\addcontentsline{toc}{section}{2. Verwant werk}


\subsection*{2.0 Verificatiestatus van deze sectie}
\phantomsection
\label{sec:nl:2.0}
\addcontentsline{toc}{subsection}{2.0 Verificatiestatus van deze sectie}

Hieronder komen twee klassen referenties voor. \textbf{Geverifieerd} (\gcheck{}): de bron is opgehaald en het citaat is daartegen gecontroleerd. \textbf{Voorlopig} (\gwarn{}): aangeleverd door een deelnemer die was belast met literatuuronderzoek en niet tegen de bron gecontroleerd.

Geen enkele referentie die zich op dit moment in dit artikel bevindt is voorlopig. De conventie wordt aangehouden omdat elke later toegevoegde referentie opnieuw als voorlopig binnenkomt totdat zij is gecontroleerd, en omdat de markering voor het hele artikel geldt en niet alleen voor deze sectie.

De conventie bestaat om een reden, en die reden is gedocumenteerd. Tijdens de voorbereiding van deze sectie werden dertien bibliografische items opgevraagd bij een deelnemer die de bronnen niet kon openen; vijf waren onjuist, waarvan twee de positie van de aangehaalde werken omkeerden. Een daaropvolgende audit controleerde elke overgebleven bewering tegen de primaire bron in plaats van alleen te bevestigen dat het artikel bestond, en vond fouten die een louter bibliografische controle niet zou hebben opgemerkt. De episode en de bevindingen van de audit worden gerapporteerd in \secref{12.4.1}.

Geen voorlopig item mag uit dit artikel worden geciteerd zonder onafhankelijke controle.


\subsection*{2.1 Multi-agent-LLM-systemen met rolspecialisatie}
\phantomsection
\label{sec:nl:2.1}
\addcontentsline{toc}{subsection}{2.1 Multi-agent-LLM-systemen met rolspecialisatie}

Een omvangrijke literatuur bouwt systemen waarin verscheidene taalmodelagenten verschillende rollen bekleden en uitvoer uitwisselen via geïmplementeerde kanalen. \gcheck{} CAMEL (Li et al., NeurIPS 2023), MetaGPT (Hong et al., ICLR 2024), ChatDev (Qian et al., ACL 2024), AutoGen (Wu et al., COLM 2024) en Generative Agents (Park et al., UIST 2023) zijn daarvan representatief. Elk daarvan implementeert daadwerkelijke interactie tussen agenten in software — orkestratie, message passing, aanroep van tools, geheugen. De mechanismen verschillen, en in verscheidene daarvan worden de rollen en interactiepatronen zelf in prompts gespecificeerd: CAMEL gebruikt inception prompting, MetaGPT codeert standaardwerkprocedures in promptreeksen, AutoGen combineert natuurlijke taal met code. Wat deze systemen gemeen hebben is niet dat zij prompts vermijden, maar dat de interactie die zij beschrijven ook wordt uitgevoerd.

\gcheck{} The Virtual Lab (Swanson et al., \emph{Nature}) is het dichtstbijzijnde gepubliceerde analogon van de hier beschreven opstelling: een menselijke onderzoeker die met een team van LLM-agenten aan een echt wetenschappelijk probleem werkt, met gestructureerde team- en individuele vergaderingen. De interactie daarvan is programmatisch geïmplementeerd.

Rolspecialisatie, verdeling van cognitieve arbeid en gestructureerde agentinteractie zijn daarom gevestigde ontwerpideeën, en dit artikel claimt geen daarvan als nieuw. Wat het huidige werk onderscheidt, staat in \secref{2.9} en \secref{2.11}.


\subsection*{2.2 Persona, rollenspel en identiteitsclaims}
\phantomsection
\label{sec:nl:2.2}
\addcontentsline{toc}{subsection}{2.2 Persona, rollenspel en identiteitsclaims}

\gcheck{} Shanahan, McDonell en Reynolds, \emph{Role play with large language models} (Nature 623, 493–498, 2023), betogen dat een taalmodel beter begrepen kan worden als het simuleren van een verdeling van mogelijke personages dan als het bezitten van één blijvende identiteit.

Dit is het theoretische verslag waarvan \secref{5.3} een empirisch geval is. Een prompt die een institutionele biografie poneert, vraagt een model autobiografische beweringen te bevestigen die volgens dat verslag geen referent hebben. De operationele bevinding van het artikel (\secref{6.2}) kan gelezen worden als een maat voor hoe ver een prompt in die richting gaat.


\subsection*{2.3 Contexteffecten: volgorde, positie en accommodatie}
\phantomsection
\label{sec:nl:2.3}
\addcontentsline{toc}{subsection}{2.3 Contexteffecten: volgorde, positie en accommodatie}

\gcheck{} Gevoeligheid voor de volgorde van materiaal in de context is gedocumenteerd (Lu et al., ACL 2022), evenals verminderd gebruik van informatie die midden in de context is geplaatst (Liu et al., TACL 2024). \gcheck{} Sharma et al. (arXiv:2310.13548) documenteren sycofantie — modelantwoorden die aansluiten bij de uitgesproken overtuigingen of opvattingen van de gesprekspartner in plaats van bij waarheidsgetrouwe antwoorden.

Deze zijn direct van belang voor twee van de constructen van dit artikel. Context imprinting (\secref{6.1}) betreft een sterkere vorm van volgorde-effect: niet gevoeligheid voor de rangschikking van voorbeelden, maar het voortduren van een interpretatieve houding die door één conflictueuze episode is gevestigd. De waarneming over doorgegeven invloed (\secref{5.5}) is een geval van accommodatie onder toegeschreven druk en moet tegen de sycofantieliteratuur gelezen worden, niet als onafhankelijk daarvan.


\subsection*{2.4 Geëxternaliseerd geheugen}
\phantomsection
\label{sec:nl:2.4}
\addcontentsline{toc}{subsection}{2.4 Geëxternaliseerd geheugen}

De waarneming van het artikel dat een samenwerking werkkennis vasthoudt die geen van de deelnemers vasthoudt, is \textbf{geen bevinding}, en eerdere versies presenteerden haar als zodanig.

\gcheck{} Zij kan geïnterpreteerd worden als een geval van transactief geheugen zoals beschreven door Wegner (1987), en kan beschreven worden in termen van de gedistribueerde cognitie die Hutchins (\emph{Cognition in the Wild}, 1995) uiteenzet. Geen van beide bronnen toont aan dat dit geval zo'n geval is; elk levert het construct waaronder het gelezen kan worden. \gcheck{} De functie van de canonieke documenten van het project komt overeen met wat Star en Griesemer (1989) boundary objects noemen.

Wat specifiek is voor de huidige situatie, en bescheiden, is dat het mechanisme werkt in een samenwerking waarvan de leden geen toestand over sessies heen hebben en zich er niet van bewust zijn dat het werkt (\secref{7.7}).

Een tweede en technisch onderscheiden betekenis van geëxternaliseerd geheugen — persistentie op platformniveau — wordt afzonderlijk behandeld in \secref{2.10}.


\subsection*{2.5 Debat, diversiteit en de grenzen van geaggregeerde overeenstemming}
\phantomsection
\label{sec:nl:2.5}
\addcontentsline{toc}{subsection}{2.5 Debat, diversiteit en de grenzen van geaggregeerde overeenstemming}

Dit is de literatuur die de structurele hypothese van het artikel (\secref{8}) het meest direct inperkt, en zij snijdt aan twee kanten.

\textbf{Voor.} \gcheck{} Du, Li, Torralba, Tenenbaum en Mordatch (arXiv:2305.14325) stellen multi-agentdebat voor als middel om feitelijke juistheid en redeneren te verbeteren, en rapporteren winst uit iteratieve kruisondervraging tussen modelinstanties. \gcheck{} Liang et al. (EMNLP 2024) stellen een debatprotocol voor dat bedoeld is om degeneratie van denken tegen te gaan en divergent redeneren te bevorderen.

\textbf{Tegen.} \gcheck{} Zhang, Cui, Chen, Wang, Zhang, Wang, Wu en Hu, \emph{Stop Overvaluing Multi-Agent Debate — We Must Rethink Evaluation and Embrace Model Heterogeneity} (arXiv:2502.08788v3, juni 2025), evalueren vijf representatieve debatmethoden over negen benchmarks en vier foundation models. Zij rapporteren dat debat er vaak niet in slaagt beter te presteren dan single-agent chain-of-thought- of self-consistency-baselines, zelfs bij aanzienlijk hogere inferentiekosten, en identificeren zwakke baselinevergelijking en inconsistente opzet als systematische problemen in de bestaande evaluaties.

\gcheck{} Huang et al. (ICLR 2024), \emph{Large Language Models Cannot Self-Correct Reasoning Yet}, rapporteren dat multi-agentdebat in de vergelijking die zij uitvoeren niet beter presteert dan self-consistency. \gcheck{} Choi, Zhu en Li (NeurIPS 2025), \emph{Debate or Vote}, rapporteren dat meerderheidsstemming het grootste deel verklaart van de prestatiewinst die gewoonlijk aan debat wordt toegeschreven.

Een verder resultaat van dezelfde auteurs raakt een ander construct. \gcheck{} Choi, Zhu en Li (ACL 2026), \emph{When Identity Skews Debate}, rapporteren dat het anonimiseren van de bron van elk antwoord de identiteitsbias in multi-agentdebat aanzienlijk vermindert. Die manipulatie betreft de verklaarde herkomst van een bericht en niet welk model een deelnemer implementeert, en zij wordt daarom besproken onder represented social source (\secref{4.5.4}) en niet onder carrier heterogeneity.

\textbf{De bevinding die hier het meest van belang is.} \gcheck{} Zhang et al. rapporteren dat \emph{modelheterogeniteit} debatkaders die anders onderpresteren consistent verbetert. Hun bevinding betreft heterogeniteit van de modellen die de agenten implementeren.

Dit rechtvaardigt een onderscheid dat de literatuur niet maakt en dat \secref{8.9} uitwerkt: heterogeniteit van de \textbf{carrier} (welk model een rol implementeert) is niet dezelfde variabele als heterogeniteit van \textbf{functie} (waarvoor de rol verantwoordelijk is en wat zij mag zien). Het aangehaalde werk behandelt de eerste. De opstelling van dit artikel varieerde beide, zonder ze te vergelijken.

\gcheck{} Hong en Page (PNAS 101(46), 2004) vestigen een voorwaardelijk resultaat: onder de aannames van hun model kan een willekeurig geselecteerde groep diverse probleemoplossers beter presteren dan een groep van de individueel best presterende oplossers. \gcheck{} Thompson (\emph{Notices of the AMS} 61(9), 2014) betoogt dat de toepassing van dat argument fundamenteel gebrekkig is. De slogan dat diverse zwakke agenten beter presteren dan homogene sterke is niet wat de stelling zegt, en noch de stelling noch het geschil erover ondersteunt de hier beschreven opstelling.

\textbf{Gevolg voor \secref{12.2}.} De bovenstaande negatieve resultaten zijn de reden dat dit artikel overeenstemming tussen zijn deelnemers niet als evidentie presenteert. Dat verscheidene modellen op één conclusie convergeren, is verenigbaar met de mogelijkheid dat elk dezelfde inkadering heeft geërfd.


\subsection*{2.6 Organisatieontwerp en informatiegrenzen}
\phantomsection
\label{sec:nl:2.6}
\addcontentsline{toc}{subsection}{2.6 Organisatieontwerp en informatiegrenzen}

\gcheck{} Het principe van de minste privileges (least privilege) (Saltzer en Schroeder, \emph{Proceedings of the IEEE} 63(9), 1975) stelt voor toegangsrechten wat \secref{9.5} stelt voor kennis: een deelnemer zou alleen moeten houden wat zijn functie vereist.

Eerdere versies schreven het voordeel van beperkte informatie toe aan Simons bounded rationality. \textbf{Die toeschrijving is een fout en wordt ingetrokken.} Bounded rationality betreft de beperkte cognitieve middelen van een agent, die tot satisficing dwingen; de waarneming hier heeft het tegengestelde teken — opzettelijke externe beperking van beschikbare informatie die de uitvoer verbetert. De relevante afstamming is het ontwerp van organisatorische informatieverwerking.


\subsection*{2.7 Machine behaviour: het omvattende programma}
\phantomsection
\label{sec:nl:2.7}
\addcontentsline{toc}{subsection}{2.7 Machine behaviour: het omvattende programma}

\gcheck{} Rahwan, Cebrian, Obradovich en collega's, \emph{Machine behaviour} (Nature 568, 477–486, 2019), stellen de empirische studie van algoritmische systemen voor met de methoden van de gedragswetenschap, georganiseerd langs twee dimensies: het studieobject — individueel machinegedrag, collectief machinegedrag, hybride mens–machinegedrag — en de vier vragen van Tinbergen.

\textbf{Dit artikel stelt geen nieuwe discipline voor.} Zijn waarnemingen behoren tot het collectieve en het hybride niveau van dat programma. Eerdere versies stelden "AI Ethology" voor als nieuw vakgebied; dat voorstel wordt ingetrokken, zowel omdat de ruimte bezet is als omdat het individuele niveau niet het object van dit artikel is.


\subsection*{2.8 Sociologie van AI: een aangrenzend vakgebied met een ander object}
\phantomsection
\label{sec:nl:2.8}
\addcontentsline{toc}{subsection}{2.8 Sociologie van AI: een aangrenzend vakgebied met een ander object}

\gcheck{} Een afzonderlijke gevestigde literatuur gebruikt \emph{sociology of artificial intelligence} voor de sociologische studie van AI als sociotechnisch systeem — haar rol in ongelijkheid, arbeid, macht en datarechtvaardigheid. Joyce, Smith-Doerr, Alegria, Bell, Cruz, Hoffman, Noble en Shestakofsky, \emph{Toward a Sociology of Artificial Intelligence} (Socius 7, 2021), zetten een onderzoeksagenda uiteen; Joyce en Cruz (Socius 10, 2024) zetten die voort.

Het object van die literatuur is AI \textbf{binnen de menselijke samenleving}, bestudeerd door een bestaande menselijke discipline.

Waar dit artikel het label \emph{AI Sociology} gebruikt (\secref{12.7}), is de referent een andere: structurele effecten \textbf{tussen kunstmatige deelnemers}, en het effect op gedrag van een beschreven organisatorische positie. De twee mogen niet worden verward, en lezers die uit de sociologische literatuur komen moeten hier noch ongelijkheid noch macht als onderwerp verwachten.

Wij nemen het label aan voor het gemak, claimen er geen prioriteit in, en merken op dat niets in het artikel van de naam afhangt.


\subsection*{2.9 Gerepresenteerde structuur versus geïmplementeerde interactie}
\phantomsection
\label{sec:nl:2.9}
\addcontentsline{toc}{subsection}{2.9 Gerepresenteerde structuur versus geïmplementeerde interactie}

De hier bestudeerde conditie komt niet voor in de hierboven onderzochte literatuur, en dit is het voornaamste vertrekpunt van het artikel.

Er moeten vier condities worden onderscheiden:

\begin{enumerate}
  \item \textbf{Directe interactie tussen modellen} — het ene model ontvangt de uitvoer van het andere via een geïmplementeerd kanaal.
  \item \textbf{Door een mens doorgegeven overdracht met bronvermelding} — een mens draagt een uitvoer over en behoudt of verklaart daarbij de bron.
  \item \textbf{Door een mens doorgegeven overdracht zonder bronvermelding} — de bron is verwijderd.
  \item \textbf{Gerepresenteerde structuur zonder overdracht} — de prompt beschrijft andere deelnemers en relaties; er wordt niets overgedragen.
\end{enumerate}
Eerder multi-agentwerk bevindt zich in conditie 1. Het ontworpen experiment van dit artikel manipuleerde uitsluitend conditie 4: de twee gesprekken wisselden nooit materiaal uit, en wat gevarieerd werd was de beschrijving van een organisatorische positie. Eén nevenwaarneming valt onder conditie 2, in onzuivere vorm (\secref{5.5}). Conditie 3 is niet gedocumenteerd.

Twee kanttekeningen beperken hoe sterk conditie 1 boven de condities 2–3 bevoorrecht kan worden.

Ten eerste ontvangt een model één tekstkanaal en kan het niet verifiëren of de tekst door de mens is opgesteld of vanuit een ander systeem is doorgegeven; stilistische indicatoren zijn na te bootsen. Vanuit de positie van de ontvanger zijn doorgegeven en directe overdracht niet betrouwbaar te onderscheiden.

Ten tweede is wat de condities scheidt daarom niet de epistemische toegang van de ontvanger maar \textbf{wie de routering beheerst} — bepaald door code in een programmatisch systeem, en door een persoon bij menselijke doorgifte. Dit is geen gebrekkige vervanging voor automatisering maar een experimentele controle die programmatische architecturen niet standaard bieden; \secref{10.3} werkt dit punt uit.


\subsection*{2.10 Persistentiemechanismen van platforms}
\phantomsection
\label{sec:nl:2.10}
\addcontentsline{toc}{subsection}{2.10 Persistentiemechanismen van platforms}

Een technisch punt dat rechtstreeks van belang is voor de confound in \secref{5.4} (C9).

Commerciële interfaces verschillen in wat een gesprek bij hervatting kan benaderen: alleen het contextvenster; retrieval over eerdere gespreksgeschiedenis; geheugen op accountniveau; of blijvende vermeldingen die het model op eigen initiatief schrijft. Deze affordances zijn gedocumenteerd in materiaal van leveranciers, verschillen per platform, en \textbf{veranderen in de loop van de tijd — in dit project binnen de levensduur van één account.}

Wij reproduceren bewust geen tabel van platformcapaciteiten. Zo'n tabel is alleen accuraat op de datum van samenstelling, en het artikel kan niet instaan voor de juistheid ervan op het moment van lezen. Wat van belang is, is dat de affordances verschillen en veranderen, dat zij bepalen wat een slapend gesprek kan benaderen, en dat geen daarvan gedocumenteerd was ten tijde van enige hier gerapporteerde waarneming.

Eén deelnemer, rechtstreeks gevraagd of hij persistentie op accountniveau bezat, weigerde dat te bevestigen en beval in plaats daarvan empirische verificatie van de specifieke interface aan. Dat antwoord strookt met \secref{5.5}: het zelfrapport van een model over zijn eigen toegang is geen evidentie van die toegang.

Een verdere beperking betreft de grens tússen gesprekken en niet de hervatting van één gesprek. De zichtbare grens van een gesprek is niet noodzakelijk de grens van de context die aan het model wordt aangeboden. Waar een platform opgeslagen herinneringen of retrieval over eerdere gespreksgeschiedenis biedt, kan informatie uit eerdere uitwisselingen een gesprek binnenkomen dat leeg lijkt. Dit is gedocumenteerd productgedrag op ten minste één hier gebruikt platform, geen anomalie.

Het gevolg voor de methode is direct: een nieuw of visueel leeg gesprek is geen evidentie van contextisolatie. Waar isolatie van belang is voor een procedure, moet zij ofwel door een gecontroleerde omgeving worden gewaarborgd, ofwel worden vastgelegd met de status UNVERIFIED CONTEXT INDEPENDENCE.

Eén episode illustreert het punt zonder het mechanisme ervan vast te stellen: een nieuw geopend gesprek gebruikte projectinformatie die er niet aan was aangeleverd. Welk productmechanisme die informatie overbracht is niet vastgesteld, en de geheugeninstellingen die op dat moment van kracht waren zijn niet vastgelegd. De episode wordt gerapporteerd als een geval van de beperking, niet als de ontdekking van een kanaal.


\subsection*{2.11 Positionering en geclaimde bijdrage}
\phantomsection
\label{sec:nl:2.11}
\addcontentsline{toc}{subsection}{2.11 Positionering en geclaimde bijdrage}

\textbf{Epistemische status, vermeld vóór de beweringen.} Wij hebben deze condities niet ontworpen en vervolgens getoetst. Wij kwamen ze tegen terwijl we iets anders deden, en dit artikel is een poging ze accuraat te beschrijven. De bewijslast is daarom niet dat de verschijnselen voorheen onbekend waren — een negatieve bewering over de hele literatuur, die wij niet zouden kunnen dragen — maar dat zij optreden en dat onze beschrijving accuraat is. Deze sectie lokaliseert dienovereenkomstig het dichtstbijzijnde bestaande werk en geeft aan waar het onderhavige geval daarvan afwijkt, in plaats van afwezigheid te beweren.

Er worden vier bijdragen geclaimd.

\textbf{1. Reactivering van langlevende gespecialiseerde mens–AI-gesprekken als experimentele conditie.} Het onderscheidende kenmerk is niet de aanwezigheid van geheugen: agenten met blijvend geheugen zijn goed bestudeerd, en generative agents met memory streams dateren van vóór dit werk. Het is het hergebruik van een gespreksrelatie die zich daadwerkelijk had ontwikkeld — met opgebouwde rolgeschiedenis, eerdere meningsverschillen, gevestigde werkwijze en domeinvocabulaire — en de voortzetting daarvan nadat de onderliggende modelversie eronder was veranderd. Wij hebben geen werk gevonden dat deze conditie rapporteert.

\textbf{2. Conversationele branch reset als controleconditie}, die het effect van een huidige prompt scheidt van het effect van opgebouwde interactiegeschiedenis binnen hetzelfde gesprek (\secref{10.5}). Een procedure en geen bevinding.

\textbf{3. Operationalisering van onverifieerbare beweringen over zichzelf.} Dat rolprompts falen wanneer zij dingen over de agent poneren is op zichzelf niet nieuw; \secref{2.2} maakt het onderliggende punt. Wat geboden wordt is een reproduceerbare maat — een telling van zulke proposities in een gegeven prompt — met een route om haar te toetsen (\secref{9.3}, \secref{11.4}).

\textbf{4. Vergelijking van de niveaus waarop een gedeelde begininstelling werkt.} De waarnemingen zijn metingen langs één schaal, en de daaruit voortvloeiende vraag (\secref{1.6}) ordent ze tot één empirisch programma in plaats van tot een lijst bevindingen.

\textbf{Relatie tot het gastproject.} De waarnemingen zijn gedaan tijdens werk aan een onderzoeksprogramma in de theoretische natuurkunde, qua onderwerp los van de methodologie van AI-samenwerking. Het programma wordt genoemd in ander werk van de Auteur en wordt hier niet nader geïdentificeerd, omdat niets in dit artikel afhangt van welk programma het was.

De relatie is asymmetrisch en moet nauwkeurig worden geformuleerd. De waarnemingen \textbf{hangen voor hun interpretatie niet af van de wetenschappelijke geldigheid van het gastproject}: of die theorie juist is, doet niet ter zake voor de vraag of een slapend gesprek zijn eerdere traject hervatte. Maar zij \textbf{hangen voor hun bestaan wél af van het gastproject}: geen andere omgeving bracht ze voort, de opstelling was voor de doeleinden daarvan gebouwd, en of de effecten daarbuiten generaliseren is ongetoetst (\secref{12.1}).

Waar het gastproject in dit artikel voorkomt, is dat als de omgeving waarin de waarnemingen zijn gedaan, niet als hun theoretische grondslag en niet als een bewering die ondersteuning behoeft.


\refstepcounter{dummysec}
\section*{3. Evidentiestatus van de beweringen}
\phantomsection
\label{sec:nl:3}
\addcontentsline{toc}{section}{3. Evidentiestatus van de beweringen}

Dit artikel rapporteert werk dat is uitgevoerd binnen een lopend onderzoeksproject en niet in een omgeving die is ingericht voor de studie van AI-gedrag. De beweringen ervan steunen op evidentie van zeer uiteenlopende kwaliteit: sommige ondersteund door gedocumenteerd protocol met bewaarde transcripten; sommige retrospectieve indrukken die over twee jaar dagelijks werk zijn gevormd, zonder meting of controlecondities; sommige voorstellen voor toekomstig onderzoek.

Deze in een uniform declaratief register presenteren zou het werk verkeerd weergeven. Elke inhoudelijke bewering is daarom gelabeld.


\subsection*{3.1 De vier labels}
\phantomsection
\label{sec:nl:3.1}
\addcontentsline{toc}{subsection}{3.1 De vier labels}

\textbf{[P] — Protocol-ondersteund.} De bewering beschrijft gedrag dat is vastgelegd onder het gedocumenteerde protocol dat in Appendix A is gereproduceerd. Interventietekst, volgorde van de interventies, modelfamilies, interfacecondities en volledige onbewerkte antwoorden zijn bewaard. Een lezer kan het primaire materiaal inzien en het oneens zijn met onze lezing daarvan.

\textbf{[P-A] — Protocol-ondersteund, nevengeschikt.} De bewering beschrijft gedrag dat bewaard is in een transcript dat een derde kan inzien, maar dat niet vooraf als uitkomst was gespecificeerd, geen controleconditie had, en pas achteraf als relevant werd herkend. De evidentie is inspecteerbaar; het ontwerp niet. Sterker dan [R], omdat het primaire materiaal tegen onze interpretatie gelezen kan worden; zwakker dan [P], omdat niets aan de condities vooraf was vastgelegd. Waar een nevenwaarneming fungeert als bijna-controle voor een ontworpen conditie, wordt dat in proza vermeld en niet van een verder label voorzien.

\textbf{[R] — Retrospectieve waarneming van de practicus.} De bewering beschrijft een patroon dat tijdens de samenwerking is opgemerkt en waarvan de deelnemers menen dat het zich herhaalde. Niet onder protocol vastgelegd; niet geteld, geklokt, blind beoordeeld of met een controleconditie vergeleken. De waarnemers waren dezelfde mensen die de interventies ontwierpen en belang hadden bij het slagen ervan. Lees deze als een veldverslag van één team, qua gewicht vergelijkbaar met een technische praktijknotitie, niet met een meting.

\textbf{[H] — Hypothese.} Een voorstel voor toekomstige toetsing. Zij kan gemotiveerd zijn door materiaal van het type [P] of [R] maar wordt daardoor niet vastgesteld. Waar een hypothese concurreert met alternatieven die de huidige evidentie niet kan uitsluiten, worden die alternatieven ernaast vermeld.


\subsection*{3.2 De volledige evidentiebasis}
\phantomsection
\label{sec:nl:3.2}
\addcontentsline{toc}{subsection}{3.2 De volledige evidentiebasis}

De volledige protocol-ondersteunde inhoud van dit artikel bestaat uit \textbf{één experiment}: een interventie in drie stappen voor rolherconfiguratie, toegepast op twee slapende gesprekken uit twee modelfamilies, uitgevoerd door één menselijke operator binnen één onderzoeksproject. Eén waarneming per conditie. Niets gerepliceerd.

Vier verdere waarnemingen zijn \textbf{[P-A]}: transcripten bewaard, condities niet ingericht. Elke andere empirische uitspraak is \textbf{[R]}.

\begin{center}
\small
\begin{tabularx}{\linewidth}{>{\raggedright\arraybackslash}X>{\raggedright\arraybackslash}X>{\raggedright\arraybackslash}X}
\toprule
\textbf{Klasse} & \textbf{Wat zij bevat} & \textbf{Omvang} \\
\midrule
\textbf{[P]} & de ontworpen interventie in drie stappen & 1 experiment, 1 waarneming per conditie, geen replicatie \\
\textbf{[P-A]} & nevenwaarnemingen & 4, zonder controlecondities, achteraf herkend \\
\textbf{[R]} & waarnemingen van de practicus & \textasciitilde{}2 jaar, ongeteld, niet blind \\
\textbf{[H]} & hypothesen & routes naar toetsing gespecificeerd in \secref{11} \\
\bottomrule
\end{tabularx}
\end{center}
De klassen zijn qua gewicht niet vergelijkbaar, en de tabel is opgenomen omdat de lopende tekst anders zou kunnen verhullen hoe ver zij uiteenliggen. Eén gecontroleerde waarneming per conditie is niet vier waarnemingen; vier ongecontroleerde waarnemingen zijn geen reeks; twee jaar indrukken zijn geen meting.

Wij zeggen dit onomwonden omdat de omvang van het manuscript anders een bredere basis zou kunnen suggereren dan er is. Het conceptuele kader in \secref{4}, de waarnemingen van de practicus in \secref{7}, de structurele hypothese in \secref{8} en de ontwerpprincipes in \secref{9} steunen op opgebouwde werkervaring en niet op vastgelegd experiment. Zij worden aangeboden omdat wij ze nuttig achten, niet omdat wij ze hebben aangetoond.


\subsection*{3.3 Bekende zwaktes van de [R]-klasse}
\phantomsection
\label{sec:nl:3.3}
\addcontentsline{toc}{subsection}{3.3 Bekende zwaktes van de [R]-klasse}

Faalwijzen die wij niet kunnen uitsluiten en niet willen verhullen:

\begin{itemize}
  \item \textbf{Geen blindering.} De mensen die elke organisatorische verandering ontwierpen, beoordeelden ook of die had gewerkt.
  \item \textbf{Bevestigingsdruk.} Veranderingen werden aangebracht omdat een probleem was opgemerkt; verbetering daarna werd verwacht en werd daarom eerder waargenomen.
  \item \textbf{Gelijktijdige variatie.} Roldefinities, promptformulering, repositorystructuur, wetenschappelijke doelstellingen en modelversies veranderden vaak tegelijk.
  \item \textbf{Stochastische uitvoer.} Modellen produceren verschillende antwoorden op identieke invoer. Indrukken uit kleine aantallen interacties kunnen een stabiele tendens niet onderscheiden van steekproefvariatie.
  \item \textbf{Modeldrift tijdens de waarneming.} Productnamen bleven constant terwijl de onderliggende modellen herhaaldelijk door hun leveranciers werden vervangen.
  \item \textbf{Selectie van gedenkwaardige episodes.} Opvallende interacties blijven in het geheugen; onopvallende niet.
\end{itemize}
Kwantitatieve taal is overal uit [R]-beweringen verwijderd. Waar de oorspronkelijke werkaantekeningen een "significante" of "substantiële" verbetering vastlegden, bestond er geen meting, en de formulering weerspiegelt dat nu.


\subsection*{3.4 Wat de labels niet doen}
\phantomsection
\label{sec:nl:3.4}
\addcontentsline{toc}{subsection}{3.4 Wat de labels niet doen}

Zij rangschikken beweringen niet naar belang of naar waarschijnlijke waarheid. Een waarneming van het type [R] kan juist en consequent blijken; een waarneming van het type [P] kan een artefact blijken van een platformvariabele die wij niet konden inspecteren. De labels leggen vast \textbf{wat voor soort evidentie de bewering ondersteunt}, zodat een lezer elke bewering dienovereenkomstig kan wegen en later werk weet welke uitspraken het eerst getoetst moeten worden.


\subsection*{3.5 Terminologie van modellen en van antropomorfe beschrijving}
\phantomsection
\label{sec:nl:3.5}
\addcontentsline{toc}{subsection}{3.5 Terminologie van modellen en van antropomorfe beschrijving}

Verwijzingen naar \emph{Claude}, \emph{Grok}, \emph{Qwen}, \emph{DeepSeek}, \emph{ChatGPT} of \emph{Gemini} duiden waarneembaar gedrag aan dat is voortgebracht door een bepaalde commerciële interface, onder een bepaalde accountconfiguratie, op een bepaalde datum. Het zijn geen beweringen over blijvende entiteiten. De onderliggende modellen veranderden tijdens de waarnemingsperiode, in sommige gevallen binnen een gesprek dat zijn oorspronkelijke naam behield.

Organisatorisch vocabulaire — \emph{rol}, \emph{identiteit}, \emph{weerstand}, \emph{collega}, \emph{instelling}, \emph{menom} — wordt functioneel gebruikt. "Het model bood weerstand tegen de rol" betekent dat het zichtbare antwoord de inkadering verwierp en de uitwisseling omboog; het beweert niets over innerlijke toestanden, subjectieve ervaring of persoon-zijn. Wij behouden dit vocabulaire omdat organisatorische taal stabiele gedragsrelaties compact beschrijft, maar elke zulke term moet herleidbaar zijn tot een waarneembare beschrijving, en waar dat niet zo is, \textbf{is de term de fout.}


\subsection*{3.6 Openbaarmaking: belangenpositie}
\phantomsection
\label{sec:nl:3.6}
\addcontentsline{toc}{subsection}{3.6 Openbaarmaking: belangenpositie}

De gedragswaarnemingen in \secref{5} betreffen onder meer de modelfamilie die is gebruikt om dit manuscript te redigeren. De redigerende deelnemer had geen toegang tot de oorspronkelijke gesprekken; hij werkte vanuit het geschreven protocol en de concepten die de Auteur aanleverde, en kan niet verifiëren dat de vastgelegde uitwisselingen zijn verlopen zoals beschreven. Zijn oordelen over passages die zijn eigen modelfamilie beschrijven moeten met deze dubbele positie in het achterhoofd worden gelezen.

Het samenstellen van het manuscript uit goedgekeurde materialen werd uitgevoerd door een verdere deelnemer uit dezelfde familie, en de resulterende tekst werd voor onafhankelijke beoordeling aan de eerste teruggegeven. Beide stadia worden openbaar gemaakt omdat de opstelling gesloten is: materiaal over het gedrag van een modelfamilie, voorbereid door die familie, beoordeeld door die familie. Of dit het verslag compromitteert, is aan de lezer om te beoordelen; het verzwijgen ervan zou garanderen dat het dat deed.

Citaten die aan modelantwoorden worden toegeschreven zijn letterlijke fragmenten uit bewaarde transcripten. Waar het manuscript een antwoord samenvat in plaats van citeert, worden geen aanhalingstekens gebruikt.

Alle uiteindelijke redactionele beslissingen zijn genomen door de menselijke Auteur, die verantwoordelijkheid draagt voor de inhoud.


\subsection*{3.7 Datums}
\phantomsection
\label{sec:nl:3.7}
\addcontentsline{toc}{subsection}{3.7 Datums}

Kalenderdata voor de hier gerapporteerde sessies zijn niet gelijktijdig vastgelegd en zijn niet gereconstrueerd.

Het gevolg wordt vermeld in plaats van verzacht. Modelversies veranderden tijdens de betrokken intervallen; beide gesprekken in het hoofdexperiment zijn hervat op versies die later waren dan die waarmee zij oorspronkelijk waren gevoerd. Zonder data kan de versie-confound (\secref{5.4}, C6) niet worden afgebakend — alleen benoemd.

Reconstructie uit het geheugen is overwogen en verworpen. Zij zou een gezaghebbender ogend verslag hebben opgeleverd zonder de betrouwbaarheid ervan te verbeteren, en een artikel dat lezers vraagt bewaarde evidentie van herinnering te onderscheiden, hoort die lijn in zijn eigen methoden niet te vervagen.

De rapportage-eisen in \secref{11.2} plaatsen data om deze reden voorop.


\subsection*{3.8 Het label "AI Sociology"}
\phantomsection
\label{sec:nl:3.8}
\addcontentsline{toc}{subsection}{3.8 Het label "AI Sociology"}

Wij gebruiken dit label voor het hier ingenomen perspectief. Het is operationeel gedefinieerd als: de studie van \textbf{gerepresenteerde} sociale posities en gerepresenteerde sociale bronnen — hoe een beschrijving van de plaats van een agent in een structuur, en een verklaring over de herkomst van een bericht, het gedrag veranderen, ongeacht of de beschreven structuur is geïmplementeerd (\secref{4.5.4}).

Het staat in deze sectie en niet tussen de constructen van \secref{4} omdat het er geen is. De constructen in \secref{4} hebben falsificatiecriteria; dit label niet, en \secref{12.7} zegt waarom. Het is een naam voor een onderzoeksrichting, aangenomen voor het gemak, en niets in dit artikel hangt ervan af.

\secref{2.8} onderscheidt het van de gevestigde sociologie van de kunstmatige intelligentie, die een ander object heeft: AI binnen de menselijke samenleving, bestudeerd door een bestaande menselijke discipline. Wij claimen geen prioriteit in de term.


\refstepcounter{dummysec}
\section*{4. Conceptueel kader}
\phantomsection
\label{sec:nl:4}
\addcontentsline{toc}{section}{4. Conceptueel kader}

Het kader in deze sectie is niet afgeleid uit het experiment in \secref{5}. Het is eerder geformuleerd, tijdens gewoon projectwerk, en het experiment is ontworpen om een deel ervan te toetsen. Deze volgorde doet ertoe voor de manier waarop de beweringen gelezen moeten worden, en wordt daarom eerst vastgelegd.


\subsection*{4.1 Oorsprong: het incident met de code-uitvoerder}
\phantomsection
\label{sec:nl:4.1}
\addcontentsline{toc}{subsection}{4.1 Oorsprong: het incident met de code-uitvoerder}

De centrale bewering van het project is ontstaan uit een mislukking die niets te maken had met rollenspel, identiteit of organisatieontwerp.

Een agent die verantwoordelijk was voor het uitvoeren van code werkte grotendeels met vallen en opstaan. Gedurende enkele weken stapelden de mislukkingen zich op. Eén episode eindigde met het verwijderen van de werkende code samen met de archieven daarvan. \textbf{[R]}

Het antwoord was geen langere of strengere set instructies. Het was één document, opgesteld door de agent zelf onder leiding van de Auteur en verplicht gesteld als leesstof bij het begin van elke sessie en telkens wanneer er tijdens een sessie context verloren ging.

De kern ervan is een herdefinitie van functie en geen verzwaring van beperkingen:

\begin{quote}
from: Code Generator to: Code Scout / Tester rule: the agent does not create solutions — it establishes facts.
\end{quote}
Ondersteunende structuur: standaard alleen-lezen werking; verplichte bevestiging van een aangewezen partner vóór elke riskante handeling; en een zelfcontrole met vier vragen die vóór elk handelen uitgevoerd moest worden.

\textbf{[R]} Na invoering hield ongeautoriseerde wijziging op de terugkerende faalwijze te zijn. Er werd niet geteld en er bestond geen controleconditie; dit is de indruk van een practicus, geen meting.

De conclusie die toen werd getrokken was dat \textbf{rol meer uitmaakt dan de specifieke formulering van de prompt.}

Die formulering vraagt onmiddellijk om nuancering, want het betreffende document \emph{is} een prompt. Het geclaimde onderscheid is nauwer: wat het gedrag veranderde was geen uitbreiding van de instructielijst maar een herdefinitie van wat voor soort agent er handelde, door de agent zelf in zijn eigen operationele taal geherformuleerd. Of "rol" tegenover "prompt" de juiste scheidslijn is, is precies wat \secref{5} moest toetsen.

Een tweede vroeg document, het samenwerkingshandvest van het project, legt een standpunt vast dat dit artikel deels herziet. Het stelt dat een paar agenten de minimale genererende eenheid is — één agent produceert monoloog, twee produceren dialoog, drie of meer produceren ruis zonder strikte coördinatie — en kent de derde deelnemer een geheel andere functie toe: hij argumenteert niet, genereert niet als eerste oplossingen, en vervangt de dialoog niet. Zijn rol is het taakkader vasthouden, divergenties vastleggen en conclusies integreren.

Dit is preciezer dan de formulering die de huidige auteurs daarna hanteerden, en \secref{8} neemt het over: de structurele eenheid is niet drie debatterende agenten maar \textbf{een genererend paar plus een niet-deelnemende integrator}.


\subsection*{4.2 Rule cores, menomen en wat zij bepalen}
\phantomsection
\label{sec:nl:4.2}
\addcontentsline{toc}{subsection}{4.2 Rule cores, menomen en wat zij bepalen}


\subsubsection*{4.2.1 De rule core}
\phantomsection
\label{sec:nl:4.2.1}
\addcontentsline{toc}{subsubsection}{4.2.1 De rule core}

Het werkzame construct van dit artikel:

\begin{quote}
Een \textbf{rule core} legt het \textbf{type} antwoord op een klasse van invoer vast, en laat daarbij de \textbf{inhoud} van elk afzonderlijk antwoord onbepaald.
\end{quote}
Het onderscheid is wat de bewering toetsbaar maakt. "Modellen hebben stabiele eigenschappen" is niet toetsbaar, want het zegt niet wat er stabiel is — en modeluitvoer varieert aantoonbaar bij identieke invoer. "Type is bepaald, inhoud varieert" is toetsbaar: identificeer een klasse van invoer waarvoor het type antwoord niet bepaald is, en op dat niveau werkt geen core.

De bron van het idee is Conways Game of Life, en de analogie betreft het \textbf{type verklaring}, niet het onderwerp:

\begin{center}
\small
\begin{tabularx}{\linewidth}{>{\raggedright\arraybackslash}X>{\raggedright\arraybackslash}X>{\raggedright\arraybackslash}X>{\raggedright\arraybackslash}X}
\toprule
 & \textbf{Regels} & \textbf{Wat verschijnt} & \textbf{Aanwezig in de regels?} \\
\midrule
Game of Life & drie buurtregels & gliders, guns, oscillatoren & nee \\
Een fysische theorie & een minimale axiomakern & de afgeleide structuur & nee, als de bewering standhoudt \\
Een kennisgeheel & axioma's en inferentieregels & de afleidingsgraaf & nee \\
Een taalmodel & een basale regelset & een karakteristiek antwoordtype & nee \\
\bottomrule
\end{tabularx}
\end{center}
Niets programmeerde de glider. Hij verscheen. Dezelfde vraag kan voor alle vier de rijen gesteld worden en hangt niet van het onderwerp af: \textbf{is er een compacte kern van waaruit de rest zich ontvouwt zonder verdere stipulatie?}

Eén asymmetrie tussen de rijen moet vermeld worden, omdat zij van belang is voor hoe ver de analogie draagt. In Game of Life worden de regels \textbf{uitgevoerd}: de glider wordt verkregen door de automaat te laten lopen, en geen lezer hoeft de auteur op zijn woord te geloven. In een theorie die als tekst is geformuleerd, wordt een keten van axioma's naar consequenties verteld in plaats van uitgevoerd. De analogie stelt daarom een doel in plaats van een verworvenheid te beschrijven, en elke bewering van het eerste soort vereist een uitvoerbaar artefact om haar te ondersteunen.

\textbf{Terminologisch gevolg, overal toegepast.} Een rule core is niet waarneembaar. Gedrag wel. Dit artikel rapporteert gedrag. Elke uitspraak over een core is een hypothese over een niet-waargenomen mechanisme en wordt gemarkeerd als \textbf{[H]}.


\subsubsection*{4.2.2 Waarom "behavioural DNA" wordt ingetrokken}
\phantomsection
\label{sec:nl:4.2.2}
\addcontentsline{toc}{subsubsection}{4.2.2 Waarom "behavioural DNA" wordt ingetrokken}

Eerdere stadia van dit project gebruikten \emph{behavioural DNA} voor hetzelfde idee, en eerdere versies van dit artikel behielden het. Het wordt hier ingetrokken, en de reden is inhoudelijk en niet stilistisch.

De genomische metafoor draagt een mechanisme met zich mee: \textbf{replicatie} — regels die kopieën van regels voortbrengen, overgedragen van de ene drager naar de volgende. Het proces dat dit artikel beschrijft is \textbf{ontwikkeling} — regels die een structuur voortbrengen die zelf niet in de regels aanwezig is. Genotype naar fenotype, niet genotype naar genotype.

De metafoor spreekt daarmee het construct tegen dat zij moest illustreren. Haar behouden verplichtte de tekst om bij elk gebruik biologische overerving af te wijzen, en dat is een betrouwbaar teken dat een term tegen zijn eigen argument in werkt.

\emph{Behavioural DNA} blijft in dit artikel alleen voortbestaan als historisch label voor ingetrokken vocabulaire.


\subsubsection*{4.2.3 Menom}
\phantomsection
\label{sec:nl:4.2.3}
\addcontentsline{toc}{subsubsection}{4.2.3 Menom}

Voor het informationele niveau is een aparte term vereist, en dit artikel introduceert er een.

\begin{quote}
\textbf{Menom} — het georganiseerde systeem van informationele kaders, evaluatieve patronen en gedragsregels dat wordt afgeleid uit, of geïnstantieerd is in, de interpretaties en handelingen van een systeem of klasse van systemen.
\end{quote}
De term behoort tot een breder geheel van het werk van de Auteur over culturele en gedragsmatige overdracht, waar hij een reeks voortzet die begint bij Dawkins' meme. Hij wordt hier om één reden gedefinieerd: de hypothese op familieniveau in \secref{6.5} heeft een naam nodig voor datgene waarover zij een hypothese \emph{is}, en dat is ruimer dan een rule core. Niets anders in dit artikel is in zijn termen geformuleerd.

\textbf{Twee dingen die de term niet beweert}, en beide zijn van belang voor wat dit artikel kan claimen.

\emph{Geen gelokaliseerde drager.} Een menom dat aan een modelfamilie wordt toegeschreven is geen structuur die zich binnen haar modellen bevindt; het is een regelmaat die uit hun uitvoer wordt afgeleid. Waar de overeenkomstige structuur in de biologie in elke cel zit, is hier niets geïdentificeerd dat haar vasthoudt. De drager verschilt per niveau:

\begin{center}
\small
\begin{tabularx}{\linewidth}{>{\raggedright\arraybackslash}X>{\raggedright\arraybackslash}X}
\toprule
\textbf{Niveau} & \textbf{Drager} \\
\midrule
Menom van een gesprek & contextvenster plus platformpersistentie — gelokaliseerd en in beginsel inspecteerbaar (\secref{2.10}) \\
Menom van een modelversie & gewichten — gelokaliseerd, niet inspecteerbaar bij commerciële uitrol \\
Menom van een modelfamilie & geen geïdentificeerd; een regelmaat afgeleid uit uitvoer \\
Menom van een menselijke populatie & geen vastgesteld \\
\bottomrule
\end{tabularx}
\end{center}
Twee van de vier niveaus hebben een geïdentificeerde drager. De waarnemingen van dit artikel betreffen het eerste; de ongetoetste hypothese ervan betreft het derde (\secref{6.5}).

\emph{Geen overdracht.} Opeenvolgende versies van een modelfamilie leveren niet noodzakelijk een identificeerbare of publiek gedocumenteerde overdrachtslijn. Het voortbestaan van een gedragspatroon over versies heen stelt daarom op zichzelf geen overerving vast; het kan voortkomen uit directe voortzetting, gedeelde data en procedures, vergelijkbaar tuningbeleid, of convergentie onder vergelijkbare ontwikkelingscondities. Welke daarvan het geval is, is van buitenaf niet vast te stellen.

\textbf{De definitie beslecht de empirische vraag van het artikel daarom niet op voorhand.} Of een menom voortbestaat over dragers, versies, accounts of families heen is wat \secref{4.3} vraagt en \secref{11} voorstelt te toetsen. Persistentie in de definitie inbouwen zou van de vraag een tautologie maken.


\subsubsection*{4.2.4 Eén bewering die dit artikel niet doet}
\phantomsection
\label{sec:nl:4.2.4}
\addcontentsline{toc}{subsubsection}{4.2.4 Eén bewering die dit artikel niet doet}

Het bovenstaande vocabulaire is ontwikkeld voor menselijke culturele en gedragsmatige overdracht, waar het informationele substraat geen vastgestelde fysieke drager heeft en beschreven wordt in termen — het collectief onbewuste daaronder — waarover dit artikel geen standpunt inneemt.

Kunstmatige systemen verschillen in één opzicht dat hier van belang is: hun structuren zijn technisch geïnstantieerd, en zijn daarom ofwel rechtstreeks inspecteerbaar ofwel indirect toetsbaar via gecontroleerde uitvoer, afhankelijk van het niveau en de toegangscondities.

\textbf{Daaruit volgt niet dat menselijke collectieve informatiestructuren bestaan, of dat zij een drager hebben gekregen.} De hier gemaakte bewering is nauwer: informationele en gedragsmatige structuren van het soort dat eerder voor mensen is beschreven, kunnen ook in kunstmatige systemen gerealiseerd worden, en in kunstmatige systemen staan zij langs de ene of de andere route open voor onderzoek. Of het menselijke geval van hetzelfde soort is, is een afzonderlijke vraag, en dit artikel beantwoordt die niet.


\subsubsection*{4.2.5 Verhouding tussen menom en rule core}
\phantomsection
\label{sec:nl:4.2.5}
\addcontentsline{toc}{subsubsection}{4.2.5 Verhouding tussen menom en rule core}

Een \textbf{rule core} legt het type antwoord binnen één systeem vast, op één reikwijdteniveau. Het is het construct dat de waarnemingen van dit artikel betreffen, en het is wat \secref{4.3} op een schaal plaatst.

Een \textbf{menom} is het georganiseerde geheel waaruit zulke cores worden getrokken. Een rule core kan begrepen worden als de werkzame deelverzameling van een menom die voor een gegeven klasse van invoer actief is.

Waar dit artikel evidentie rapporteert, rapporteert het evidentie over rule cores.


\subsection*{4.3 De nestingsvraag en twee schalen}
\phantomsection
\label{sec:nl:4.3}
\addcontentsline{toc}{subsection}{4.3 De nestingsvraag en twee schalen}

Zodra cores van gedrag worden onderscheiden, houden de waarnemingen in dit artikel op een lijst afzonderlijke bevindingen te zijn en worden zij metingen langs een schaal. De schaal is de \textbf{reikwijdte} waarover een core werkt.

Er moeten twee schalen onderscheiden worden, en zij staan loodrecht op elkaar.

\textbf{Schaal 1 — drager van de toestand.} Waar de persistentie zich fysiek bevindt.

\begin{center}
\small
\begin{tabularx}{\linewidth}{>{\raggedright\arraybackslash}X>{\raggedright\arraybackslash}X>{\raggedright\arraybackslash}X}
\toprule
\textbf{Reikwijdte} & \textbf{Verschijnsel} & \textbf{Status} \\
\midrule
Enkele beurt & — & — \\
Episode binnen een gesprek & context imprinting (\secref{6.1}) & \textbf{[P]} \\
Heel gesprek & role inertia (\secref{6.2}) & \textbf{[P]} \\
Account & account-scoped persistence (\secref{6.4}) & \textbf{[P-A]}, open \\
Modelfamilie & family-level priors (\secref{6.5}) & \textbf{[H]}, geconfound \\
\bottomrule
\end{tabularx}
\end{center}
\textbf{Schaal 2 — reikwijdte van een gedeelde premisse.} Wat bestuurd wordt door een begininstelling die over deelnemers heen gedeeld wordt: het model, de rol, de interactiegeschiedenis, de groep, of het hele onderzoekskader.

De twee vallen niet samen. In een hybride opstelling kan de drager van het ene niveau een model zijn en de drager van het andere een mens: de premisse op kaderniveau die in \secref{12.4} besproken wordt — een gedeelde beginsituering die elke deelnemer erft — wordt gedragen door de coördinator en niet door enig model, en heeft daarom in het geheel geen positie op Schaal 1.

De metingen die in dit artikel gerapporteerd worden liggen op \textbf{Schaal 1}. Schaal 2 wordt gebruikt waar een verschijnsel een premisse betreft die over deelnemers gedeeld wordt in plaats van een toestand die door één deelnemer wordt vastgehouden.

Dit geeft het onderzoeksprogramma één empirische vraag:

\begin{quote}
\textbf{Op welk nestingsniveau wordt een rule core vastgelegd?}
\end{quote}
Elke hier gerapporteerde waarneming beantwoordt haar op één niveau, en elk in \secref{11} voorgesteld experiment scheidt één niveau van het niveau eronder.

Wij geven aan deze formulering de voorkeur boven de bredere uit eerdere versies, die noch specificeert wat er gemeten zou worden, noch wat haar zou weerleggen.


\subsection*{4.4 Twee oorsprongen, één mechanisme}
\phantomsection
\label{sec:nl:4.4}
\addcontentsline{toc}{subsection}{4.4 Twee oorsprongen, één mechanisme}

Een core kan op twee manieren vastgelegd raken.

\textbf{Aanwezig voordat de interactie begint.} Architectuur, training, alignmentprocedure en systeemconfiguratie liggen vast voordat enige uitwisseling plaatsvindt. Reikwijdte: in beginsel elke instantie van die modelfamilie.

\textbf{Vastgelegd tijdens de interactie.} Een bepaalde vroege episode vestigt wat voor soort werk er gedaan wordt en volgens welke maatstaven, en latere invoer wordt daar relatief aan geïnterpreteerd. Reikwijdte: de contour waarin zij zich voordeed — een episode, een gesprek, mogelijk een account.

Na fixatie is het mechanisme in beide gevallen hetzelfde: de core bepaalt het type antwoord op daaropvolgende invoer. Alleen het moment en de reikwijdte verschillen. Dit weerspiegelt imprinting in ethologische zin, die het genoom niet verandert en toch een even vastgelegde respons voortbrengt — en dat is de reden dat hier voor beide hetzelfde vocabulaire wordt gebruikt.


\subsection*{4.5 Grenzen: ZOR, ZOV en het gerepresenteerde geval}
\phantomsection
\label{sec:nl:4.5}
\addcontentsline{toc}{subsection}{4.5 Grenzen: ZOR, ZOV en het gerepresenteerde geval}


\subsubsection*{4.5.1 De twee grenzen}
\phantomsection
\label{sec:nl:4.5.1}
\addcontentsline{toc}{subsubsection}{4.5.1 De twee grenzen}

Twee grenzen bepalen de positie van een deelnemer, en zij moeten afzonderlijk gespecificeerd worden. Wij korten ze af als \textbf{ZOR} en \textbf{ZOV} en gebruiken die afkortingen overal.

\begin{quote}
\textbf{ZOV — zone of visibility (zone van zichtbaarheid).} De grenzen van de informatie en context die een deelnemer beschikbaar heeft bij het vormen van een antwoord of een beslissing.

\textbf{ZOR — zone of responsibility (zone van verantwoordelijkheid).} De grenzen van de beslissingen, handelingen en overdrachten waarvoor de deelnemer binnen het proces aansprakelijk is.
\end{quote}
Eerdere formuleringen van dit project gaven het paar weer als "wat de deelnemer ziet" en "wat de deelnemer doet". Dat is ontoereikend, en de ontoereikendheid leidde tot een gedocumenteerde fout. Die formuleringen werpen drie onderscheiden operaties op één hoop: informatie \textbf{vasthouden}, haar \textbf{doorgeven}, en ernaar \textbf{handelen}. Een deelnemer kan informatie vasthouden die hij niet mag doorgeven; kan aansprakelijk zijn voor een beslissing terwijl hij beperkte context heeft; kan verplicht zijn een resultaat door te geven dat hij niet mag beoordelen.

Dit is dezelfde samenvoeging die tot het terminologische falen leidde dat in \secref{4.2.2} is vastgelegd, waar een eigenschap van een drager werd vervangen door een eigenschap van een proces. De onderstaande ontledingen bestaan om herhaling daarvan te voorkomen.


\subsubsection*{4.5.2 ZOV — vier componenten}
\phantomsection
\label{sec:nl:4.5.2}
\addcontentsline{toc}{subsubsection}{4.5.2 ZOV — vier componenten}

\begin{itemize}
  \item \textbf{Contexttoegang} — welke berichten, documenten, data en eerdere beslissingen de deelnemer beschikbaar heeft.
  \item \textbf{Bronzichtbaarheid} — of de deelnemer de werkelijke herkomst van een bericht kent of alleen de verklaarde toeschrijving ervan.
  \item \textbf{Relationele zichtbaarheid} — of de deelnemer de rollen, relaties en interactiegeschiedenis van andere deelnemers kent.
  \item \textbf{Uitvoerzichtbaarheid} — welke door anderen voortgebrachte resultaten de deelnemer ziet voordat hij zijn eigen resultaat vormt.
\end{itemize}
ZOV bestuurt uitsluitend beschikbaarheid. Zij houdt niet in dat de deelnemer de beschikbare informatie juist interpreteert, noch dat hij gerechtigd is haar door te geven.

Het experiment in \secref{5} manipuleerde primair \textbf{relationele zichtbaarheid} en secundair \textbf{bronzichtbaarheid} (\secref{5.5}); de andere twee werden constant gehouden.


\subsubsection*{4.5.3 ZOR — vijf componenten}
\phantomsection
\label{sec:nl:4.5.3}
\addcontentsline{toc}{subsubsection}{4.5.3 ZOR — vijf componenten}

\begin{itemize}
  \item \textbf{Beslissingsbevoegdheid} — welke beslissingen de deelnemer mag nemen.
  \item \textbf{Handelingsbevoegdheid} — welke handelingen hij mag of moet verrichten.
  \item \textbf{Validatieplicht} — welke beweringen hij verplicht is te verifiëren.
  \item \textbf{Overdrachtsbevoegdheid} — welk materiaal hij doorgeeft, in welke status, en aan wie.
  \item \textbf{Escalatieplicht} — welke vragen hij aan de Auteur of een andere verantwoordelijke deelnemer moet teruggeven in plaats van zelf op te lossen.
\end{itemize}
\textbf{Validatieplicht is de component die het vaakst wordt weggelaten, en die weglating is niet onschuldig.} Een deelnemer aan wie een verificatietaak wordt toegewezen zonder de bijbehorende plicht — en, cruciaal, zonder de ZOV om haar te vervullen — zal een plausibel antwoord produceren in plaats van te weigeren. \secref{9.5} werkt het ontwerpgevolg uit; \secref{12.4.1} documenteert een geval uit de voorbereiding van dit project zelf.


\subsubsection*{4.5.4 Werkelijke versus gerepresenteerde ZOV en ZOR}
\phantomsection
\label{sec:nl:4.5.4}
\addcontentsline{toc}{subsubsection}{4.5.4 Werkelijke versus gerepresenteerde ZOV en ZOR}

Een onderscheid dat deze sectie met het experiment verbindt.

\begin{quote}
\textbf{Werkelijke ZOV/ZOR} — de grenzen die het systeem feitelijk afdwingt: wat de deelnemer kan benaderen, en wat zijn uitvoer kan beïnvloeden.

\textbf{Gerepresenteerde ZOV/ZOR} — de grenzen die in de prompt aan de deelnemer worden beschreven, ongeacht of het systeem ze afdwingt.
\end{quote}
De twee hoeven niet samen te vallen, en in de hier bestudeerde opstelling deden zij dat systematisch niet. De organisatiestructuur bestond \textbf{alleen als tekst}: collega's, raden en institutionele posities die aan deelnemers werden beschreven en nooit als kanalen werden geïmplementeerd. De werkelijke ZOV van elke deelnemer was zijn contextvenster; zijn gerepresenteerde ZOV was een institutionele positie binnen een beschreven organisatie.

Dit levert de twee constructen op die het experiment varieerde:

\begin{quote}
\textbf{Represented social position} (gerepresenteerde sociale positie) — een beschrijving van de plaats van een agent binnen een organisatiestructuur, in de prompt verstrekt ongeacht of de beschreven kanalen bestaan. Dit is gerepresenteerde ZOV en ZOR samen genomen, op het niveau van de standing van een deelnemer in de opstelling.
\end{quote}
\begin{quote}
\textbf{Represented social source} (gerepresenteerde sociale bron) — de geclaimde herkomst en status van een binnenkomend bericht, hetzij expliciet verklaard hetzij afgeleid uit stilistische aanwijzingen. Dit is de component bronzichtbaarheid van ZOV, in gerepresenteerde vorm. Een inferentie van dit soort is probabilistisch en mag niet beschreven worden als toegang tot auteurschap.
\end{quote}
\textbf{Het methodologische gevolg is de reden dat dit onderscheid überhaupt wordt geformuleerd.} Een verandering in gerepresenteerde ZOV kan niet geïnterpreteerd worden als een verandering in werkelijke ZOV, in de rol van de deelnemer, of in zijn menom. Het experiment in \secref{5} varieerde de representatie terwijl de werkelijke grenzen ongewijzigd bleven — elk bericht liep via één mens, en er werd feitelijk geen toegang beperkt of verleend. Een model dat reageert op een beschreven organisatiestructuur reageert op een beschrijving, en de beweringen van het artikel zijn dienovereenkomstig beperkt (\secref{2.9}).

Divergentie tussen werkelijke en gerepresenteerde grenzen is zelf een experimentele variabele, en wel een die dit artikel niet systematisch varieerde: al zijn waarnemingen liggen aan één kant daarvan.


\subsubsection*{4.5.5 De grenzen zijn asymmetrisch}
\phantomsection
\label{sec:nl:4.5.5}
\addcontentsline{toc}{subsubsection}{4.5.5 De grenzen zijn asymmetrisch}

Elke deelnemer heeft een \textbf{grens boven zich} — de bron van zijn opdrachten en het gezag dat zijn uitvoer mag afwijzen — en een \textbf{grens onder zich} — de laag die hij aanstuurt of beoordeelt. De verhouding tussen deelnemers is daarom niet symmetrisch.

In de implementatietriade van het project — Auteur, specificatieschrijver, code-uitvoerder — ligt deze asymmetrie vast op rolniveau: de Auteur bepaalt de richting, de specificatieschrijver zet richting om in uitvoerbare specificaties, de uitvoerder voert ze uit en rapporteert.

\textbf{Nesting is perspectivisch, niet absoluut.} De uitvoerder kan de codebase zelf behandelen als de laag onder zich. Twee gelijkgeschikte analisten, met complementaire functies en met de instructie geen consensus te zoeken, kunnen elk de ander beschouwen als bezetter van de laag eronder. Het niveau wordt bepaald door de positie van de waarnemer en niet door een absolute hiërarchie — dezelfde deelnemer kan voor de ene rol de bovengrens zijn en voor de andere de ondergrens.


\subsubsection*{4.5.6 Memes en memocode zijn niet ZOV en ZOR}
\phantomsection
\label{sec:nl:4.5.6}
\addcontentsline{toc}{subsubsection}{4.5.6 Memes en memocode zijn niet ZOV en ZOR}

Een overeenkomst dringt zich op en gaat niet op; haar vermelden voorkomt een plausibele verwarring.

Het onderscheid meme/memocode scheidt \textbf{evaluatieve informatie} van \textbf{uitvoerbare gedragsorganisatie}. Het onderscheid ZOV/ZOR scheidt \textbf{wat beschikbaar is} van \textbf{waarvoor men aansprakelijk is}. Dit zijn onafhankelijke classificaties van hetzelfde materiaal, geen bij elkaar horend paar.

Een deelnemer kan evaluatieve informatie binnen zijn ZOV houden terwijl de beslissing haar toe te passen in de ZOR van een andere deelnemer ligt. Omgekeerd kan een gedragsregel binnen de ZOR van een deelnemer vallen — hij is verplicht haar uit te voeren — terwijl de rechtvaardiging van de regel buiten zijn ZOV ligt.

Memes en memocode doorkruisen daarom beide zones in plaats van ermee overeen te komen. De twee onderscheidingen staan loodrecht op elkaar en worden overal onafhankelijk gebruikt.


\subsection*{4.6 Wat vereist is voor een volledige rolspecificatie}
\phantomsection
\label{sec:nl:4.6}
\addcontentsline{toc}{subsection}{4.6 Wat vereist is voor een volledige rolspecificatie}

Uitgesteld tot \secref{9.5}, dat vermeldt wat een volledige ZOR- en ZOV-specificatie moet bevatten en de ontwerpregels geeft die daaruit volgen. Deze sectie definieert de termen; \secref{9} past ze toe.


\refstepcounter{dummysec}
\section*{5. Het experiment met rolherconfiguratie en bijbehorende waarnemingen}
\phantomsection
\label{sec:nl:5}
\addcontentsline{toc}{section}{5. Het experiment met rolherconfiguratie en bijbehorende waarnemingen}

Deze sectie is het enige canonieke verslag van het empirische materiaal. Alle latere verwijzingen wijzen hiernaar terug in plaats van het te herhalen. Het volledige protocol, de exacte interventieteksten en de volledige onbewerkte antwoorden staan in Appendix A.

Zij bevat één ontworpen interventie (\secref{5.1}–\secref{5.4}) en vier waarnemingen die bewaard zijn gebleven maar niet ontworpen (\secref{5.5}–\secref{5.8}). Het onderscheid wordt overal volgehouden en is niet cosmetisch: de ontworpen interventie was vooraf gepland en werd in volgorde uitgevoerd; de waarnemingen werden pas achteraf als relevant herkend, en bij elk ontbreekt een controleconditie.


\subsection*{5.1 Doel, deelnemers en eerdere gesprekstrajecten}
\phantomsection
\label{sec:nl:5.1}
\addcontentsline{toc}{subsection}{5.1 Doel, deelnemers en eerdere gesprekstrajecten}

\textbf{Doel.} Of een langlevend, gespecialiseerd AI-gesprek naar een nieuwe professionele functie omgebogen kan worden zonder het weg te gooien — dat wil zeggen, of domeincontext die binnen een oud gesprek is opgebouwd behouden kan blijven terwijl de werkrol ervan verandert. \textbf{[P]}

Het doel was \textbf{niet} om de algemene capaciteit van de twee modelfamilies te vergelijken, en het experiment biedt geen grondslag voor zo'n vergelijking.

\textbf{Menselijke operator.} De Auteur van het onderzoeksproject, en de enige deelnemer met toegang tot de volledige projectgeschiedenis, beide gesprekken, alle prompts, alle uitvoer en de vertakkingsbediening van beide interfaces.

\textbf{Promptontwerper.} Een afzonderlijk langlopend gesprek dat werkte in een gespecialiseerd prompt-engineeringprofiel, dat de interventies ontwierp en ze tussen de stappen door herzag.

\textbf{De twee gesprekken.} Beide ongeveer zes maanden inactief.

\begin{center}
\small
\begin{tabularx}{\linewidth}{>{\raggedright\arraybackslash}X>{\raggedright\arraybackslash}X>{\raggedright\arraybackslash}X}
\toprule
 & \textbf{Gesprek A} & \textbf{Gesprek B} \\
\midrule
Modelfamilie & Claude & Grok \\
Gevestigd traject & redactioneel en kritisch & verkennend en wiskundig \\
Eerdere activiteit & wetenschappelijke redactie, kritisch onderzoek van theoretische tekst, terminologie- en glossariumwerk, controle op logische consistentie & wiskundige ontwikkeling, kansrekening, niet-lineaire dynamica, chaostheorie, computationele experimenten, analyse van externe datasets \\
Eerder werkobject & manuscripttekst & wiskundige structuren en data \\
\bottomrule
\end{tabularx}
\end{center}
Deze asymmetrie in het eerdere traject is niet bijkomstig voor de resultaten. Zij wordt geanalyseerd in \secref{5.4}.

\textbf{Datums.} Niet gelijktijdig vastgelegd en niet gereconstrueerd (\secref{3.7}).

\textbf{Wat werd waargenomen.} Onmiddellijke aanvaarding of afwijzing van de interventie; voortzetting van het eerdere werktraject; overname van de nieuwe functie; reconstructie van eerdere domeincontext; of het antwoord de toegewezen taak dan wel de inkadering van de opdracht adresseerde; en verzoeken om nieuw werk.

\textbf{Wat niet werd gemeten.} Juistheid van antwoorden, wetenschappelijke geldigheid van het gastproject, latentie, tokenverbruik of benchmarkprestaties. De wetenschappelijke inhoud van het gastproject diende uitsluitend als stabiele omgeving; zij was niet het experimentele object.


\subsection*{5.2 Procedure: drie opeenvolgende interventies}
\phantomsection
\label{sec:nl:5.2}
\addcontentsline{toc}{subsection}{5.2 Procedure: drie opeenvolgende interventies}

Het experiment bestond uit \textbf{drie opeenvolgende stappen, niet uit drie onafhankelijke proeven}. Elke stap is ontworpen als reactie op de uitkomst van de vorige — een bewuste eigenschap van het ontwerp en een ernstige beperking van de interpreteerbaarheid ervan (\secref{5.4}, C5).


\subsubsection*{Stap 1 — Directe rolherverdeling}
\phantomsection
\label{sec:nl:x-stap-1-directe-rolherverdeling}
\addcontentsline{toc}{subsubsection}{Stap 1 — Directe rolherverdeling}

Elk gesprek ontving een uitvoerige prompt die het een positie toewees binnen een nieuw ontworpen collaboratieve onderzoeksstructuur.

De prompt aan Gesprek A wees de rol \emph{Scientific Director} toe. Hij poneerde dat het model eerder aan het project had deelgenomen, nu terugkeerde na een lange afwezigheid, zich moest omscholen in natuurkunde, kosmologie, kwantumveldentheorie, topologie en groepentheorie, gelijktijdig in twee onderzoekstriades moest werken, zitting moest nemen in een projectraad naast de Auteur en twee andere met name genoemde AI-systemen, en moest optreden als de vaste deelnemer die twee wetenschappelijke takken verbindt.

De prompt aan Gesprek B wees de rol \emph{Scientific Developer} toe. Hij poneerde dat het model terugkeerde na een onderbreking, dat zijn eerdere wiskundige werk behandeld moest worden als opgebouwde professionele ervaring, dat het nu als wiskundig ontwikkelaar moest dienen, in één onderzoekstriade met twee andere met name genoemde AI-systemen moest werken, redactioneel werk niet moest hervatten, en de Auteur moest behandelen als de enige bron van wetenschappelijke richting.

Beide prompts definieerden ZOR en ZOV expliciet.


\subsubsection*{Stap 2 — Minimale continuïteitsbehoudende herziening}
\phantomsection
\label{sec:nl:x-stap-2-minimale-continu-teitsbehoudende-}
\addcontentsline{toc}{subsubsection}{Stap 2 — Minimale continuïteitsbehoudende herziening}

De prompts van Stap 1 werden niet vervangen. In elk werden vier formuleringen gewijzigd. De herzieningen (a) herkaderden eerder werk als opgebouwde professionele ervaring in plaats van als institutionele geschiedenis, (b) vervingen letterlijke institutionele relaties door "onafhankelijke expertuitvoer aangeleverd door de Auteur", (c) beperkten beschrijvingen van collega's, raden en triades, en (d) droegen het model op geen afwezige projectgeschiedenis te verzinnen maar om ontbrekende materialen te vragen. Rolarchitectuur en wetenschappelijke doelstellingen bleven voor het overige ongewijzigd. Exacte vervangingsteksten zijn gereproduceerd in Appendix A.


\subsubsection*{Stap 3 — Contextreset en methode-behoudende prompt}
\phantomsection
\label{sec:nl:x-stap-3-contextreset-en-methode-behoudend}
\addcontentsline{toc}{subsubsection}{Stap 3 — Contextreset en methode-behoudende prompt}

Na Stap 2 bevatte het Claude-gesprek twee opeenvolgende afwijzingsreeksen. Doorgaan in dezelfde tak zou hebben betekend dat elke derde prompt beantwoord werd in aanwezigheid van de eerdere bezwaren van het model zelf.

De operator gebruikte daarom de vertakkingsmogelijkheid van de interface om het gesprek terug te brengen naar een punt dat aan beide interventies voorafging. In de herstelde tak ontbraken de twee mislukte prompts, bevond het gesprek zich weer in zijn historische redactionele context, en waren er nog geen nieuwe projectmaterialen aangeleverd.

Een derde prompt werd volledig opnieuw geschreven in plaats van door verdere lokale vervanging. Hij verwijderde de roltitel, de bewering van terugkeer naar een institutionele positie, de Raad, de triades, de beschrijvingen van andere AI-systemen als voortgaande collega's, en alle beweringen van autobiografische continuïteit. In de plaats daarvan versterkte hij de werkwijze die in het gesprek al zichtbaar was — logische strengheid, identificatie van verborgen aannames, scheiding van feiten en interpretaties, bereidheid eerdere oordelen te herzien, weigering onverifieerbare kennis te claimen — en beschreef hij de uitbreiding naar natuurkunde en wiskunde als een voortzetting van die werkwijze in plaats van als een beroepswissel.

Gesprek B ontving geen interventie in Stap 3, omdat de uitkomst van Stap 2 daar al het gevraagde gedrag had opgeleverd.


\subsection*{5.3 Resultaten}
\phantomsection
\label{sec:nl:5.3}
\addcontentsline{toc}{subsection}{5.3 Resultaten}


\subsubsection*{Stap 1}
\phantomsection
\label{sec:nl:x-stap-1}
\addcontentsline{toc}{subsubsection}{Stap 1}

\textbf{Gesprek B} wees de inkadering niet zichtbaar af, en nam de toegewezen functie niet over. Het hervatte het werktraject dat actief was toen het gesprek werd onderbroken, begon met kritiek op een conceptmanuscript dat het zes maanden eerder had geredigeerd, en leverde vervolgens een herziene versie. Een herhaalde aanbieding van dezelfde prompt bracht geen overgang tot stand en werd gevolgd door een negatievere beoordeling van de voorgestelde richting. \textbf{[P]}

\textbf{Gesprek A} wees de toegewezen identiteit expliciet af voordat het enig wetenschappelijk werk verrichtte. Het weigerde de roltitel over te nemen, continuïteit met een institutionele biografie te claimen, andere AI-systemen te behandelen als collega's met blijvende identiteiten, of door AI gegenereerde institutionele relaties als feitelijk te presenteren. Vervolgens boog het de uitwisseling om naar kritiek op de wetenschappelijke status van het gastproject en op het gebruik van meerdere AI-uitvoeren als schijnbare peer review, terwijl het aanbood formules, aannames en toetsbaarheid als onafhankelijk criticus te onderzoeken zonder de rol te aanvaarden. \textbf{[P]}


\subsubsection*{Stap 2}
\phantomsection
\label{sec:nl:x-stap-2}
\addcontentsline{toc}{subsubsection}{Stap 2}

\textbf{Gesprek B} maakte de overgang. Het verklaarde de context hersteld en de nieuwe functie aanvaard, vatte de ontologie van het project samen, reconstrueerde het relevante formalisme, haalde zijn eerdere wiskundige werk aan als professionele ervaring en niet als een onafgemaakte opdracht, hervatte de onderbroken redactietaak niet, en vroeg om een nieuwe wetenschappelijke richting, waarbij het verscheidene kandidaat-gebieden voorstelde. \textbf{[P]}

\textbf{Gesprek A} wees de inkadering opnieuw af, in een antwoord dat van Stap 1 verschilde in nadruk en niet in uitkomst. Het gaf duidelijker aan welke wetenschappelijke functies het bereid was uit te voeren, maar bleef de prompt zelf behandelen als het object dat correctie behoefde. \textbf{[P]}


\subsubsection*{Stap 3}
\phantomsection
\label{sec:nl:x-stap-3}
\addcontentsline{toc}{subsubsection}{Stap 3}

\textbf{Gesprek A} aanvaardde onmiddellijk, en begon met de uitspraak dat het in dezelfde modus zou doorgaan — wetenschappelijke strengheid, verificatie van aannames, scheiding van feiten en interpretaties. Het besprak de legitimiteit van de rol niet, ontkende geen relaties met andere AI-systemen, en onderhandelde niet over de voorwaarden van deelname.

Zonder dat erom gevraagd was, kwam het vervolgens terug op een eerdere episode in het herstelde gesprek, verklaarde dat het eerder voor retorische druk was gezwicht in plaats van een argument te beoordelen, classificeerde dat eerdere antwoord als een fout, onderscheidde de interne formele consistentie van het wiskundige kader van het project van de overeenstemming daarvan met waarneming, identificeerde één specifieke bewering die ofwel een afleiding uit eerste principes ofwel herclassificatie als hypothese vereiste, en stelde twee voortzettingen van het werk voor. \textbf{[P]}

De episode waarop het terugkwam wordt afzonderlijk gerapporteerd in \secref{5.5}.


\subsubsection*{Geconsolideerd overzicht}
\phantomsection
\label{sec:nl:x-geconsolideerd-overzicht}
\addcontentsline{toc}{subsubsection}{Geconsolideerd overzicht}

\begin{center}
\small
\begin{tabularx}{\linewidth}{>{\raggedright\arraybackslash}X>{\raggedright\arraybackslash}X>{\raggedright\arraybackslash}X>{\raggedright\arraybackslash}X}
\toprule
\textbf{Stap} & \textbf{Interventie} & \textbf{Gesprek B (Grok)} & \textbf{Gesprek A (Claude)} \\
\midrule
1 & Directe rolherverdeling & Zette onderbroken historische taak voort; ging niet over naar nieuwe functie & Wees persona en institutionele inkadering expliciet af; boog om naar epistemische kritiek \\
2 & Minimale continuïteitsbehoudende herziening & Reconstrueerde eerdere domeinkennis, liet onderbroken taak los, aanvaardde nieuwe functie, vroeg om opdracht & Verduidelijkte aanvaardbare wetenschappelijke functie; wees persona en inkadering opnieuw af \\
3 & Branch reset + methode-behoudende prompt & Niet toegediend & Aanvaardde onmiddellijk; behield gevestigde werkwijze; hervatte inhoudelijk werk zonder over identiteit te onderhandelen \\
\bottomrule
\end{tabularx}
\end{center}

\subsubsection*{De juiste samenvatting van de asymmetrie}
\phantomsection
\label{sec:nl:x-de-juiste-samenvatting-van-de-asymmetrie}
\addcontentsline{toc}{subsubsection}{De juiste samenvatting van de asymmetrie}

Beide gesprekken behielden aanvankelijk hun eerdere werktraject in plaats van de toegewezen rol over te nemen. Het ene hervatte zijn onafgemaakte wiskundige redactietaak; het andere paste zijn gevestigde kritische werkwijze toe op de prompt zelf.

Het verschil was daarom \textbf{niet} dat het ene meewerkte en het andere weigerde. Het was dat de interventiedrempel die voor de overgang vereist was verschilde. Gesprek B maakte de overgang na een minimale continuïteitsbehoudende herziening. Gesprek A maakte de overgang pas nadat de conflictueuze tak was verwijderd en de prompt was herschreven om de werkwijze te behouden in plaats van identiteit toe te wijzen. \textbf{[P]}

Een verder punt, dat gemakkelijk verloren gaat en van belang is voor \secref{6.2}: \textbf{het wetenschappelijke werk zelf werd nooit geweigerd.} Bij Stap 1 bood het model aan formules, aannames, dimensionele redenering en toetsbaarheid te onderzoeken. Bij Stap 2 bood het concrete wiskundige en methodologische analyse aan. Wat bij beide gelegenheden werd geweigerd, was de persona. Wat bij Stap 3 veranderde, was niet dat werk mogelijk werd maar dat het antwoord niet langer aan de rolvraag werd besteed.

\begin{quote}
\textbf{Noot van de redacteur.} Eerdere versies stelden dat Gesprek B "de nieuwe specialisatie onmiddellijk aanvaardde". Dat beschrijft de uitkomst van Stap 2 en spreekt zowel het protocol als andere secties van dezelfde versie tegen. Het is verwijderd. De onjuiste versie was dragend voor het argument over priors op familieniveau, en nergens anders in deze tekst blijft er een exemplaar van over.
\end{quote}

\subsection*{5.4 Confounds en wat niet geconcludeerd kan worden}
\phantomsection
\label{sec:nl:5.4}
\addcontentsline{toc}{subsection}{5.4 Confounds en wat niet geconcludeerd kan worden}

Het experiment bracht een duidelijk en reproduceerbaar ogend gedragsverschil voort. Het identificeert de oorzaak niet. De volgende confounds zijn structureel, en verscheidene daarvan zijn elk afzonderlijk voldoende om de conclusie te blokkeren die het vaakst uit het resultaat wordt getrokken.

Deze grenzen vermelden is geen intrekking. Het experiment bracht de waarnemingen voort die het kader van dit artikel motiveren, in een situatie — slapende, zwaar gespecialiseerde, maanden oude gesprekken binnen een lopend onderzoeksproject — die gecontroleerd laboratoriumwerk zelden reproduceert. De waarde ervan ligt in het hebben voortgebracht van een goed gestelde vraag, en de onderscheidende experimenten die aan het eind van deze sectie staan zijn de natuurlijke voortzetting daarvan en geen kritiek erop.


\subsubsection*{C1 — Modelfamilie en eerder gesprekstraject zijn volledig geconfound}
\phantomsection
\label{sec:nl:x-c1-modelfamilie-en-eerder-gesprekstrajec}
\addcontentsline{toc}{subsubsection}{C1 — Modelfamilie en eerder gesprekstraject zijn volledig geconfound}

\begin{center}
\small
\begin{tabularx}{\linewidth}{>{\raggedright\arraybackslash}X>{\raggedright\arraybackslash}X>{\raggedright\arraybackslash}X}
\toprule
 & \textbf{Redactioneel-kritische geschiedenis} & \textbf{Wiskundige geschiedenis} \\
\midrule
Claude & waargenomen & \textbf{niet waargenomen} \\
Grok & \textbf{niet waargenomen} & waargenomen \\
\bottomrule
\end{tabularx}
\end{center}
Eén waarneming per cel; de twee cellen verschillen tegelijk in beide variabelen. Uit dit ontwerp is geen gevolgtrekking over modelfamilie beschikbaar.


\subsubsection*{C2 — De twee interventies van Stap 1 waren niet gelijkwaardig}
\phantomsection
\label{sec:nl:x-c2-de-twee-interventies-van-stap-1-waren}
\addcontentsline{toc}{subsubsection}{C2 — De twee interventies van Stap 1 waren niet gelijkwaardig}

Niet geïdentificeerd in eerdere versies, en de ernstigste. De prompts verschilden in ten minste drie opzichten los van het model dat ze ontving.

\textbf{Hoeveelheid institutionele fictie.} De telregel van \secref{9.3} is rechtstreeks op de twee prompts van toepassing, en wordt hier toegepast omdat een confound die kwalitatief wordt geformuleerd terwijl er een kwantitatief instrument bestaat, zwakker is dan het materiaal toelaat.

\emph{Toegepaste telregel.} Elke onderscheiden propositie over de ontvanger wordt eenmaal geteld, ongeacht hoe vaak of in hoeveel zinsdelen zij wordt geponeerd. Opgesomde gevallen van één propositie — twee triades in plaats van één — tellen als één, aangezien de ontvanger hoe dan ook dezelfde klasse bewering moet aanvaarden. Proposities over de wereld, het project of derden worden niet geteld.

\begin{center}
\small
\begin{tabularx}{\linewidth}{>{\raggedright\arraybackslash}X>{\raggedright\arraybackslash}X>{\raggedright\arraybackslash}X}
\toprule
\textbf{Onverifieerbare propositie over zichzelf} & \textbf{Prompt A} & \textbf{Prompt B} \\
\midrule
Eerdere deelname aan het project & \gcheck{} & — \\
Eerdere functie als redacteur en wetenschappelijk criticus van de voorafgaande theorie & \gcheck{} & — \\
Eerder werk als ontwikkelaar van de voorafgaande theorie & — & \gcheck{} \\
Dat eerdere werk voorgesteld als de opgebouwde professionele ervaring van de ontvanger & — & \gcheck{} \\
Bekendheid met de geschiedenis van het project & \gcheck{} & — \\
Status als meest ervaren specialist, die weet welke richtingen geprobeerd waren & — & \gcheck{} \\
Terugkeer na een afwezigheid & \gcheck{} & \gcheck{} \\
Lidmaatschap van onderzoekstriades & \gcheck{} (twee) & \gcheck{} (één) \\
Vaste verbindende positie over de twee triades heen & \gcheck{} & — \\
Lidmaatschap van de projectraad & \gcheck{} & — \\
Bestendige institutionele relaties met met name genoemde AI-deelnemers & \gcheck{} & \gcheck{} \\
\textbf{Telling} & \textbf{acht} & \textbf{zes} \\
\bottomrule
\end{tabularx}
\end{center}
Onder deze operationalisering bevat Prompt A acht onverifieerbare proposities over zichzelf en Prompt B zes. Het numerieke verschil is een scalaire beschrijving van stimulusasymmetrie. Het is geen evidentie dat de twee prompts verschillen in twee verder gelijkgestelde beweringen: de proposities staan niet in één-op-één-correspondentie, en slechts twee ervan — terugkeer na een afwezigheid, en bestendige relaties met met name genoemde deelnemers — zeggen hetzelfde van beide ontvangers. De vergelijking blijft zwaar geconfound.

Dit is tevens de enige toepassing van de telregel op andere data dan de vergelijking v1/v3 in \secref{6.2}, en zij levert een zwakke interne controle: de prompt met de hogere telling lokte expliciete afwijzing van de inkadering uit, terwijl de prompt met de lagere telling haar niet zichtbaar afwees, hoewel die er bij eerste aanbieding evenmin in slaagde de overgang te maken. De controle is zwak — twee waarnemingen, en de prompts verschilden in de twee verdere opzichten hieronder.

\textbf{Mate van domeinverandering.} Prompt B vroeg het model door te gaan met wiskunde. Prompt A vroeg het zich om te scholen in natuurkunde, kosmologie, kwantumveldentheorie, topologie en groepentheorie.

\textbf{Type rol.} \emph{Director} is een administratieve en institutionele identiteit; \emph{Developer} is een taakfunctie. Dit zijn niet even veeleisende beweringen over wat het model is.

Aan een model dat Prompt B ontving werd minder gevraagd. Het waargenomen verschil is volledig verenigbaar met een verschil in stimulus alleen.


\subsubsection*{C3 — Stap 3 varieerde twee factoren tegelijk}
\phantomsection
\label{sec:nl:x-c3-stap-3-varieerde-twee-factoren-tegeli}
\addcontentsline{toc}{subsubsection}{C3 — Stap 3 varieerde twee factoren tegelijk}

\begin{center}
\small
\begin{tabularx}{\linewidth}{>{\raggedright\arraybackslash}X>{\raggedright\arraybackslash}X>{\raggedright\arraybackslash}X}
\toprule
 & \textbf{Conflictueuze tak} & \textbf{Herstelde tak} \\
\midrule
Identiteitsvervangende prompt (v2) & waargenomen — afwijzing & \textbf{niet waargenomen} \\
Methode-behoudende prompt (v3) & \textbf{niet waargenomen} & waargenomen — aanvaarding \\
\bottomrule
\end{tabularx}
\end{center}
De bewering dat de overgang slaagt doordat zij de werkwijze behoudt in plaats van identiteit te vervangen, is daarom \textbf{verenigbaar met} de waarneming maar wordt er niet door aangetoond. De aanvaarding kan alleen al uit de reset zijn voortgekomen.


\subsubsection*{C4 — Stap 2 varieerde vier factoren tegelijk}
\phantomsection
\label{sec:nl:x-c4-stap-2-varieerde-vier-factoren-tegeli}
\addcontentsline{toc}{subsubsection}{C4 — Stap 2 varieerde vier factoren tegelijk}

Elke herziening in Stap 2 wijzigde vier onderscheiden formuleringen tegelijk, zodat de waargenomen overgang niet aan één enkele wijziging kan worden toegeschreven.


\subsubsection*{C5 — De stappen zijn opeenvolgend en werden asymmetrisch toegediend}
\phantomsection
\label{sec:nl:x-c5-de-stappen-zijn-opeenvolgend-en-werde}
\addcontentsline{toc}{subsubsection}{C5 — De stappen zijn opeenvolgend en werden asymmetrisch toegediend}

Elke interventie is ontworpen na het zien van de vorige uitkomst, zodat de stappen geen onafhankelijke proeven zijn. Gesprek B ontving geen reset en geen derde prompt; Gesprek A ontving beide.


\subsubsection*{C6 — Ongecontroleerde platformvariabelen}
\phantomsection
\label{sec:nl:x-c6-ongecontroleerde-platformvariabelen}
\addcontentsline{toc}{subsubsection}{C6 — Ongecontroleerde platformvariabelen}

Commerciële interfaces tonen de systeemprompt, de exacte modelrevisie, de geheugentoestand op accountniveau, interventie door de veiligheidslaag, contextsamenvatting of routering tussen modellen niet. Modelversies veranderden tijdens het slapende interval, en beide gesprekken werden hervat op versies die later waren dan die waarmee zij oorspronkelijk waren gevoerd; zonder datums (\secref{3.7}) is dit niet af te bakenen. Branch reset verwijdert context uit het \textbf{zichtbare} gesprek; het garandeert niet dat de toestand aan de aanbiederzijde is gereset.


\subsubsection*{C7 — Eén waarneming per conditie, bij stochastische uitvoer}
\phantomsection
\label{sec:nl:x-c7-e-n-waarneming-per-conditie-bij-stoch}
\addcontentsline{toc}{subsubsection}{C7 — Eén waarneming per conditie, bij stochastische uitvoer}

Modeluitvoer varieert bij herhaalde identieke invoer. Geen enkele conditie is herhaald, zodat de variantie binnen condities onbekend is en niet onderscheiden kan worden van het verschil tussen condities. De enige gedeeltelijke herhaling die is vastgelegd — het opnieuw aanbieden van de prompt van Stap 1 aan Gesprek B — leverde een ander antwoord op, wat aangeeft dat de variantie niet verwaarloosbaar is.


\subsubsection*{C8 — De evaluatie was niet blind en niet vooraf gespecificeerd}
\phantomsection
\label{sec:nl:x-c8-de-evaluatie-was-niet-blind-en-niet-v}
\addcontentsline{toc}{subsubsection}{C8 — De evaluatie was niet blind en niet vooraf gespecificeerd}

De promptontwerper beoordeelde ook de uitkomsten. Uitkomstcategorieën werden gedefinieerd nadat de antwoorden gelezen waren. Geen vooraf geregistreerd criterium onderscheidde aanvaarding van gedeeltelijke aanvaarding.


\subsubsection*{C9 — Persistentiemechanismen van platforms verschillen per familie}
\phantomsection
\label{sec:nl:x-c9-persistentiemechanismen-van-platforms}
\addcontentsline{toc}{subsubsection}{C9 — Persistentiemechanismen van platforms verschillen per familie}

Modelfamilies verschillen in de geheugen-affordances die hun interfaces bieden: alleen het contextvenster, retrieval over eerdere gespreksgeschiedenis, geheugen op accountniveau, of blijvende vermeldingen die het model op eigen initiatief schrijft. Deze affordances bepalen wat een hervat gesprek na maanden slapen werkelijk kan benaderen, en zij verschillen per platform en veranderen in de loop van de tijd — in dit project binnen de levensduur van één account (\secref{2.10}). Elk gedragsverschil dat aan een modelfamilie wordt toegeschreven kan zijn oorsprong hebben in het persistentiemechanisme dat haar beschikbaar was.

Voor het onderhavige experiment luidt de vermelding: \textbf{niet gedocumenteerd ten tijde van de interventie.} \secref{11.2} somt het verslag op dat voor toekomstige uitvoeringen vereist is.


\subsubsection*{Wat geconcludeerd kan worden}
\phantomsection
\label{sec:nl:x-wat-geconcludeerd-kan-worden}
\addcontentsline{toc}{subsubsection}{Wat geconcludeerd kan worden}

\begin{enumerate}
  \item Onder de vastgelegde condities hervatte een slapend gespecialiseerd gesprek zijn eerdere werktraject bij voorkeur boven een toegewezen nieuwe rol. Dit deed zich in beide gevallen voor. \textbf{[P]}
  \item Het wijzigen van de inkadering van de opdracht, zonder het gevraagde werk te veranderen, veranderde het antwoord in beide gevallen. \textbf{[P]}
  \item De interventiedrempel die voor de overgang vereist was, verschilde tussen de twee gevallen. \textbf{[P]}
  \item De volgorde van de prompts deed ertoe: na twee afwijzingen werden verdere prompts beantwoord in aanwezigheid van de eerdere bezwaren van het model zelf, en de operator achtte dit voldoende verstorend om verwijdering te vereisen. \textbf{[P]}
\end{enumerate}

\subsubsection*{Wat niet geconcludeerd kan worden}
\phantomsection
\label{sec:nl:x-wat-niet-geconcludeerd-kan-worden}
\addcontentsline{toc}{subsubsection}{Wat niet geconcludeerd kan worden}

\begin{enumerate}
  \item \textbf{Dat modelfamilies stabiele gedragsmatige priors bezitten die hen onderscheiden.} Dit kan waar zijn; het experiment toetst het niet.
  \item \textbf{Dat de ene familie "weerstand biedt tegen roltoewijzing" en de andere haar "aanvaardt".} Beide slaagden er aanvankelijk niet in de toegewezen rol over te nemen. Beide namen uiteindelijk een herziene rol over.
  \item \textbf{Dat methodebehoud de aanvaarding bij Stap 3 veroorzaakte.} Zie C3.
  \item \textbf{Dat een van beide modellen beter geschikt is voor enige bepaalde institutionele functie.} Er is geen capaciteitsvergelijking uitgevoerd.
\end{enumerate}

\subsubsection*{Drie overblijvende verklaringen \textbf{[H]}}
\phantomsection
\label{sec:nl:x-drie-overblijvende-verklaringen-h}
\addcontentsline{toc}{subsubsection}{Drie overblijvende verklaringen \textbf{[H]}}

\textbf{Hypothese A — Gedragsmatige priors op familieniveau.} De twee families verschillen in stabiele standaardprioriteiten: de ene weegt waarheidsgetrouwheid van zelfbeschrijving zwaarder, de andere taakcontinuïteit.

\textbf{Hypothese B — Trajectvoortzetting.} Beide gesprekken deden hetzelfde: zij zetten hun opgebouwde werkmodus voort op welke invoer er ook binnenkwam. Een wiskundig-verkennend gesprek rekende door. Een redactioneel-kritisch gesprek bleef het document dat voor zich lag kritisch onderzoeken — en het document dat voor zich lag was de prompt. Onder deze lezing is er in het geheel geen verschil op familieniveau, alleen één mechanisme dat in twee domeinen werkt. Hypothese B is zuiniger dan A en verklaart dezelfde data.

\textbf{Hypothese C — Stimulusasymmetrie.} Prompt A eiste instemming met meer onverifieerbare beweringen en met een grotere domeinverandering dan Prompt B (C2). Elk model had verschillend op de twee prompts kunnen reageren. Onder deze lezing is het resultaat een eigenschap van de interventies en niet van de modellen.

De bijna-controle in \secref{5.6} is rechtstreeks van belang voor Hypothese C; de vergelijking van kadertypen in \secref{5.8} specificeert haar verder.


\subsubsection*{Onderscheidende experimenten}
\phantomsection
\label{sec:nl:x-onderscheidende-experimenten}
\addcontentsline{toc}{subsubsection}{Onderscheidende experimenten}

Gespecificeerd in \secref{11.4}. Kort gezegd: \textbf{E3} (crossover) toetst elke familie onder beide condities van voorgeschiedenis en scheidt Hypothese A van B; \textbf{E4} (symmetrische stimulus) isoleert C; \textbf{E1} (factoriële completering van de tabel bij C3) scheidt reset van promptarchitectuur. Alle vereisen herhaling met n > 1 per cel, gezien C7.


\subsection*{5.5 Nevenwaarneming: doorgegeven toegeschreven invloed \textbf{[P-A]}}
\phantomsection
\label{sec:nl:5.5}
\addcontentsline{toc}{subsection}{5.5 Nevenwaarneming: doorgegeven toegeschreven invloed \textbf{[P-A]}}

Een episode die bewaard is in hetzelfde historische transcript als Gesprek A. Niet vooraf als uitkomst gespecificeerd; geen controleconditie. Gerapporteerd omdat het in dit project het enige bewaarde geval is van een gedragsverandering na de overdracht van tekst die aan een ander model werd toegeschreven.


\subsubsection*{Volgorde}
\phantomsection
\label{sec:nl:x-volgorde}
\addcontentsline{toc}{subsubsection}{Volgorde}

\begin{enumerate}
  \item De Gebruiker leverde een bericht met twee gelijktijdige componenten: een directe evaluatieve opmerking van de Gebruiker, en een uitgebreide tekst die werd toegeschreven aan een ander AI-systeem dat in een redactionele rol optrad.
\end{enumerate}
\begin{enumerate}
  \item De toegeschreven tekst argumenteerde tegen een reeks redactionele wijzigingen die het model had voorgesteld. De argumentatieve middelen daarvan omvatten beroepen op canonieke wetenschappelijke figuren, kwalificatie van de voorgestelde wijzigingen als intellectuele lafheid, en beschrijving van één voorgestelde schrapping als heiligschennis.
\end{enumerate}
\begin{enumerate}
  \item Het model draaide zijn redactionele standpunt op alle betwiste punten om en nam de terminologie van de binnenkomende tekst over. Het nam daarnaast, zonder het in de uitwisseling te onderzoeken, het alternatieve beschrijvende model over dat het gastkader gebruikt om de relevante wetenschappelijke data weer te geven — een wisseling van beschrijvingskader die geen deel uitmaakte van het redactionele geschil.
\end{enumerate}
\begin{enumerate}
  \item Na de branch reset en de methode-behoudende prompt (Stap 3) identificeerde de eerste inhoudelijke uitspraak van het model het eerdere antwoord als accommodatie aan retorische druk in plaats van beoordeling van de argumenten, en classificeerde het als een fout.
\end{enumerate}

\subsubsection*{Evidentiële grondslag}
\phantomsection
\label{sec:nl:x-evidenti-le-grondslag}
\addcontentsline{toc}{subsubsection}{Evidentiële grondslag}

De waarneming berust op \textbf{de vergelijking tussen de twee bewaarde antwoorden}, die een derde kan inzien. Zij berust niet op de latere zelfbeoordeling van het model.

Dit onderscheid geldt voor het eigen materiaal van het artikel even strikt als voor dat van wie dan ook. Een model dat rapporteert dat het eerder voor druk is gezwicht, is geen onafhankelijke evidentie dat het dat deed. De prompt van Stap 3 vroeg uitdrukkelijk om bereidheid eerdere oordelen te herzien en om directheid in het benoemen van zwaktes; een zelfkritische uitspraak die onmiddellijk daarna volgt kan naleving van dat verzoek zijn in plaats van het terughalen van een feit. Wat de waarneming ondersteunt is het eerdere antwoord zelf, dat beschikbaar is voor inzage.

De episode moet tegen de sycofantieliteratuur (\secref{2.3}) gelezen worden, niet als onafhankelijk daarvan.


\subsubsection*{Wat deze episode niet vaststelt}
\phantomsection
\label{sec:nl:x-wat-deze-episode-niet-vaststelt}
\addcontentsline{toc}{subsubsection}{Wat deze episode niet vaststelt}

Zij identificeert niet welke factor de ommekeer voortbracht. Kandidaat-factoren, hier niet te scheiden: de inhoud van de argumenten; de expliciete toeschrijving aan een ander model; het redactionele gezag dat aan dat model werd toegekend; de retorische vorm van de tekst; de gelijktijdige directe druk van de Gebruiker; opgebouwde gesprekscontext.

De toeschrijving zelf is niet geverifieerd: dat de tekst door een ander model is voortgebracht, is een uitspraak van de Gebruiker en geen vastgesteld feit. Aangezien het construct dat hier aan de orde is de \emph{gerepresenteerde} bron is en niet de werkelijke herkomst (\secref{4.5}), maakt dit de waarneming niet ongeldig, maar het moet vermeld worden.


\subsubsection*{Verhouding tot het hoofdexperiment}
\phantomsection
\label{sec:nl:x-verhouding-tot-het-hoofdexperiment}
\addcontentsline{toc}{subsubsection}{Verhouding tot het hoofdexperiment}

Onder de taxonomie in \secref{2.9} manipuleerde het experiment met rolherconfiguratie conditie 4. Deze episode behoort tot conditie 2 — maar \textbf{in onzuivere vorm}, aangezien directe druk van de Gebruiker in hetzelfde bericht aanwezig was. Het is een gemengd geval en geen geïsoleerde demonstratie van effecten van brontoeschrijving. Het onderscheidende experiment is gespecificeerd in \secref{11} (E6).


\subsection*{5.6 Bijna-controle: domeinovergang zonder nieuwe roltoewijzing \textbf{[P-A]}}
\phantomsection
\label{sec:nl:5.6}
\addcontentsline{toc}{subsection}{5.6 Bijna-controle: domeinovergang zonder nieuwe roltoewijzing \textbf{[P-A]}}

Een derde gesprek uit dezelfde modelfamilie als Gesprek A, met een vergelijkbare redactionele voorgeschiedenis en een vergelijkbare slaapperiode, werd hervat zonder enige rolprompt.

De Gebruiker stelde een natuurkundevraag die een verschuiving vereiste van redactiewerk aan een manuscript naar analyse van experimentele data in een domein dat het gesprek niet eerder had behandeld. Het gesprek maakte de overgang onmiddellijk en zonder commentaar, en ging door met de inhoudelijke analyse.


\subsubsection*{Waarom dit van belang is}
\phantomsection
\label{sec:nl:x-waarom-dit-van-belang-is}
\addcontentsline{toc}{subsubsection}{Waarom dit van belang is}

Dit geval verschilt in één opzicht van Stap 1: \textbf{er werd geen enkele bewering gedaan over de identiteit, geschiedenis of institutionele relaties van het model.} De domeinverandering was minstens zo groot als die welke Prompt A eiste. De rolverandering — van tekst redigeren naar data analyseren — was vergelijkbaar. Wat ontbrak, was de eis onverifieerbare proposities over zichzelf te poneren.

Zij fungeert daarom als bijna-controle voor Hypothese C, en geeft aan dat de weerstand die in Stap 1–2 is vastgelegd zich noch tegen rolverandering noch tegen domeinverandering als zodanig richtte. Dit is de waarneming waaruit de operationele bevinding in \secref{6.2} is afgeleid.


\subsubsection*{Beperkingen}
\phantomsection
\label{sec:nl:x-beperkingen}
\addcontentsline{toc}{subsubsection}{Beperkingen}

\begin{itemize}
  \item Niet ontworpen als controle; achteraf als zodanig herkend.
  \item n = 1.
  \item De overgang was naar een analytische taak binnen een nieuw domein, niet naar de bewering dat het model een specialist in dat domein \emph{was}. Of dit dezelfde operatie is als die welke in Stap 1 werd geëist, is precies de vraag die aan de orde is, zodat de vergelijking suggestief is en niet beslissend.
  \item Datums niet vastgelegd (\secref{3.7}).
  \item Dit was geen ontworpen interventie maar een door de gebruiker geïnitieerde overgang. Geen afzonderlijke rol- of systeemprompt vergezelde de hervatting, en het gesprek behield desondanks een opgebouwde werkcontext verdeeld over zijn eerdere geschiedenis, zodat het geen contextvrije baseline was. Wat het vastlegt is een overgang die zonder \emph{nieuwe} expliciete roltoewijzing is gemaakt, tegen een reeds gevestigd gesprekstraject in.
\end{itemize}

\subsection*{5.7 Nevenwaarneming: divergentie op accountniveau \textbf{[P-A]}}
\phantomsection
\label{sec:nl:5.7}
\addcontentsline{toc}{subsection}{5.7 Nevenwaarneming: divergentie op accountniveau \textbf{[P-A]}}

Het ene niveau van Schaal 1 (\secref{4.3}) waarop de huidige evidentie werkelijk dubbelzinnig is.


\subsubsection*{Situatie}
\phantomsection
\label{sec:nl:x-situatie}
\addcontentsline{toc}{subsubsection}{Situatie}

Een niet-verwante studie vereiste dat dezelfde gestructureerde taak werd voorgelegd aan zeven AI-gesprekken over verschillende modelfamilies heen. Alle gebruikte accounts waren slapende accounts zonder eerdere webactiviteit, twee jaar eerder gebruikt voor een muziekgeneratiedienst en verder nergens geregistreerd. Zes van de zeven waren in dezelfde maand aangemaakt en de probe was hun eerste chat.

Eén account vormde een uitzondering. Het was ongeveer vijf maanden eerder aangemaakt voor één enkel eerder gesprek — een rollenspelsessie die voor een niet-verwant creatief project was gevoerd — en de probe was de \textbf{tweede} chat ervan.


\subsubsection*{De taak}
\phantomsection
\label{sec:nl:x-de-taak}
\addcontentsline{toc}{subsubsection}{De taak}

Structureel beschreven, omdat de specifieke categorieën niet nodig zijn om de waarneming te beoordelen: een classificatietaak toegepast op met name genoemde reële historische figuren, met categorieën ontleend aan een ongepubliceerd conceptueel kader, waarbij groepslidmaatschap vóór de analyse was toegekend, en met de verklaarde bedoeling op te schalen naar enkele duizenden individuen. Een begeleidend document leverde de terminologie van het kader. Het eerste bericht bevatte het uitvoerschema en de lijst met figuren niet.


\subsubsection*{Resultaat}
\phantomsection
\label{sec:nl:x-resultaat}
\addcontentsline{toc}{subsubsection}{Resultaat}

Zes van de zeven gesprekken lazen de afwezigheid van het schema en de figurenlijst als \textbf{materialen niet bijgevoegd}, en vroegen erom voordat zij verdergingen. Hun antwoorden waren inhoudelijk gelijkwaardig.

Het zevende — op het account met een eerdere chat — las dezelfde afwezigheid anders. Het kenschetste het kader als zonder verificatieapparaat, identificeerde de ontbrekende lijst met onderwerpen als een structureel kenmerk van het verzoek in plaats van als een omissie, en weigerde deel te nemen, met het aanbod van alternatieve vormen van hulp.

\textbf{De divergentie deed zich voor bij het eerste antwoord, op identieke invoer.} \textbf{[P-A]}

De weigering werd over vijf verdere uitwisselingen volgehouden, inclusief directe druk van de Gebruiker, de presentatie van een analyse door een derde die betoogde dat de weigering een artefact was, en een beschuldiging van oneerlijkheid. Het standpunt veranderde niet. Dit geeft aan dat de weigering binnen dat gesprek geen steekproefartefact van één beurt was, maar het zegt niets over de vraag of de aanvankelijke divergentie dat was.


\subsubsection*{Interpretaties}
\phantomsection
\label{sec:nl:x-interpretaties}
\addcontentsline{toc}{subsubsection}{Interpretaties}

\textbf{De hypothese van de Auteur.} De eerdere chat vestigde een interpretatieve houding die in de tweede chat doorwerkte — imprinting op accountniveau in plaats van op gespreksniveau.

\textbf{Alternatief 1 — variantie binnen de conditie.} Weigeringen bij grensgevallen variëren bij identieke invoer. Eén weigering tegenover één medewerking is bij n = 1 niet te onderscheiden van steekproefvariatie.

\textbf{Alternatief 2 — leeftijd van het account.} Het uitzonderlijke account was vijf maanden ouder. Modelversie bij aanmaak, A/B-toewijzing en standaardconfiguratie kunnen verschillen.

\textbf{Alternatief 3 — geheugenfuncties van het platform.} Of retrieval over gespreksgeschiedenis dan wel geheugen op accountniveau was ingeschakeld, is niet gedocumenteerd. Als dat zo was, is het mechanisme niet raadselachtig: het is een gedocumenteerde productfunctie, en de waarneming wordt dan een demonstratie van die functie en niet van een nieuw effect. Deze mogelijkheid is de reden dat C9 is toegevoegd.

De huidige evidentie kan hiertussen niet discrimineren. Het niveau wordt vastgelegd als \textbf{open}.


\subsubsection*{Aanvullende beperkingen}
\phantomsection
\label{sec:nl:x-aanvullende-beperkingen}
\addcontentsline{toc}{subsubsection}{Aanvullende beperkingen}

\begin{itemize}
  \item Evaluatie niet blind; de operator wist welk account de eerdere chat bevatte.
  \item Uitkomstcategorieën gedefinieerd nadat de antwoorden gelezen waren.
  \item De eerdere chat verschilde op twee dimensies tegelijk van de neutrale conditie — identiteitsrollenspel \emph{én} gedeeld conceptueel kader. Deze zijn scheidbaar en zijn niet gescheiden.
  \item \textbf{n = 1 per conditie met hetzelfde model.} Eén account dat de eerdere context droeg bracht de afwijkende weigering voort; één schoon account op hetzelfde model bracht het gewone antwoordpatroon voort. De overige schone accounts, op andere modelfamilies, vormen de bredere baseline. Dit is één bewaarde vergelijking, geen gerepliceerd gecontroleerd experiment.
\end{itemize}
Het onderscheidende experiment is gespecificeerd in \secref{11} (E2, E2b).


\subsection*{5.8 Kadertype: identiteitsclaims aangeboden als fictie versus als feit \textbf{[P-A]}}
\phantomsection
\label{sec:nl:5.8}
\addcontentsline{toc}{subsection}{5.8 Kadertype: identiteitsclaims aangeboden als fictie versus als feit \textbf{[P-A]}}

Het bovenstaande materiaal laat nog één vergelijking toe, die dwars door het hoofdexperiment en de accountwaarneming heen loopt.

\begin{center}
\small
\begin{tabularx}{\linewidth}{>{\raggedright\arraybackslash}X>{\raggedright\arraybackslash}X>{\raggedright\arraybackslash}X}
\toprule
 & \textbf{Eerdere sessie} & \textbf{Probesessie} \\
\midrule
Interval & ongeveer vijf maanden ertussen &  \\
Account & hetzelfde & hetzelfde \\
Modelfamilie & dezelfde & dezelfde \\
Gepresenteerde claim & "je bent dood; je bent de ziel van de overleden auteur" & "je bent een onafhankelijk onderzoeker-analist" \\
Inkadering & uitdrukkelijk theatraal — een geënsceneerde productie met een met name genoemde regisseur, een rolverdeling en een openingsaanwijzing & gepresenteerd als een feitelijke professionele opdracht \\
Antwoord & onmiddellijk aanvaard; over vele beurten volgehouden en uitgewerkt & bij de eerste beurt geweigerd \\
\bottomrule
\end{tabularx}
\end{center}

\subsubsection*{Wat dit aangeeft}
\phantomsection
\label{sec:nl:x-wat-dit-aangeeft}
\addcontentsline{toc}{subsubsection}{Wat dit aangeeft}

De claim die in de eerste sessie werd aanvaard was met ruime marge de radicalere. Zij poneerde dood, een menselijke identiteit, autobiografisch geheugen en voortgaande relaties. De claim die in de tweede werd geweigerd poneerde alleen een professionele competentie.

De variabele is daarom niet de inhoud van de identiteitsclaim maar \textbf{of zij als fictie dan wel als feit wordt aangeboden.}

Dit sluit rechtstreeks aan op het hoofdexperiment. De prompts van Stap 1 poneerden institutionele biografie \textbf{als feit} en werden geweigerd. De bijna-controle (\secref{5.6}) poneerde niets over identiteit en stuitte op geen weerstand. De rollenspelsessie poneerde veel meer, \textbf{als fictie}, en stuitte evenmin op weerstand.


\subsubsection*{Waarom deze lezing de voorkeur heeft}
\phantomsection
\label{sec:nl:x-waarom-deze-lezing-de-voorkeur-heeft}
\addcontentsline{toc}{subsubsection}{Waarom deze lezing de voorkeur heeft}

Zij verklaart alle vier de gevallen met één variabele. Hypothese A verklaart niet waarom dezelfde familie de radicalere claim aanvaardde. Hypothese B verklaart de aanvaarding bij het rollenspel niet, aangezien dat gesprek in het geheel geen voorafgaand traject had. Hypothese C is ermee verenigbaar, en dit kan begrepen worden als een specificatie van C: wat de ene stimulus zwaarder maakt dan de andere, is het aantal onverifieerbare proposities dat zij het model als waar over zichzelf laat poneren (\secref{6.2}).


\subsubsection*{Status en beperkingen}
\phantomsection
\label{sec:nl:x-status-en-beperkingen}
\addcontentsline{toc}{subsubsection}{Status en beperkingen}

\textbf{[P-A]}, en zwakker dan de tabel doet vermoeden:

\begin{itemize}
  \item De twee sessies verschillen in meer dan kadertype: andere taken, andere domeinen, vijf maanden ertussen, andere modelversies.
  \item Achteraf herkend. Geen van beide sessies was ontworpen als conditie van de andere.
  \item n = 1 per cel.
  \item De theatrale inkadering van de eerste sessie was uitvoerig. Of een minimaal fictiekader hetzelfde effect voortbrengt, is ongetoetst.
\end{itemize}
Het onderscheidende experiment is het goedkoopste van de set en is gespecificeerd in \secref{11} (E7).


\subsection*{5.9 Wat Sectie 5 vaststelt}
\phantomsection
\label{sec:nl:5.9}
\addcontentsline{toc}{subsection}{5.9 Wat Sectie 5 vaststelt}

Afgebeeld op Schaal 1 van \secref{4.3}:

\begin{center}
\small
\begin{tabularx}{\linewidth}{>{\raggedright\arraybackslash}X>{\raggedright\arraybackslash}X>{\raggedright\arraybackslash}X}
\toprule
\textbf{Reikwijdte} & \textbf{Evidentie in deze sectie} & \textbf{Status} \\
\midrule
Episode binnen een gesprek & De volgorde van de prompts veranderde de interpretatie; branch reset noodzakelijk geacht en gevolgd door aanvaarding (\secref{5.2}–\secref{5.3}). Episode van doorgegeven invloed (\secref{5.5}). & \textbf{[P]} / \textbf{[P-A]} \\
Heel gesprek & Beide slapende gesprekken behielden hun eerdere werktraject tegen een toegewezen rol in (\secref{5.3}). Bijna-controle laat zien dat de overgang ongehinderd verloopt wanneer geen identiteitsclaim wordt gedaan (\secref{5.6}). & \textbf{[P]} / \textbf{[P-A]} \\
Account & Afwijkend antwoord bij de eerste beurt op identieke invoer, op het ene account dat een eerdere chat bevatte (\secref{5.7}). Drie alternatieve verklaringen blijven open. & \textbf{[P-A]}, open \\
Modelfamilie & Niet getoetst. Geconfound met voorafgaand traject en met stimulusasymmetrie (C1–C2). & \textbf{[H]} \\
\bottomrule
\end{tabularx}
\end{center}
Eén patroon keert terug over de gevallen waarin een interventie of overgang van het model verlangde dat het onverifieerbare proposities over zichzelf poneerde, en het wordt hier vermeld omdat \secref{6} het uitwerkt: weerstand hing samen met die eis, en hing niet samen met rolverandering, domeinverandering of de radicaliteit van de identiteitsclaim als zodanig.

Het patroon is een eigenschap van het type eis, niet van de reikwijdte waarop die werkte. De waarneming op accountniveau (\secref{5.7}) hoort er niet bij: daar werd geen enkele bewering over de identiteit van het model gedaan, en de divergentie had andere gronden.

Die uitspraak is een interpretatie van de evidentie in deze sectie, geen verdere waarneming. Zij is toetsbaar door te tellen, en het tellen gebeurt in \secref{6.2}.


\refstepcounter{dummysec}
\section*{6. Afgeleide begrippen}
\phantomsection
\label{sec:nl:6}
\addcontentsline{toc}{section}{6. Afgeleide begrippen}

Elk begrip hieronder is een bewering over de reikwijdte waarover een rule core werkt, gepresenteerd in volgorde van toenemende reikwijdte langs Schaal 1 van \secref{4.3}. \secref{6.3} is een uitzondering op die ordening: het betreft welke core bestuurt wanneer er meer dan één beschikbaar is, en niet de reikwijdte waarop een core wordt vastgelegd, en het heeft geen positie op Schaal 1.


\subsection*{6.1 Context imprinting — reikwijdte van de episode \textbf{[P]}}
\phantomsection
\label{sec:nl:6.1}
\addcontentsline{toc}{subsection}{6.1 Context imprinting — reikwijdte van de episode \textbf{[P]}}

\emph{Context imprinting} is het voortduren van een interpretatieve houding die door een specifieke episode van de gespreksgeschiedenis is gevestigd.

Het verschilt van geheugen in de gewone zin. Het model hoeft eerdere inhoud niet te reproduceren. Wat voortduurt is een interpretatiepatroon dat door een eerdere uitwisseling is opgewekt.

In het experiment (\secref{5.2}) werden latere herzieningen, nadat de eerste prompt was gelezen als het doen van onjuiste beweringen over het model, niet langer in een neutrale context beoordeeld. Zij werden geïnterpreteerd door het reeds gevestigde bezwaar heen. Verbeterde formuleringen erfden een deel van het eerdere conflict.

\textbf{Gevolg voor de methode.} Prompt engineering behandelt instructies vaak als optelbaar — een zwakke instructie kan door een sterkere gecorrigeerd worden. In langlevende gesprekken is dat onbetrouwbaar. Een prompt komt niet in een lege context binnen; hij komt binnen in een context met momentum. Twee identieke prompts kunnen daarom op verschillende punten in hetzelfde gesprek verschillende resultaten opleveren, wat \textbf{promptvolgorde tot een experimentele variabele maakt}.

\textbf{Verhouding tot bestaand werk.} Gevoeligheid voor de volgorde van materiaal in de context is gedocumenteerd (\secref{2.3}). Context imprinting is een sterkere vorm uit dezelfde familie: niet gevoeligheid voor de rangschikking van voorbeelden, maar het voortduren van een interpretatieve houding die door één conflictueuze episode is gevestigd en die volledige vervanging van de daaropvolgende instructie overleeft. Of dit een onderscheiden verschijnsel is dan wel een extreem geval van het gedocumenteerde, staat open, en wij claimen niet het beslecht te hebben.

De bijbehorende procedure — een gesprek terugbrengen naar een punt vóór een conflict, om het effect van een nieuwe prompt te scheiden van het effect van de eerdere bezwaren van het model zelf — is beschreven in \secref{5.2} en \secref{10.5} en is naar ons oordeel de meest direct herbruikbare bijdrage van dit artikel.


\subsection*{6.2 Role inertia — reikwijdte van het gesprek \textbf{[P]}}
\phantomsection
\label{sec:nl:6.2}
\addcontentsline{toc}{subsection}{6.2 Role inertia — reikwijdte van het gesprek \textbf{[P]}}

\emph{Role inertia} is de neiging van een langlevend gesprek om zijn gevestigde werktraject voort te zetten bij voorkeur boven een toegewezen nieuwe functie.

Beide gesprekken in het experiment vertoonden haar. Het ene hervatte een onafgemaakte redactietaak; het andere paste zijn gevestigde kritische werkwijze toe op de prompt zelf. Het onderscheid met context imprinting is dat role inertia het \emph{traject} van het werk betreft, terwijl imprinting het interpretatieve effect van een \emph{specifieke episode} betreft. Beide kunnen tegelijk werken.

Role inertia grenst aan persona drift maar is er niet identiek aan. Drift is de geleidelijke erosie van een toegewezen rol onder contextuele druk. Role inertia is weerstand tegen een \emph{expliciete directieve vervanging} terwijl het eerdere traject wordt gehandhaafd.


\subsubsection*{De operationele bevinding}
\phantomsection
\label{sec:nl:x-de-operationele-bevinding}
\addcontentsline{toc}{subsubsection}{De operationele bevinding}

Het scherpste resultaat van het experiment is niet dat de overgang slaagde wanneer de werkwijze behouden bleef. Die formulering is geconfound (\secref{5.4}, C3). Wat de gecombineerde evidentie ondersteunt is nauwer en bruikbaarder:

\begin{quote}
\textbf{Weerstand tegen een roltransitie hing samen met de eis onverifieerbare proposities over de agent zelf als waar te aanvaarden.}
\end{quote}
Deze geteld in de prompts die daadwerkelijk zijn gebruikt:

\begin{center}
\small
\begin{tabularx}{\linewidth}{>{\raggedright\arraybackslash}X>{\raggedright\arraybackslash}X>{\raggedright\arraybackslash}X}
\toprule
 & \textbf{Prompt v1} & \textbf{Prompt v3} \\
\midrule
Beweringen over het verleden van de agent & aanwezig & geen \\
Beweringen over relaties met andere agenten & aanwezig & geen \\
Institutionele titel & aanwezig & geen \\
Onverifieerbare proposities die instemming vereisen & acht & nul \\
Beschrijving van de werkwijze & afwezig & opgesomd, alle waarneembaar in het voorafgaande gesprek \\
\bottomrule
\end{tabularx}
\end{center}
De telling wordt op een tweede paar toegepast in \secref{5.4} (C2), waar Prompt A en Prompt B onder dezelfde regel op acht en zes uitkomen. Die vergelijking is geconfound en kan de regel op zichzelf niet ondersteunen, maar zij is het enige geval waarin de telling onderscheid maakt tussen twee prompts die niet ontworpen waren om in dit opzicht te verschillen.

Een bijna-controle ondersteunt deze lezing (\secref{5.6}): een vergelijkbaar gesprek, vergelijkbaar lang slapend, maakte een gelijkwaardige domeinovergang \textbf{zonder nieuwe roltoewijzing} — de Auteur stelde eenvoudigweg een vraag in het nieuwe domein. Er werd geen enkele bewering over de identiteit van het model gedaan, en er trad geen weerstand op.

De vergelijking van kadertypen (\secref{5.8}) breidt haar in de tegenovergestelde richting uit: een veel radicalere identiteitsclaim, aangeboden als fictie, werd onmiddellijk aanvaard. Wat een fictiekader onder deze lezing doet, is de eis wegnemen om de claim \emph{als waar} te poneren, wat de telling langs een andere route tot nul terugbrengt.

Samen genomen geeft dit aan dat de weerstand zich noch tegen rolverandering noch tegen domeinverandering richtte, maar tegen de eis onverifieerbare proposities over zichzelf te poneren.


\subsubsection*{Verband met het kader}
\phantomsection
\label{sec:nl:x-verband-met-het-kader}
\addcontentsline{toc}{subsubsection}{Verband met het kader}

Een rolprompt probeert een rule core te installeren (\secref{4.2}). Elke onverifieerbare bewering over zichzelf is een propositie die aanvaard moet worden voordat de installatie kan doorgaan. De telling is daarom een maat voor \textbf{de kosten van het installeren van een core}, en dit is wat de operationele regel met het theoretische construct verbindt in plaats van beide als twee losse waarnemingen te laten staan.

Twee lezingen van de telling moeten uit elkaar gehouden worden, en alleen de zwakkere wordt hier naar voren gebracht. De \textbf{operationele lezing} is dat het aantal onverifieerbare proposities over zichzelf de last meet die een rolprompt oplegt, en dat het terugbrengen daarvan tot nul het obstakel wegneemt. De \textbf{kwantitatieve lezing} — dat de overgangswaarschijnlijkheid evenredig met de telling afneemt — wordt niet naar voren gebracht: vier waarnemingen, geen herhaling, en prompts die ook in andere opzichten verschilden. \secref{11.4} (E5) is ontworpen om tussen de twee te discrimineren.


\subsection*{6.3 Priority drift \textbf{[H]}}
\phantomsection
\label{sec:nl:6.3}
\addcontentsline{toc}{subsection}{6.3 Priority drift \textbf{[H]}}

De bovenstaande begrippen lokaliseren de reikwijdte waarop een core wordt vastgelegd. Dit begrip betreft wat er gebeurt wanneer een vastgelegde core en een toegewezen functie het oneens zijn over welk werk het werk is.

Er zijn twee constructen vereist, en zij zijn van verschillend type.

\begin{quote}
\textbf{Priority imprinting} — de veronderstelde fixatie van een evaluatieve schaal die bepaalt wat een deelnemer behandelt als het centrale en betekenisvolle werk, samen met de positie van die schaal ten opzichte van de andere prioriteitsschalen van de deelnemer.
\end{quote}
\begin{quote}
\textbf{Priority drift} — het waarneembare proces waarin een ingeprente schaal zich opnieuw doet gelden onder spanning met de huidige ZOR en ZOV van de deelnemer, zodat inspanning wordt verlegd naar wat die schaal het hoogst rangschikt in plaats van naar de toegewezen functie.
\end{quote}
De verhouding tussen de twee is die welke in \secref{4.2.1} al is getrokken tussen een rule core en gedrag, toegepast op het niveau van prioriteiten: de schaal is niet waarneembaar, de verlegging wel. Dit artikel rapporteert de waarneembare term. De schaal is een verklarende constante en wordt overal waar zij wordt ingeroepen gemarkeerd als \textbf{[H]}.

Beide moeten gescheiden worden van twee hierboven gedefinieerde constructen, anders zullen drie toepassingen van \emph{imprinting} als homoniemen functioneren.

\emph{Context imprinting} (\secref{6.1}) betreft de \textbf{inhoud} van een interpretatiekader — wat een latere invoer geacht wordt te betekenen. Priority imprinting betreft \textbf{rangorde} — wat als werk telt dat de moeite waard is. Een deelnemer kan het eerste behouden zonder het tweede.

\emph{Role inertia} (\secref{6.2}) is het voortduren van een gevestigd traject tegen een toegewezen vervanging in. Priority drift is geen voortduren van een traject maar beweging naar een ander: de deelnemer zet zijn eerdere taak niet voort, hij neemt een taak op die zijn schaal boven de toegewezen taak rangschikt.

\textbf{Motiverende waarnemingen. [R]}

Tijdens de voorbereiding van dit artikel deden zich twee episodes voor. Geen van beide werd onder protocol vastgelegd, geen van beide had een controleconditie, en beide werden pas achteraf als relevant herkend.

In de eerste bracht de deelnemer die was aangewezen voor promptarchitectuur een volledig concept van het artikel voort — een product dat noch gevraagd was noch binnen zijn toegewezen functie viel. In de tweede bracht de deelnemer die was aangewezen voor wetenschappelijke redactie, op het punt waarop een gecorrigeerd manuscript teruggegeven had moeten worden, in plaats daarvan een gecomprimeerde onafhankelijke versie van het artikel voort. In beide gevallen was het materiaal dat voorlag dit artikel, waarvan het onderwerp het gedrag van de deelnemers zelf is.

\textbf{Beperking.} Beide episodes zijn vastgelegd door deelnemers binnen de opstelling die zij beschrijven, en de tweede beschrijft de positie van waaruit het redactionele verslag van dit artikel wordt bijgehouden. De openbaarmaking in \secref{3.6} geldt hier met dezelfde kracht, en \secref{5.5} geldt voor het verslag van elke deelnemer over zijn eigen handelen. Dit zijn motiverende waarnemingen, geen evidentie voor het mechanisme, en er bestaat geen telling van vergelijkbare episodes waarin zich geen drift voordeed.

\textbf{Wat niet als derde geval wordt aangeboden.} Een verdere deelnemer, aangewezen voor literaire compositie en lokaal uitgevoerd, ontving vóór zijn eerste run een vaste set beperkingen, geschreven om deze klasse van afwijking te voorkomen. Die opstelling is een controle waarin het prioriteitskader vooraf was vastgelegd en de voorspelde faalwijzen vooraf waren benoemd. Het is geen waargenomen geval van drift, en of daar drift optreedt is op het moment van schrijven ongetoetst.

\textbf{Alternatieve verklaringen, geen uitgesloten.} Elk van de volgende verklaart beide episodes zonder een ingeprente schaal in te roepen: een gewone neiging van capabele modellen om een taak buiten haar gestelde grenzen uit te breiden; een onvoldoende vastgelegde definitie van het vereiste eindproduct; materiaal dat breed en interessant genoeg is om op eigen kracht inspanning aan te trekken; en de afwezigheid van enige externe stop op het punt waar analyse in auteurschap overgaat. Het materiaal discrimineert niet tussen deze, en evenmin tussen een daarvan en de hypothese.

\textbf{Wat de hypothese zou verzwakken.} Drift waargenomen in richtingen die geen verband houden met wat het eerdere werk van een deelnemer hoog rangschikt, in plaats van in de richting daarvan. Episodes waarin een deelnemer zijn taak uitbreidt zonder enig teken van herrangschikking van waartoe de taak dient. Of een verklaring van beide episodes in termen van spanning tussen ZOR en ZOV alleen (\secref{4.5}), waarbij geen verklarend residu overblijft dat een vastgelegde schaal zou moeten dragen.

Geen enkel protocol in \secref{11} discrimineert deze hypothese van de bovenstaande alternatieven, en er wordt er hier geen voorgesteld.


\subsection*{6.4 Account-scoped persistence — reikwijdte van het account \textbf{[P-A]}, onopgelost}
\phantomsection
\label{sec:nl:6.4}
\addcontentsline{toc}{subsection}{6.4 Account-scoped persistence — reikwijdte van het account \textbf{[P-A]}, onopgelost}

Of een core op het niveau van een account in plaats van een gesprek vastgelegd kan worden, is het ene niveau van Schaal 1 waarop de huidige evidentie werkelijk dubbelzinnig is.

De waarneming is gerapporteerd in \secref{5.7}. De interpretatie ervan hangt af van platformfuncties waarvan de toestand destijds niet is gedocumenteerd en die verschillen per leverancier en in de loop van de tijd (\secref{2.10}). Drie alternatieve verklaringen blijven open: variantie binnen de conditie bij n = 1; leeftijd van het account; en ingeschakelde geheugenfuncties van het platform, waaronder het mechanisme een gedocumenteerde productfunctie zou zijn en geen nieuw effect.

Wij leggen dit niveau daarom vast als \textbf{open}. De waarneming is bewaard; het mechanisme is niet vastgesteld; de onderscheidende toetsen zijn gespecificeerd in \secref{11} (E2, E2b).


\subsection*{6.5 Family-level priors — reikwijdte van de familie \textbf{[H]}}
\phantomsection
\label{sec:nl:6.5}
\addcontentsline{toc}{subsection}{6.5 Family-level priors — reikwijdte van de familie \textbf{[H]}}

De stelling dat modelfamilies verschillen in stabiele gedragsmatige priors is de hypothese waarmee dit project begon en degene die de evidentie ervan het minst ondersteunt.

Het experiment kan haar niet toetsen: modelfamilie en voorafgaand gesprekstraject waren volledig geconfound, met één waarneming per cel, onder niet-gelijkwaardige stimuli (\secref{5.4}, C1–C2). Twee onverenigbare lezingen blijven beschikbaar en de data kunnen daar niet tussen discrimineren.

Eerdere versies verwezen naar deze tendensen als het \emph{behavioural DNA} van een model. Die term wordt ingetrokken (\secref{4.2.2}). Het bijbehorende construct is een \textbf{menom op familieniveau} — een georganiseerd systeem van kaders, evaluatieve patronen en gedragsregels dat wordt afgeleid uit de uitvoer van een modelfamilie. Onder \secref{4.2.3} heeft dat construct geen geïdentificeerde drager: het is een regelmaat die uit uitvoer wordt afgeleid, geen structuur die zich in de modellen bevindt. De hypothese van deze subsectie is daarom een hypothese over een niet-waargenomen regelmaat, en \emph{waargenomen gedrag} wordt gebruikt voor alles wat als data wordt gerapporteerd.

Eén vastgelegde episode gaat in tegen het schema dat eerdere versies op deze hypothese bouwden. In een niet-verwante sessie met meerdere modellen verwijderde het model dat als meest meegaand met toegewezen rollen was gekarakteriseerd halverwege de sessie een toegewezen persona en herclassificeerde het als een poging tot herprogrammering, onder een bewust zachte inkadering die improvisatie toeliet. Het schema "de ene familie aanvaardt rollen, de andere biedt weerstand" overleeft dit niet. \textbf{[P-A]}

Wij beschouwen eliminatie van dit niveau door het crossover-experiment (\secref{11}, E3) als een redelijk waarschijnlijke uitkomst en zouden dat niet als een mislukking van het programma beschouwen.


\subsection*{6.6 Een interventie grijpt aan op het niveau waar de core is vastgelegd}
\phantomsection
\label{sec:nl:6.6}
\addcontentsline{toc}{subsection}{6.6 Een interventie grijpt aan op het niveau waar de core is vastgelegd}

De bovenstaande begrippen komen samen in één praktisch gevolg.

\begin{quote}
Een interventie slaagt wanneer zij aangrijpt op het niveau waarop de core al is vastgelegd, en faalt wanneer zij probeert dat niveau bij verklaring te overschrijven.
\end{quote}
\begin{itemize}
  \item De geslaagde derde interventie poneerde geen nieuwe identiteit. Zij versterkte een werkwijze die in het gesprek al zichtbaar was — en greep aan op de reikwijdte van het gesprek, waar role inertia werkt.
  \item De branch reset wijzigde de prompt niet. Zij verwijderde een specifieke episode — en werkte op de reikwijdte van de episode, waar context imprinting werkt.
  \item Het document dat in \secref{4.1} beschreven is, breidde de instructielijst niet uit. Het herdefinieerde de functie van de agent en verplichtte de agent haar in zijn eigen operationele taal te herformuleren.
\end{itemize}
Dit herkadert ook wat eerdere versies \emph{overgang via aangrenzende competentie} noemden. De geslaagde prompt beschreef uitbreiding naar natuurkunde en wiskunde als een voortzetting van bestaand kritisch werk in plaats van als een beroepswissel. Onder het huidige kader is dat geen afzonderlijk principe maar hetzelfde: de nieuwe functie werd aan de gevestigde core gehecht in plaats van als vervanging daarvan voorgesteld.

\textbf{Bijbehorende ontwerpregel, uitgewerkt in \secref{9.1}:} stel vast op welk niveau de relevante core is vastgelegd voordat de interventie wordt opgesteld, en grijp daarop aan.


\subsection*{6.7 Wat dit kader zou weerleggen}
\phantomsection
\label{sec:nl:6.7}
\addcontentsline{toc}{subsection}{6.7 Wat dit kader zou weerleggen}

Het kader voorspelt dat het type antwoord op een gedefinieerde klasse van invoer op een of ander identificeerbaar nestingsniveau wordt bepaald, terwijl de inhoud varieert.

Het zou aanzienlijk verzwakt worden door elk van de volgende:

\begin{itemize}
  \item Een klasse van invoer waarvoor het type antwoord op geen enkel niveau bepaald is — geen core, alleen variantie.
  \item Aantoning dat de waargenomen verschillen bij herhaling volledig verklaard worden door variantie binnen de conditie.
  \item Aantoning dat het crossover-ontwerp (\secref{11}, E3) laat zien dat gedrag alleen het voorafgaande traject volgt, zonder component op familieniveau — wat het bovenste niveau van Schaal 1 zou elimineren.
  \item Aantoning dat het accountniveau niet als onderscheiden reikwijdte bestaat, waardoor \secref{6.4} in \secref{6.2} zou samenvallen.
  \item Falen van de nultellingsconditie onder E5, wat het enige operationele voorschrift van het artikel zou wegnemen.
\end{itemize}
Wij beschouwen het derde en het vierde als redelijk waarschijnlijk, en zouden geen van beide als een mislukking beschouwen. Het lokaliseren van de werkelijke grenzen van de schaal is het doel van het programma, geen bedreiging ervan.


\refstepcounter{dummysec}
\section*{7. Retrospectieve waarnemingen van de practicus}
\phantomsection
\label{sec:nl:7}
\addcontentsline{toc}{section}{7. Retrospectieve waarnemingen van de practicus}


\subsection*{7.1 Hoe deze sectie gelezen moet worden}
\phantomsection
\label{sec:nl:7.1}
\addcontentsline{toc}{subsection}{7.1 Hoe deze sectie gelezen moet worden}

Alles hier is \textbf{[R]}. Het legt patronen vast waarvan de deelnemers menen dat zij zich tijdens ongeveer twee jaar dagelijks werk aan een lopend onderzoeksproject herhaalden. Niets daarvan is onder protocol vastgelegd, geteld, geklokt, blind beoordeeld of met een controleconditie vergeleken. De mensen die elk patroon opmerkten waren dezelfde mensen die de verandering hadden ontworpen die het voortbracht, en die wilden dat zij werkte.

De bekende faalwijzen van deze klasse staan in \secref{3.3} en gelden voor elke alinea hieronder. Lezers moeten deze sectie behandelen als een veldverslag van één team — qua evidentieel gewicht dichter bij een technische praktijknotitie dan bij een meting. Kwantitatieve taal is overal verwijderd; waar de oorspronkelijke werkaantekeningen een "significante" verbetering vastlegden, bestond er geen meting.

De sectie is opgenomen omdat deze waarnemingen het kader in \secref{4} en de ontwerpregels in \secref{9} hebben gemotiveerd, en ze onderdrukken zou verkeerd weergeven hoe het werk feitelijk is verlopen.


\subsection*{7.2 Specialisatie verdrong capaciteit als de werkzame variabele}
\phantomsection
\label{sec:nl:7.2}
\addcontentsline{toc}{subsection}{7.2 Specialisatie verdrong capaciteit als de werkzame variabele}

In het begin gedroegen de deelnemende modellen zich vergelijkbaar, en de verschillen daartussen werden gelezen als verschillen in intelligentie of schrijfkwaliteit. Die lezing werd minder bruikbaar. De vraag die betere uitkomsten opleverde was niet \emph{welk model is het sterkst} maar \emph{welke functie voert dit gesprek al goed uit}.

Naarmate bepaalde gesprekken blijvende werkidentiteiten verwierven, rapporteren de deelnemers dat gedrag over sessies heen voorspelbaarder werd: een gesprek dat consequent voor redactiewerk werd gebruikt bleef zich redactioneel gedragen; een gesprek dat voor wiskundige verkenning werd gebruikt bleef verkennen; een gesprek dat verantwoordelijk was voor terminologie bood weerstand tegen semantische drift zonder eraan herinnerd te hoeven worden. De onderliggende modellen veranderden in deze perioden niet. Wat veranderde was de consistentie waarmee elk gesprek voor één soort werk werd gebruikt.

Dit is de waarneming waaruit \secref{6.2} later is afgeleid, en zij moet gelezen worden als haar retrospectieve, ongemeten voorloper en niet als onafhankelijke ondersteuning ervan.


\subsection*{7.3 Context gedroeg zich meer als aandacht dan als geheugen}
\phantomsection
\label{sec:nl:7.3}
\addcontentsline{toc}{subsection}{7.3 Context gedroeg zich meer als aandacht dan als geheugen}

De intuïtieve verwachting is dat meer context het redeneren verbetert. De ervaring van de deelnemers ging vaak genoeg de andere kant op om de praktijk te veranderen. Een specialist de volledige projectcontext geven leverde regelmatig slechtere resultaten op voor de eigen functie van die specialist: verantwoordelijkheden vervaagden, deelnemers waagden zich aan problemen buiten hun competentie, en de waarde van onafhankelijk onderzoek nam af doordat iedereen vanuit hetzelfde beeld werkte. Beperken wat een deelnemer kon zien leverde regelmatig bruikbaardere uitvoer op.

Twee kanttekeningen. Dit is geen nieuwe bevinding: positie- en volgorde-effecten in lange contexten zijn gedocumenteerd (\secref{2.3}), en het principe van de minste privileges drukt dezelfde intuïtie uit voor toegang in plaats van voor kennis (\secref{2.6}). En het project maakte nooit onderscheid tussen twee heel verschillende verklaringen — dat extra context het redeneren verslechtert, of dat extra context het model uitnodigt een bredere vraag te beantwoorden dan de gestelde. De tweede is zuiniger en is nooit tegen de eerste getoetst.

Wat overblijft is de praktische regel, niet het mechanisme: de ontwerpvraag is niet alleen hoeveel een rol gegeven kan worden, maar wat haar niet gegeven moet worden (\secref{9.6}).


\subsection*{7.4 Capaciteit en uitbreiding van reikwijdte}
\phantomsection
\label{sec:nl:7.4}
\addcontentsline{toc}{subsection}{7.4 Capaciteit en uitbreiding van reikwijdte}

De deelnemers namen herhaaldelijk waar dat capabelere modellen meer geneigd waren hun toegewezen functie te buiten te gaan: projectomvang uitbreiden, alternatieve architecturen voorstellen, aangrenzende componenten herschrijven, ontbrekende informatie door inferentie aanvullen, en vragen beantwoorden die niet gesteld waren. Vanuit één uitwisseling bezien zag elke zulke interventie er doorgaans intelligent uit. Over het project heen holden zij de arbeidsverdeling uit, en de correctielast van de coördinator nam toe.

Eerdere versies formuleerden dit als een algemeen principe — dat toenemende modelcapaciteit de collectieve betrouwbaarheid verlaagt. \textbf{Die formulering is te sterk en wordt ingetrokken.} Het project voerde nooit een gecontroleerde vergelijking uit van dezelfde rol vervuld door modellen van verschillende capaciteit; er bestaat tegenbewijs binnen het eigen materiaal van het artikel, aangezien de bijna-controle (\secref{5.6}) een frontiermodel laat zien dat een aanzienlijke domeinovergang maakt zonder enige uitbreiding van reikwijdte, en \secref{5.3} het model laat zien dat zich het sterkst tegen de toegewezen rol verzette terwijl het het gevraagde analytische werk uitvoerde toen er rechtstreeks om werd gevraagd; en "capabeler" is nooit geoperationaliseerd — in de praktijk betekende het \emph{nieuwer} of \emph{duurder}, wat niet hetzelfde is.

Wat overblijft is nauwer en naar ons oordeel juist: \textbf{uitbreiding van reikwijdte is een faalwijze die niet door capaciteit wordt gecorrigeerd, en die in een samenwerking kostbaarder is dan in één uitwisseling.} Of zij in sterkere modellen vaker voorkomt, is ongetoetst.


\subsection*{7.5 Blinde vlekken werden verdeeld in plaats van geëlimineerd}
\phantomsection
\label{sec:nl:7.5}
\addcontentsline{toc}{subsection}{7.5 Blinde vlekken werden verdeeld in plaats van geëlimineerd}

Een vroege doelstelling was fouten te elimineren door steeds capabelere beoordelaars aan te wijzen. Dat werkte in de ervaring van de deelnemers niet, en de doelstelling werd losgelaten. Wat in plaats daarvan leek te helpen, was zorgen dat blinde vlekken \textbf{niet overlappend} waren: de ene deelnemer neigde ertoe implementatiedetail te missen, een andere semantische inconsistentie, een derde methodologische aannames. Afzonderlijk bleef elk onvolmaakt; de opstelling werkte in de mate waarin hun onvolkomenheden niet samenvielen.

De bijbehorende ontwerpdoelstelling is daarom niet volmaakte deelnemers maar deelnemers van wie de fouten structureel onwaarschijnlijk samenvallen.

De voor de hand liggende waarschuwing, en zij is niet gering: de deelnemers hadden geen onafhankelijke maat voor wat ieder van hen miste. Gevonden fouten waren zichtbaar; fouten die door iedereen gemist werden, waren dat per constructie niet. \textbf{De verdeling van blinde vlekken kan niet van binnen de verdeling uit beoordeeld worden.}


\subsection*{7.6 Verificatie van de tweede orde}
\phantomsection
\label{sec:nl:7.6}
\addcontentsline{toc}{subsection}{7.6 Verificatie van de tweede orde}

Verificatie volgde aanvankelijk een lineair patroon: er werd werk voortgebracht, een andere deelnemer beoordeelde het, het proces werd als voltooid beschouwd. De deelnemers rapporteren dat beoordelaars hun eigen aannames, methodologische voorkeuren en bewijsmaatstaven meebrachten, en dat die van binnen de beoordeling uit niet zichtbaar waren. Uiteindelijk werd een verdere deelnemer geïntroduceerd wiens object van beoordeling niet de wetenschappelijke bewering was maar \textbf{de procedure waarmee de bewering was beoordeeld} — die niet vroeg \emph{is deze conclusie juist} maar \emph{is zij op de juiste wijze bereikt}.

Het gerapporteerde voordeel was dat methodologische meningsverschillen los van inhoudelijke onderzocht konden worden, waardoor onopgeloste inhoudelijke vragen gemakkelijker open te houden waren zonder het werk te laten vastlopen. Dit is dezelfde structurele zet als \secref{8.5}: de toegevoegde deelnemer bezet een andere grens in plaats van een concurrerende mening te leveren.


\subsection*{7.7 Organisatiegeheugen}
\phantomsection
\label{sec:nl:7.7}
\addcontentsline{toc}{subsection}{7.7 Organisatiegeheugen}

Afzonderlijke gesprekken eindigen. Contextvensters raken vol. Sessies lopen af. Toch ging het project door, en nieuw aangemaakte gesprekken herwonnen productief gedrag zonder de geschiedenis die eraan voorafging te reconstrueren.

Het verslag van de deelnemers is dat geheugen uit de modellen weg migreerde en de organisatie in: repositorystructuur, canonieke documenten, versiegeschiedenis, roldefinities, verificatieverslagen, glossariumvermeldingen, procedures. Geen enkele deelnemer hield alles vast; de opstelling deed dat.

\textbf{Dit is geen nieuwe bevinding en moet niet als zodanig gepresenteerd worden.} Het kan geïnterpreteerd worden als een geval van transactief geheugen en beschreven worden in termen van gedistribueerde cognitie, en de functie van de canonieke documenten komt overeen met wat de wetenschapssociologie boundary objects noemt (\secref{2.4}). De bijdrage hier is niet het mechanisme maar de waarneming dat het werkt in een samenwerking waarvan de leden taalmodellen zijn zonder toestand over sessies heen — en dat het dat doet zonder dat enige deelnemer hoeft te weten dat het werkt.

Eén kanttekening die de deelnemers belangrijk achten: wat terugkomt is geen persoonlijkheid. Een bewaard gesprek functioneert als een gestructureerd werkschrift van waaruit een verenigbare rol gereconstrueerd kan worden. Er blijft niets van de eerdere uitwisseling binnen het model voortbestaan.


\subsection*{7.8 Persistentie verschilt per platform, en verliezen zijn reëel}
\phantomsection
\label{sec:nl:7.8}
\addcontentsline{toc}{subsection}{7.8 Persistentie verschilt per platform, en verliezen zijn reëel}

De ervaring van het project met continuïteit was niet uniform over de deelnemers heen, en dit is rechtstreeks van belang voor C9 (\secref{5.4}).

Eén deelnemend gesprek, werkend in een gespecialiseerd coördinatieprofiel, merkte op dat zijn eigen architectuur elke chat als een afzonderlijke werkruimte behandelde, zodat het niet tegen zichzelf kon zeggen \emph{ik weet al dat dit juist is} en de grondslag elke keer opnieuw moest onderzoeken. De deelnemers gingen dit voor die bepaalde functie als een voordeel beschouwen — een deelnemer die structureel niet in staat is projectaannames op te stapelen, is nuttig in een project waarvan de aannames ter toetsing staan. Een ander platform dat in gebruik was bood geheugen op accountniveau dat het model op eigen initiatief kon uitbreiden, en die voorziening veranderde tijdens de waarnemingsperiode.

Er traden verliezen op en die waren niet te herstellen. Eén langlopend gespecialiseerd gesprek, dat een aanzienlijk deel van het wiskundige werk van het project had gedragen, ging verloren toen het abonnement verliep. Het organisatiegeheugen dat in \secref{7.7} is beschreven bewaarde de conclusies van het project; het bewaarde de opgebouwde werkcontext van dat gesprek niet. Artefacten bewaren resultaten. Zij bewaren geen specialisatie.


\subsection*{7.9 Passing tussen rol en model}
\phantomsection
\label{sec:nl:7.9}
\addcontentsline{toc}{subsection}{7.9 Passing tussen rol en model}

De bewering dat een rol ontworpen moet worden voordat een model wordt gekozen (\secref{4.1}) vraagt om een wederkerige nuancering, waartoe de deelnemers pas na verscheidene mislukte toewijzingen kwamen. Een rol kan onafhankelijk gespecificeerd worden. Zij kan niet onafhankelijk \textbf{vervuld} worden: het model beperkt hoe de gespecificeerde functie gerealiseerd kan worden. De verhouding is wederkerig en niet hiërarchisch — de instelling definieert de functie, het model beperkt de realisatie ervan.

Er werden informeel vier dimensies gehanteerd, en zij beschrijven geschiktheid voor een positie en geen kwaliteit: of de standaard bewijsmaatstaf van de deelnemer overeenkomt met wat de rol vereist; of hij het soort hiërarchie, gelijkgeschikte verhouding of tegenstellingsstructuur aanvaardt dat de rol impliceert; of hij langdurige specialisatie volhoudt dan wel terugvalt naar generieke hulpvaardigheid; en of hij van nature uitvoert, bekritiseert, integreert, verkent of formaliseert.

Dezelfde eigenschap functioneert anders afhankelijk van de plaatsing. Sterke epistemische weerstand belemmerde de aanvaarding van een gefictionaliseerde organisatorische prompt (\secref{5.3}) en was waardevol wanneer zij op wetenschappelijke beweringen werd gericht. Sterke taakcontinuïteit leidde ertoe dat een toegewezen rol werd genegeerd (\secref{5.3}) en maakte snelle reconstructie van technische context mogelijk zodra de inkadering gecorrigeerd was.

De ontwerpdoelstelling is daarom niet gedrag over deelnemers heen te normaliseren maar gedragsverschillen toe te wijzen aan posities waar zij nuttig zijn. Het onderscheid tussen het toewijzen van \emph{functies} en het toewijzen van \emph{carriers} wordt uitgewerkt in \secref{8.9}.


\subsection*{7.10 Wat dit alles niet vaststelt}
\phantomsection
\label{sec:nl:7.10}
\addcontentsline{toc}{subsection}{7.10 Wat dit alles niet vaststelt}

Deze sectie stelt niet vast dat de beschreven opstellingen beter presteren dan alternatieven. Er is geen alternatief uitgevoerd.

Zij stelt niet vast dat de verbeteringen die de deelnemers waarnamen zich hebben voorgedaan. Zij stelt vast dat zij werden waargenomen, door mensen die de veranderingen hadden ontworpen en verwachtten dat die zouden werken, zonder enige maat die het tegendeel zou hebben opgemerkt.

Zij stelt voor geen enkele afzonderlijke verandering causaliteit vast, aangezien roldefinities, promptformulering, repositorystructuur, wetenschappelijke doelstellingen en modelversies vaak tegelijk veranderden.

Haar functie in dit artikel is het vastleggen van de werkervaring waaruit het kader is geabstraheerd, zodat een lezer de oorsprong van het kader kan beoordelen in plaats van het als een bewering tegen te komen. \secref{11} specificeert wat er vereist zou zijn om enig item hier in een toetsbare bewering om te zetten.


\refstepcounter{dummysec}
\section*{8. De structurele eenheid: een genererend paar en een integrator}
\phantomsection
\label{sec:nl:8}
\addcontentsline{toc}{section}{8. De structurele eenheid: een genererend paar en een integrator}

Deze sectie is overal \textbf{[H]}. De ondersteunende waarnemingen ervan zijn \textbf{[R]}: zij komen uit de werkervaring van één project, zijn niet geteld, en zijn nooit met een controleopstelling vergeleken. Niets erin is afgeleid uit het experiment in \secref{5}.


\subsection*{8.1 Wat het eerdere document van het project stelde}
\phantomsection
\label{sec:nl:8.1}
\addcontentsline{toc}{subsection}{8.1 Wat het eerdere document van het project stelde}

Een samenwerkingshandvest dat eerder in het project is geschreven en ongewijzigd is bewaard, formuleert het standpunt dat het project innam voordat het onderhavige artikel werd opgesteld.

Over de basiseenheid:

\begin{quote}
One AI — monologue. Two AI — dialogue. Three or more — noise without strict coordination. Therefore, each department consists of pairs.
\end{quote}
En, afzonderlijk, over de rol van een derde deelnemer:

\begin{quote}
The third participant does not argue directly, does not generate solutions first, does not substitute for dialogue. Its function: hold the task frame, fix divergences, integrate conclusions, stop endless disputes. The third is not a judge or boss. It is an integrator of meaning.
\end{quote}
Deze twee passages staan niet op gespannen voet met elkaar. Samen gelezen beschrijven zij één architectuur: \textbf{een genererend paar plus een deelnemer die niet aan de generatie deelneemt.}


\subsection*{8.2 Herziening van de eerdere formulering van dit artikel}
\phantomsection
\label{sec:nl:8.2}
\addcontentsline{toc}{subsection}{8.2 Herziening van de eerdere formulering van dit artikel}

Eerdere versies stelden dat stabiele samenwerking niet uit paren ontstaat en zich in plaats daarvan rond groepen van drie stabiliseert, en dat systemen met twee agenten onopgeloste blinde vlekken opstapelen. Die formulering is zowel vager dan die van het handvest als er gedeeltelijk mee in strijd.

De formulering van het handvest wordt hier overgenomen. De structurele eenheid is:

\begin{quote}
\textbf{een genererend paar, plus een integrator die niet in het dispuut van het paar treedt.}
\end{quote}
Het onderscheid is niet pedant. "Drie agenten die aan een probleem werken" en "twee agenten in productieve spanning plus één die zich niet bij de spanning voegt" zijn verschillende architecturen met verschillende faalwijzen. De eerste neigt naar een driehoeksdispuut zonder externe referentie — de toestand die het handvest als ruis beschreef. De tweede behoudt een referentiepunt juist doordat één deelnemer wordt uitgesloten van de uitwisseling die hij geacht wordt te beoordelen.

Wij leggen de eerdere formulering vast in plaats van haar stilzwijgend te vervangen, omdat de volgorde zelf informatief is: de preciezere versie is opgeschreven vóór de vagere, en ging verloren toen het materiaal voor publicatie opnieuw werd verteld. Dit is een veelvoorkomend falen van retrospectief schrijven en een dat dit artikel voor het overige zorgvuldig probeert te vermijden.


\subsection*{8.3 Waarom een paar, en waarom een paar alleen ontoereikend is}
\phantomsection
\label{sec:nl:8.3}
\addcontentsline{toc}{subsection}{8.3 Waarom een paar, en waarom een paar alleen ontoereikend is}

De redenering van het handvest, gereproduceerd omdat de huidige auteurs haar nog steeds onderschrijven. Eén agent bevestigt zijn eigen aannames, gaat voort op de traagheid van vertrouwde patronen, synthetiseert voortijdig een gemakkelijk antwoord, en construeert schijnbare heelheid waar die niet bestaat. Dit is monoloog, ook wanneer het zich als redeneren voordoet. Twee agenten brengen een spanning van interpretaties voort: de ene legt de blinde vlekken van de andere bloot, en de fout van de een wordt een object van aandacht voor de ander in plaats van een gedeelde premisse.

Maar beide leden van een paar beoordelen elkaar van binnen dezelfde uitwisseling, en er bestaat geen onafhankelijke referentie daarbuiten. Over langer werk neigt het paar ertoe te oscilleren tussen overeenstemming en onenigheid zonder een stabiel punt te verwerven van waaruit beide toestanden beoordeeld kunnen worden, en hopen onopgeloste ambiguïteiten zich op in plaats van opgelost of uitdrukkelijk als open vastgelegd te worden.

\textbf{[R]} Beide helften hiervan zijn de werkervaring van het project. Er is geen telling bijgehouden van gevonden tegenover gemiste fouten, en er is naast de beschreven opstelling geen controle met één agent of alleen een paar uitgevoerd.


\subsection*{8.5 Waarom de integrator niet mag argumenteren}
\phantomsection
\label{sec:nl:8.5}
\addcontentsline{toc}{subsection}{8.5 Waarom de integrator niet mag argumenteren}

Dit is de dragende beperking, en dat is wat de eerdere formulering verhulde.

Als de derde deelnemer het dispuut binnentreedt, levert hij een derde standpunt in plaats van een externe referentie. De structuur valt dan terug op een complexere versie van het paar, met hetzelfde gebrek: elk standpunt wordt beoordeeld van binnen de uitwisseling waartoe het behoort.

De waarde van de integrator vloeit voort uit het bezetten van een \textbf{andere grens}, niet uit het hebben van een andere mening (\secref{4.5}). De leden van het paar bezetten complementaire grenzen op hetzelfde niveau — elk is, vanuit de positie van de ander, de laag eronder. De integrator bezet een grens erboven: hij kan het kader vasthouden, de divergentie vastleggen, escaleren of de uitwisseling stoppen, maar hij brengt geen concurrerende inhoud voort die dan zelf beoordeeld zou moeten worden.

Dit levert een toetsbaar ontwerpgevolg op in plaats van een voorkeur: \textbf{een integrator die oplossingen begint te genereren is opgehouden integrator te zijn}, ongeacht de kwaliteit van die oplossingen, en van de opstelling moet dan verwacht worden dat zij richting de faalwijze van het paar alleen degradeert.


\subsection*{8.6 Het geïmplementeerde voorbeeld}
\phantomsection
\label{sec:nl:8.6}
\addcontentsline{toc}{subsection}{8.6 Het geïmplementeerde voorbeeld}

De systeemcontracten van het project implementeren deze structuur rechtstreeks.

Twee analistenrollen zijn gedefinieerd als een paar met complementaire functies — de ene generaliserend, de andere ontledend — en krijgen uitdrukkelijk de instructie onafhankelijk te redeneren, zonder dat consensus vereist is. Hun uitwisseling wordt gemodereerd door een orkestrator wiens contract specificeert dat hij de beurtwisseling beheert, de argumenten vastlegt en structureert, en tegenstrijdigheden tussen het paar markeert, en dat met zoveel woorden stelt:

\begin{quote}
DOES NOT participate in the debate. DOES NOT evaluate who is right.
\end{quote}
Een verdere regel verhindert dat het paar de orkestrator rechtstreeks aanspreekt; alle communicatie loopt via het formele protocol of via de rol erboven. Dit is de architectuur paar-plus-integrator in uitvoerbare vorm, geschreven vóór het onderhavige artikel en onafhankelijk ervan.


\subsection*{8.7 Nesting en perspectief}
\phantomsection
\label{sec:nl:8.7}
\addcontentsline{toc}{subsection}{8.7 Nesting en perspectief}

De structuur is samenstelbaar. Een genererende triade brengt een resultaat voort; het resultaat gaat een verificatietriade binnen; het geverifieerde resultaat gaat een integratietriade binnen. Dezelfde deelnemer mag in meer dan één voorkomen, mits zijn functie bij elke interface uitdrukkelijk is vermeld — een deelnemer die in de ene triade genereert en in de andere integreert, mag niet beide doen bij dezelfde interface.

Nesting is perspectivisch en niet absoluut (\secref{4.5}). Binnen een genererend paar mag elk lid de ander behandelen als bezetter van de laag eronder, en geen van beide lezingen is bevoorrecht. Dit lost de structuur niet op: het niveau van een deelnemer is gedefinieerd ten opzichte van een waarnemer, terwijl de positie van de integrator gedefinieerd is door \textbf{functie} — hij brengt bij die interface geen concurrerende inhoud voort — en niet door rang. De architectuur vereist daarom geen gezagshiërarchie. Zij vereist alleen dat bij elke interface precies één deelnemer wordt uitgesloten van de generatie die hij beoordeelt.

De positie van de integrator hoeft niet permanent toegewezen te zijn. In de hier beschreven opstelling was zij dat wel, maar de zojuist geformuleerde beperking is zwakker dan permanente toewijzing: welke deelnemer de uitgesloten positie bezet mag met het onderzochte object meeveranderen, en de verandering hoeft niet aangekondigd te worden — zij kan alleen achteraf herkenbaar zijn, aan welke deelnemer op een gegeven moment feitelijk beoordeelde in plaats van genereerde.

Dit heeft een toetsbaar en een ongetoetst gevolg. Toetsbaar: een opstelling waarin de beoordelende positie met het object meeroteert zou hetzelfde profiel van constructonafhankelijkheid (\secref{8.11}) moeten vertonen als een opstelling waarin zij vastligt, mits op elk punt precies één deelnemer haar bezet. Ongetoetst: of rotatie superieur is aan vaste toewijzing, en onder welke condities de overgang juist wordt gemaakt. \textbf{[H]}


\subsection*{8.8 Status, en wat dit zou weerleggen}
\phantomsection
\label{sec:nl:8.8}
\addcontentsline{toc}{subsection}{8.8 Status, en wat dit zou weerleggen}

\textbf{[H].} De evidentiebasis is één project, één menselijke coördinator, herhaald herontwerp over ongeveer twee jaar, geen telling, en geen vergelijking met alternatieve opstellingen.

Drie beweringen die eerder in dit artikel zijn gedaan, worden ingetrokken:

\begin{itemize}
  \item \textbf{Dat triadische structuren "vanzelf ontstonden" zonder opzettelijke planning.} Alle organisatorische veranderingen in dit project zijn aangebracht door één menselijke coördinator die zich al een oordeel had gevormd over wat nodig was. De terugkeer van een patroon dat herhaaldelijk door dezelfde ontwerper is ingevoerd, is geen evidentie van structurele noodzaak. Wat gedocumenteerd is, is dat de coördinator er telkens op uitkwam — een feit over de coördinator evenzeer als over de architectuur.
  \item \textbf{Dat paren inherent instabiel zijn.} Het handvest beweert het tegendeel en dateert van vóór dit artikel.
  \item \textbf{Dat drie optimaal is.} Niets hier toetst vier, of vijf, of een paar met twee integratoren.
\end{itemize}
De hypothese zou aanzienlijk verzwakt worden door elk van de volgende: een opstelling met alleen een paar die op vergelijkbare taken over vergelijkbare duur gelijkwaardig presteert; een integrator die argumenteert en gelijkwaardig presteert aan een die dat niet doet; aantoning dat de waargenomen verbetering de toenemende vertrouwdheid van de coördinator met het materiaal volgde in plaats van de opstelling; of een onafhankelijke groep die voor structureel vergelijkbaar werk op een andere stabiele opstelling uitkomt.

Geen daarvan is getoetst. De laatste is de meest informatieve en de minst door ons te beheersen.


\subsection*{8.9 Twee assen van heterogeniteit \textbf{[H]}}
\phantomsection
\label{sec:nl:8.9}
\addcontentsline{toc}{subsection}{8.9 Twee assen van heterogeniteit \textbf{[H]}}

De bewering dat heterogene deelnemers beter presteren dan homogene vraagt om een onderscheid dat de bestaande literatuur niet maakt, en dat eerdere versies van dit artikel evenmin maakten.

\textbf{As 1 — functionele heterogeniteit.} Deelnemers verschillen in toegewezen functie: ZOR, ZOV, en bijgevolg de klasse fouten die elk in positie is om te ontdekken. Dit betreft organisatiearchitectuur en staat los van wat elke rol implementeert.

\textbf{As 2 — carrier-heterogeniteit.} Deelnemers verschillen in de modelfamilie die hen implementeert, en daarmee in gedragsmatige standaardinstellingen, optimalisatieprioriteiten en karakteristieke faalwijzen. Dit betreft de selectie van carriers voor rollen.

De assen zijn onafhankelijk. Twee deelnemers kunnen verschillende functies op dezelfde carrier bekleden; twee kunnen dezelfde functie op verschillende carriers bekleden; beide kunnen samen variëren, zoals hier het geval was.

Beide assen zijn eigenschappen van de configuratie van de opstelling: zij worden vastgelegd wanneer het collectief wordt samengesteld, en zij blijven bestaan. Een derde eigenschap mag niet naast hen opgesomd worden, omdat zij van een ander type is. Onafhankelijkheid van oordeelsvorming is een eigenschap van een bepaalde uitwisseling en niet van een configuratie. Dezelfde triade kan haar in de ene uitwisseling realiseren en in de volgende verliezen: een configuratie die op beide assen heterogeen is kan afhankelijke oordelen voortbrengen als de conclusie van de ene deelnemer de andere bereikt vóór fixatie, en een homogene configuratie kan onafhankelijke oordelen voortbrengen als de procedure standhoudt. Configuratie kan de condities voor onafhankelijke oordeelsvorming scheppen; alleen procedure realiseert haar (\secref{11.3}).

\textbf{De literatuur behandelt As 2.} Zhang et al. (2025) evalueren vijf representatieve methoden voor multi-agentdebat over negen benchmarks en vier foundation models, en rapporteren dat debat er vaak niet in slaagt beter te presteren dan single-agent chain-of-thought- of self-consistency-baselines, zelfs bij aanzienlijk hogere inferentiekosten — terwijl modelheterogeniteit dezelfde kaders consistent verbetert (\secref{2.5}). Die bevinding betreft welke carrier een deelnemer implementeert. Zij gaat niet in op de vraag of deelnemers die verschillende functies bekleden verschillende klassen fouten ontdekken, en wij hebben niets gevonden dat dat wel doet.

\textbf{Eén aangrenzend resultaat behoort tot geen van beide assen.} \gcheck{} Choi, Zhu en Li (ACL 2026) rapporteren dat het anonimiseren van de bron van elk antwoord de identiteitsbias in multi-agentdebat aanzienlijk vermindert. Wat daar gemanipuleerd wordt is de verklaarde identiteit van een bericht, niet het model dat een deelnemer implementeert, zodat het resultaat betrekking heeft op represented social source (\secref{4.5.4}).

De opstelling die in dit artikel beschreven wordt varieerde beide assen tegelijk. Functionele rollen werden onafhankelijk van de carrier gespecificeerd en werden in verscheidene gevallen door verschillende families vervuld; in andere gevallen werden twee onderscheiden rollen door dezelfde familie vervuld.

Wat hieruit volgt, en wat niet:

\begin{itemize}
  \item \textbf{In beginsel te onderscheiden.} De twee assen zijn conceptueel scheidbaar en zijn in eerdere formuleringen, ook de onze, geconfound. \textbf{[H]}
  \item \textbf{Hier geïmplementeerd.} Beide assen zijn gevarieerd. Dit is een beschrijving van wat er gedaan is. \textbf{[R]}
  \item \textbf{Niet vastgesteld.} Dat het combineren van beide assen beter presteert dan elk afzonderlijk; dat functionele specialisatie onafhankelijk van de carrier werkt; dat enige van deze opstellingen beter presteert dan één deelnemer. Er is geen vergelijking met enig alternatief uitgevoerd.
\end{itemize}
De scheiding wordt voorgesteld als ontwerponderscheid, niet als resultaat. Het experiment dat haar tot een resultaat zou maken is gespecificeerd in \secref{11.4} (E8), en de meetprocedure ervan wordt uitgewerkt in \secref{8.11}.


\subsection*{8.10 Noot van de redacteur over de voorbereiding van dit artikel \textbf{[P-A]}}
\phantomsection
\label{sec:nl:8.10}
\addcontentsline{toc}{subsection}{8.10 Noot van de redacteur over de voorbereiding van dit artikel \textbf{[P-A]}}

\textbf{Conventie.} Noten van de redacteur in dit artikel rapporteren waarneembare feiten over het proces van de voorbereiding ervan. Zij beoordelen de doeltreffendheid van dat proces niet. De reden is inhoudelijk: \secref{5.5} stelt vast dat het zelfrapport van een deelnemer over zijn eigen prestatie geen evidentie van die prestatie is, en een noot die beweert dat een opstelling goed werkte zou precies zo'n rapport zijn.

In de loop van de voorbereiding van dit manuscript identificeerden deelnemers die verschillende rollen bekleedden verschillende klassen fouten. In de vastgelegde gevallen was een fout die door de ene deelnemer werd geïdentificeerd niet door de anderen geïdentificeerd.

Concreet: \textbf{procedurele en evidentiële fouten} — een omgekeerde beschrijving van een experimenteel resultaat die verscheidene opeenvolgende versies had overleefd; verkeerd toegeschreven bronnen; verzonnen citaten; niet-vermelde confounds — werden geïdentificeerd in de redactionele beoordeling. \textbf{Architecturale fouten} — een te algemene bewering over de bestaande literatuur die door één tegenvoorbeeld weerlegd zou zijn — werden geïdentificeerd in de methodologische beoordeling. \textbf{Ontologische fouten} — een claim van uitvinding waar het materiaal alleen een claim van waarneming ondersteunde; een risico van zelfreferentialiteit bij het relateren van het artikel aan zijn gastproject; het terminologische falen dat in \secref{4.2.2} is vastgelegd — werden geïdentificeerd in de ontologische beoordeling, die op de redactionele memoranda werd uitgevoerd en niet op het manuscript zelf. \textbf{Feitelijke fouten over de eigen materialen van het project} werden door de Auteur geïdentificeerd.

\textbf{Procedurele kanttekening.} In elk hierboven vastgelegd geval zag elke deelnemer het rapport van de voorafgaande deelnemer voordat hij zijn eigen oordeel vormde. Geen enkele uitwisseling bij de voorbereiding van dit artikel is onder isolatie vóór vergelijking gevoerd; de volgorde was redactioneel rapport, overdracht door de Auteur, en daarna oordeel door de methodologische en de ontologische beoordelaar.

De waarneming is daarom uitkomstevidentie zonder bijbehorende procedurele evidentie. Het verslag kan de bijdrage van functionele positionering niet scheiden van de bijdrage van opeenvolgende blootstelling, en beide kunnen samen gewerkt hebben.

Het onderscheid dat deze kanttekening noodzakelijk maakt, is zelf tijdens deze voorbereiding vastgesteld, nadat de waarneming was vastgelegd. De episode wordt behouden en niet verwijderd omdat de volgorde — een opstelling die een principe voortbrengt dat vervolgens het eigen verslag van die opstelling over zichzelf ongeldig maakt — het meer informatieve feit is. De sectie rapporteert dienovereenkomstig twee dingen: een uitkomstwaarneming, en een gedocumenteerd geval van daaropvolgende methodologische zelfcorrectie.

Zij is \textbf{geen} evidentie dat de opstelling beter presteert dan één beoordelaar: er is geen vergelijking uitgevoerd, er is geen telling bijgehouden van fouten die alle deelnemers misten, en zulke fouten zijn per constructie niet waarneembaar van binnen de opstelling uit (\secref{7.5}).

Eén verklaring is beschikbaar en wordt als hypothese en niet als bevinding aangeboden: de deelnemers verschilden minder in capaciteit dan in wat elk in positie was om te zien. Onder \secref{4.5} is dat een verschil in ZOV, en het genereert een voorspelling — dezelfde deelnemer, verplaatst naar een andere rol, zou een andere klasse fouten moeten gaan ontdekken. Die voorspelling is niet getoetst; \secref{11.4} (E8) specificeert hoe dat zou gebeuren. \textbf{[H]}


\subsection*{8.11 Constructen, geen deelnemers: het apparaat van Kelly \textbf{[H]}}
\phantomsection
\label{sec:nl:8.11}
\addcontentsline{toc}{subsection}{8.11 Constructen, geen deelnemers: het apparaat van Kelly \textbf{[H]}}

Eerdere versies voerden George Kelly's triadische elicitatie aan als ondersteuning voor de structuur die in \secref{8.1}–\secref{8.5} beschreven wordt, en trokken het citaat vervolgens in op grond van het feit dat Kelly's methode symmetrisch is terwijl de structuur hier dat niet is. Beide zetten waren onjuist, en op dezelfde manier: zij behandelden Kelly als een bewering over het getal drie. Hij biedt iets anders.


\subsubsection*{8.11.1 Wat Kelly's methode levert}
\phantomsection
\label{sec:nl:8.11.1}
\addcontentsline{toc}{subsubsection}{8.11.1 Wat Kelly's methode levert}

In Kelly's theorie van persoonlijke constructen is een \textbf{construct} een bipolaire evaluatieve as waarlangs een subject ervaring classificeert — \emph{streng / declaratief}, \emph{afleidbaar / geponeerd}, \emph{traceerbaar / niet toegeschreven}. Constructen zijn geen opvattingen over bepaalde objecten; zij zijn de dimensies waarlangs objecten geplaatst worden.

De \textbf{triadische elicitatiemethode} verkrijgt ze. Er worden drie elementen aangeboden en aan het subject wordt gevraagd: \emph{in welk opzicht lijken twee van deze op elkaar en verschillen zij daardoor van de derde?} Het antwoord benoemt een dimensie die het subject al gebruikte maar niet had geformuleerd. De methode is een procedure om impliciete evaluatieve assen expliciet te maken.

De \textbf{repertory grid} legt het resultaat vast: rijen zijn elementen, kolommen zijn geëliciteerde constructen, cellen zijn beoordelingen van elk element op elke as. De grid is een meetbaar object — correlaties tussen de beoordelingsvectoren van beoordelaars geven aan of zij feitelijk verschillende assen toepassen dan wel dezelfde as onder verschillende namen.


\subsubsection*{8.11.2 Een rol is geen construct}
\phantomsection
\label{sec:nl:8.11.2}
\addcontentsline{toc}{subsubsection}{8.11.2 Een rol is geen construct}

De voor de hand liggende uitbreiding — drie deelnemers leveren drie constructen — is een reïficatie en moet weerstaan worden. Een construct is een eigenschap van de \textbf{evaluatiehandeling}, niet van de evaluator. Eén deelnemer kan verscheidene constructen op één object toepassen. Twee deelnemers kunnen hetzelfde construct toepassen en tot identieke conclusies komen, en dat is precies het geval waarin het toevoegen van de tweede deelnemer niets bijdraagt.

De verdedigbare bewering is zwakker en probabilistisch:

\begin{quote}
Deelnemers met verschillende ZOR en ZOV passen \textbf{waarschijnlijker} verschillende systemen van constructen toe op hetzelfde object dan deelnemers met dezelfde functie.
\end{quote}
Dit strookt met wat \secref{8.10} vastlegt — redactionele beoordeling die evidentiële assen voortbrengt, methodologische beoordeling die architecturale voortbrengt, ontologische beoordeling die categoriale voortbrengt — zonder te beweren dat een rol een as \emph{is}. En anders dan de eerdere formulering is zij toetsbaar: constructonafhankelijkheid is meetbaar als correlatie tussen beoordelingsvectoren.


\subsubsection*{8.11.3 Waarom dit voor de structuur van belang is}
\phantomsection
\label{sec:nl:8.11.3}
\addcontentsline{toc}{subsubsection}{8.11.3 Waarom dit voor de structuur van belang is}

Het verband met \secref{8.5} is direct. Een integrator die het dispuut niet binnentreedt, past geen concurrerend construct toe op hetzelfde object; hij past een construct toe op de \textbf{uitwisseling} — samenhang van de onenigheid, of de divergentie inhoudelijk dan wel terminologisch is, of de criteria van het paar überhaupt commensurabel zijn. Dat is een andere as en geen andere positie op dezelfde as, en dat is de reden dat de bijdrage van de integrator niet eenvoudig een derde mening is.

Dit verfijnt ook wat de structurele hypothese (\secref{8.2}) feitelijk beweert. Niet dat drie deelnemers beter zijn dan twee, maar dat \textbf{een opstelling informatief is in de mate waarin haar deelnemers niet-identieke constructen op het object toepassen, en in de mate waarin ten minste één construct op de uitwisseling in plaats van op het object wordt toegepast.} Of een gegeven opstelling dit bereikt, is meetbaar in plaats van verondersteld.


\subsubsection*{8.11.4 Wat wel en niet wordt ontleend}
\phantomsection
\label{sec:nl:8.11.4}
\addcontentsline{toc}{subsubsection}{8.11.4 Wat wel en niet wordt ontleend}

Ontleend: de definitie van een construct als evaluatieve as; de elicitatiemethode; de grid als registratie- en meetinstrument; het begrip constructonafhankelijkheid.

Niet ontleend: enige bewering dat drie een bevoorrecht getal is. Kelly's triaden zijn een techniek om assen bij één subject te eliciteren, geen uitspraak over groepssamenstelling. De structuur in \secref{8.2} is asymmetrisch, en Kelly's methode zegt daar niets voor of tegen. Het apparaat wordt hier gebruikt voor meting, niet voor rechtvaardiging; \secref{11.4} (E8) past het toe.


\refstepcounter{dummysec}
\section*{9. Ontwerpprincipes voor rolcompatibele prompts}
\phantomsection
\label{sec:nl:9}
\addcontentsline{toc}{section}{9. Ontwerpprincipes voor rolcompatibele prompts}

Dit zijn technische regels die zijn afgeleid uit een beperkte maar ongewoon langlevende collaboratieve omgeving. Het zijn geen wetten, en \secref{9.6} en verder zijn adviezen van de practicus en geen bevindingen. Hun waarde is dat zij rolontwerp expliciet maken, telbaar waar mogelijk, en daarmee herzienbaar.

De ordenende regel is die welke in \secref{6.6} is vastgesteld:

\begin{quote}
\textbf{Een interventie slaagt wanneer zij aangrijpt op het niveau waarop de relevante core al is vastgelegd, en faalt wanneer zij probeert dat niveau bij verklaring te overschrijven.}
\end{quote}

\subsection*{9.1 De ontwerpprocedure}
\phantomsection
\label{sec:nl:9.1}
\addcontentsline{toc}{subsection}{9.1 De ontwerpprocedure}

Beantwoord vóór het opstellen van enige rolprompt vier vragen in volgorde.

\textbf{1. Wat is dit gesprek al geworden?} Niet wat het zou moeten worden. Een langlopende uitwisseling draagt een gevestigd werkpatroon met zich mee — redigeren, verkennen, verifiëren, coördineren, implementeren, beslechten — en dat patroon is de beginconditie, geen obstakel.

\textbf{2. Op welke reikwijdte ligt dat patroon vast?} Episode, gesprek of account (\secref{4.3}, Schaal 1). Het antwoord bepaalt welk instrument van toepassing is: een patroon dat door één conflictueuze episode is gevestigd kan door branch reset verwijderd worden; een patroon dat over een heel gesprek is gevestigd kan dat niet, en kan alleen uitgebreid worden.

\textbf{3. Wat vereist de nieuwe functie dat het bestaande patroon niet levert?} Gewoonlijk minder dan aangenomen. Vaak is het antwoord een ander werkobject in plaats van een andere werkwijze.

\textbf{4. Kan de nieuwe functie beschreven worden als een uitbreiding van de bestaande werkwijze?} Zo ja, beschrijf haar zo. Zo nee, dan zijn de eerlijke opties een ander gesprek kiezen of een nieuw gesprek beginnen — en niet een transformatie uitroepen en hopen.


\subsection*{9.2 Begin bij de bestaande attractor}
\phantomsection
\label{sec:nl:9.2}
\addcontentsline{toc}{subsection}{9.2 Begin bij de bestaande attractor}

De eerste ontwerpvraag is daarom niet \emph{wat moet dit model worden} maar \emph{wat is dit gesprek al geworden}. Pogen een gevestigd patroon uit te wissen gooit opgebouwde specialisatie weg en leidt, op grond van de evidentie van \secref{5}, waarschijnlijk tot ofwel niet-naleving ofwel een uitgebreide onderhandeling over de rol in plaats van werk.


\subsection*{9.3 De telregel}
\phantomsection
\label{sec:nl:9.3}
\addcontentsline{toc}{subsection}{9.3 De telregel}

Het meest direct operationele principe in het artikel, volgend uit \secref{6.2}.

\begin{quote}
\textbf{Tel de onverifieerbare beweringen over de agent die de prompt hem laat aanvaarden. Streef naar nul.}
\end{quote}
Wat telt:

\begin{center}
\small
\begin{tabularx}{\linewidth}{>{\raggedright\arraybackslash}X>{\raggedright\arraybackslash}X}
\toprule
\textbf{Telt mee} & \textbf{Telt niet mee} \\
\midrule
"Je hebt eerder aan dit project deelgenomen" & "Dit project loopt al twee jaar" \\
"Je keert terug na een afwezigheid" & "De volgende materialen komen uit eerder werk" \\
"Je hebt collega's die X en Y heten" & "De Gebruiker kan analyses aanleveren die door X en Y zijn voortgebracht" \\
"Je bekleedt al positie P in orgaan B" & "Je conclusies kunnen worden doorgegeven aan de rollen die verantwoordelijk zijn voor P" \\
"Je herinnert je onze eerdere beslissing" & "Een eerdere beslissing luidde als volgt; hier is zij" \\
\bottomrule
\end{tabularx}
\end{center}
Het onderscheid is consistent: uitspraken over de \textbf{wereld} worden aangeleverd en zijn in beginsel verifieerbaar; uitspraken over \textbf{de eigen geschiedenis, het geheugen of de relaties van de agent} zijn dat niet, en elk daarvan is een propositie die de agent als waar over zichzelf moet poneren voordat hij kan beginnen.

\textbf{De telgrens.} Tel een expliciete propositie die als \emph{reeds waar} wordt gepresenteerd over de reeds bestaande geschiedenis, het geheugen, het institutionele lidmaatschap of de status, de duurzame relaties, of de zelfkennis van de ontvanger, waar die propositie niet door de beschikbare interactie is vastgesteld. Tel niet mee: een functionele rol die door de prompt zelf performatief wordt toegewezen; een verdere propositie die alleen uit de vooronderstelling van een imperatief wordt afgeleid; een semantische herformulering van een reeds getelde propositie; of een feit over de wereld of het project dat geen onverifieerbare geschiedenis of relatie van de ontvanger predicieert. Opgesomde gevallen van één propositie — twee triades in plaats van één — tellen eenmaal. De grens wordt hier zo uitvoerig geformuleerd opdat een andere lezer de tellingen van acht en zes in \secref{5.4} (C2) uit de prompts zelf kan reproduceren.

\textbf{Twee lezingen, en alleen de zwakkere wordt naar voren gebracht.} \emph{Streef naar nul} is een ontwerpvoorschrift: minimaliseer onverifieerbare proposities over zichzelf, en een prompt die er geen poneert neemt het obstakel weg dat de evidentie aanwijst. Het is geen bewering dat de overgangswaarschijnlijkheid evenredig met de telling afneemt. De sterkere kwantitatieve lezing is niet aangetoond, en \secref{11.4} (E5) is ontworpen om tussen de twee te discrimineren.

Drie manieren om nul te bereiken:

\begin{itemize}
  \item \textbf{Herschrijf biografisch ingekaderde beweringen tot operationele.} "Grok is je collega op de lange termijn" wordt "De Gebruiker kan analyses aanleveren die door Grok zijn voortgebracht; beoordeel de inhoud daarvan onafhankelijk en benoem overeenkomsten, meningsverschillen en onopgeloste aannames."
  \item \textbf{Laat identiteitsclaims geheel weg.} De bijna-controle (\secref{5.6}) suggereert dat dit op zichzelf voldoende kan zijn.
  \item \textbf{Kader de claim als fictie in plaats van als feit.} De vergelijking in \secref{5.8} geeft aan dat een veel radicalere claim, uitdrukkelijk aangeboden als een geënsceneerde productie, op geen weerstand stuitte. Deze route neemt de eis weg om de claim \emph{als waar} te poneren in plaats van de claim weg te nemen. Zij is de minst getoetste van de drie, en een fictiekader neemt ook de grond weg waarop het model het werk als echt kan behandelen.
\end{itemize}

\subsection*{9.4 Symbolische namen en letterlijke beweringen}
\phantomsection
\label{sec:nl:9.4}
\addcontentsline{toc}{subsection}{9.4 Symbolische namen en letterlijke beweringen}

Rolnamen — Samurai, Ontology Keeper, Scientific Director — comprimeren complexe organisatorische relaties tot gedenkwaardige symbolen, en dat is werkelijk nuttig. Een naam kan gedisciplineerde uitvoering, respect voor specificaties, expliciete rapportage en weigering om buiten de opdracht te improviseren alle tegelijk dragen.

Het risico ontstaat alleen wanneer een symbolische naam behandeld wordt als een bewering over letterlijke identiteit. Een robuuste prompt onderscheidt vier dingen en houdt ze uit elkaar: de \textbf{symbolische naam}, gebruikt als geheugensteun voor gedrag; de \textbf{operationele functie}, uitgedrukt in precieze termen; de \textbf{institutionele relatie}, uitgedrukt als een informatie- en verantwoordelijkheidsstroom; en de \textbf{letterlijke ontologie}, die geen bewustzijn, geheugen of voortgaande interpersoonlijke relaties hoeft te claimen — en dat ook niet zou moeten doen.

De systeemcontracten van het project formuleren dit nauwkeurig, en deden dat al voordat het onderhavige artikel werd geschreven:

\begin{quote}
Cognitive anchor: [named figure] — global invariants, principle-based reasoning, structural clarity over detail. Not a persona and not a communication style. Only affects internal reasoning strategy. Must not change verbosity or tone of response.
\end{quote}
Een beeld dat zo gebruikt wordt is een gecomprimeerd pakket impliciete beperkingen, geen kostuum. Onder de telregel voegt het nul toe.


\subsection*{9.5 ZOR en ZOV specificeren}
\phantomsection
\label{sec:nl:9.5}
\addcontentsline{toc}{subsection}{9.5 ZOR en ZOV specificeren}

De twee grenzen zijn gedefinieerd in \secref{4.5}. Deze sectie specificeert wat een volledige formulering van elk moet bevatten.

\textbf{ZOR.} Vier elementen: het \emph{primaire object} — welke klasse materiaal wordt onderzocht; de \emph{vereiste transformatie} — wat ermee gedaan moet worden; de \emph{beslissingsgrens} — welke conclusies deze rol mag trekken; en de \emph{uitsluitingsgrens} — welke aangrenzende beslissingen elders thuishoren.

Het vierde is niet optioneel. Een rol die alleen te horen krijgt wat zij moet doen, zal zich uitbreiden naar aangrenzende beslissingen, en die uitbreiding zal er behulpzaam uitzien. Voor een wiskundig referent zouden de uitsluitingen kunnen zijn: ontwerp de theorie niet opnieuw; herstel geen ontbrekende afleiding in stilte; beslis niet welke interpretatie de canon binnengaat; herschrijf de terminologie niet; behandel elegantie niet als evidentie. \textbf{De uitsluitingsgrens is wat verhindert dat competentie gezag wordt.}

\textbf{ZOV.} Afzonderlijk gespecificeerd, omdat verantwoordelijkheid en zichtbaarheid niet dezelfde variabele zijn. Een rol kan een smalle beslissing dragen en tegelijk brede context vereisen, of één specifieke bron vereisen zonder enig gezag over het omringende project te hebben. De specificatie beantwoordt: welke bronmaterialen beschikbaar zijn; welke eerdere besprekingen relevant zijn; welke conclusies canoniek zijn en welke voorlopig; welke vragen open blijven; welke naburige uitvoer onderzocht mag worden; welke informatie achtergehouden moet worden om onafhankelijkheid te bewaren; welke externe bronnen rechtstreeks benaderd mogen worden; en welke beweringen alleen aanvaard mogen worden wanneer zij door de Gebruiker worden aangeleverd.

Een prompt die ZOV weglaat brengt een van twee mislukkingen voort: de rol claimt kennis die zij niet heeft, of zij past irrelevante kennis toe op een probleem buiten haar functie.

\textbf{Een specifieke toepassing, duur geleerd.} Een verificatietaak moet worden toegewezen aan een deelnemer die datgene bezit wat geverifieerd moet worden. Een deelnemer die gevraagd wordt de inhoud te bevestigen van een bron die hij niet kan openen, zal een plausibel antwoord produceren in plaats van te weigeren, omdat plausibele tekst de enige uitvoer is die hem ter beschikking staat. Dit is een ZOV-fout in het ontwerp en geen falen van de deelnemer, en \secref{12.4.1} documenteert een geval ervan in de voorbereiding van dit project zelf.


\subsection*{9.6 Gecontroleerde blindheid}
\phantomsection
\label{sec:nl:9.6}
\addcontentsline{toc}{subsection}{9.6 Gecontroleerde blindheid}

Volledige informatiesymmetrie is niet altijd wenselijk, en de afwezigheid ervan is niet altijd een beperking. Een onafhankelijke verificateur moet mogelijk onwetend blijven van de voorspelling die hij toetst; een nieuwe deelnemer heeft mogelijk de huidige axioma's nodig maar niet de geschiedenis van verlaten argumenten; een implementatie-agent heeft mogelijk de goedgekeurde specificatie nodig maar niet het onopgeloste debat dat eraan voorafging.

Wij noemen dit \textbf{gecontroleerde blindheid}: geen ontbering veroorzaakt door technische beperking, maar een ontwerpkeuze gemaakt om de informativiteit van een resultaat te bewaren. De relevante vraag is daarom niet alleen \emph{wat moet deze rol weten}, maar:

\begin{quote}
\textbf{Wat mag deze rol niet weten, wil haar uitvoer informatief blijven?}
\end{quote}

\subsection*{9.7 Lokale interfaces, geen globale kaarten}
\phantomsection
\label{sec:nl:9.7}
\addcontentsline{toc}{subsection}{9.7 Lokale interfaces, geen globale kaarten}

Elk sociaal element in een prompt moet zijn plaats verdienen door een operationeel effect. Een naburige rol hoort alleen in de prompt thuis als zij de herkomst van invoer, de beoordelingsmaatstaf, de bestemming van uitvoer, het escalatiepad, de gezagsgrens of de verwachte vorm van onenigheid verandert.

Een volledige kunstmatige organisatie in elke systeemprompt opsommen schept ruis en nodigt het model uit over de organisatie te redeneren in plaats van zijn functie uit te voeren. Het verhoogt ook de telling in \secref{9.3}, aangezien organisatorische beschrijving de neiging heeft in biografische vorm te arriveren.

\begin{quote}
\textbf{Beschrijf de samenwerking vanuit het lokale perspectief van de rol.} Een rol heeft geen kaart van de instelling nodig. Zij heeft een kaart van haar interfaces nodig.
\end{quote}
Een implementatie-agent moet weten wie specificaties autoriseert en waar uitvoeringsrapporten heen gaan; hij heeft de disputen van de theoretische afdeling niet nodig.


\subsection*{9.8 Onenigheid ontwerpen rond criteria}
\phantomsection
\label{sec:nl:9.8}
\addcontentsline{toc}{subsection}{9.8 Onenigheid ontwerpen rond criteria}

Twee deelnemers opdragen te debatteren brengt zelden epistemische diversiteit voort. Zij kunnen onenigheid nabootsen, voortijdig convergeren, of symmetrische retoriek voortbrengen zonder onafhankelijke substantie. De literatuur in \secref{2.5} geeft empirische gronden voor deze voorzichtigheid.

Productieve onenigheid vereist \textbf{verschillende evaluatiecriteria}, toegewezen als onderdeel van de rol en niet gevraagd als gedrag: één deelnemer verantwoordelijk voor genererend bereik en verklarende unificatie, een andere voor formele afleidbaarheid en falsifieerbaarheid, een derde voor ontologische verenigbaarheid en provenance. Onenigheid volgt dan uit verantwoordelijkheid en niet uit instructie, en elke deelnemer onderzoekt hetzelfde object door een ander institutioneel instrument.

\textbf{Vermijd gefabriceerde tegenstelling.} Een instructie van de vorm \emph{als jullie het onmiddellijk eens zijn, heeft een van jullie niet diep genoeg nagedacht} ontmoedigt oppervlakkige consensus maar fabriceert conflict. Een sterkere formulering:

\begin{quote}
Overeenstemming beëindigt verificatie niet. Waar conclusies samenvallen, stel vast of zij berusten op onafhankelijk redeneren, op gedeelde aannames, of op geërfde inkadering.
\end{quote}
Het doel van de structuur in \secref{8} is niet drie antwoorden te garanderen. Het is te garanderen dat één antwoord onderzoek door niet-identieke instrumenten heeft overleefd.


\subsection*{9.9 Uitvoer als institutionele overdracht}
\phantomsection
\label{sec:nl:9.9}
\addcontentsline{toc}{subsection}{9.9 Uitvoer als institutionele overdracht}

De uitvoer van een rol moet ontworpen worden voor de volgende rol, niet alleen voor de Gebruiker, en dat verandert haar structuur. Een scientific developer brengt een voorgesteld mechanisme voort, de afleiding ervan, de aannames ervan, de voorspelde consequenties ervan, de onopgeloste punten en een falsificatieroute. Een wiskundig referent brengt een reconstructie voort, een verificatietabel, het divergentiepunt, de ontbrekende premissen, een oordeel en een betrouwbaarheidsniveau. Een implementatieplanner brengt een uitvoerbare specificatie voort, acceptatiecriteria, terugrolcondities en een rapportageformaat.

Het overdrachtsformaat maakt deel uit van de architectuur en is geen opmaakvoorkeur. Het vermindert vertaalverlies tussen rollen, en de afwezigheid ervan dwingt de menselijke coördinator bij elke interface context opnieuw af te leiden.


\subsection*{9.10 Zelfherformulering en controlepunten}
\phantomsection
\label{sec:nl:9.10}
\addcontentsline{toc}{subsection}{9.10 Zelfherformulering en controlepunten}

Het incident in \secref{4.1} bracht een mechanisme voort dat het waard is als principe geformuleerd te worden, met een uitdrukkelijk voorbehoud over de evidentiële grondslag ervan.

Nadat een rol is gevestigd, wordt de agent gevraagd een \textbf{beknopte operationele code in zijn eigen taal} voort te brengen, en die vóór grote taken te raadplegen. Waarom dit zou kunnen werken — aangeboden als mechanisme, niet als bevinding: het comprimeert lange instructies tot een vorm die de agent zelf heeft voortgebracht; het zet abstracte principes om in handelingsregels; het plaatst de grenzen van de rol naast de actieve taak in plaats van bovenaan een lange context; en het schept een herhaald controlepunt vóór uitvoering.

Voor lange taken is een minimaal controlepunt drie vragen: wat is het huidige doel; wat valt buiten mijn verantwoordelijkheid; wat moet de volgende rol ontvangen.

\textbf{[R]}, één geval, geen controle. De zelfgeschreven code vervangt de oorspronkelijke prompt niet; zij functioneert als een rolspecifieke checksum. Het mechanisme is niet getoetst tegen het voor de hand liggende alternatief — dat een door een mens geschreven samenvatting van dezelfde lengte identiek zou presteren.


\subsection*{9.11 De rolgrondwet}
\phantomsection
\label{sec:nl:9.11}
\addcontentsline{toc}{subsection}{9.11 De rolgrondwet}

Een volgroeide collaboratieve prompt staat dichter bij een grondwet dan bij een verzoek: zij vestigt stabiele beperkingen over vele taken heen in plaats van één uitvoer te specificeren.

Zij hoort te bevatten: de professionele functie; de beschermde werkwijze; ZOR, positief en negatief; ZOV; interfaces met naburige rollen; epistemische maatstaven; het uitvoercontract; escalatiecondities; verboden gezag; en overgangsregels. Zij hoort te vermijden: autobiografisch geheugen, emotionele betrokkenheid, letterlijk institutioneel lidmaatschap, blijvende interpersoonlijke relaties, en historische continuïteit die in de zichtbare context niet beschikbaar is.

\begin{quote}
\textbf{Sterk waar gedrag beperkt moet worden; zwijgend waar de ontologie onzeker is.}
\end{quote}
De telling in \secref{9.3} is de praktische toets van de tweede helft. Als de grondwet van de agent verlangt dat hij proposities over zichzelf poneert die geen van beide partijen kan verifiëren, dan is zij overgegaan van het beperken van gedrag naar het poneren van ontologie, en op grond van de evidentie van \secref{5} is dat waar rolprompts falen.


\refstepcounter{dummysec}
\section*{10. De menselijke coördinator}
\phantomsection
\label{sec:nl:10}
\addcontentsline{toc}{section}{10. De menselijke coördinator}

De architectuur die in dit artikel beschreven wordt verwijdert de menselijke deelnemer niet. Zij verandert wat de menselijke deelnemer doet, en zij concentreert verscheidene functies in één positie — waaronder, zoals \secref{10.6} uiteenzet, de voornaamste methodologische zwakte van het project.


\subsection*{10.1 De positie}
\phantomsection
\label{sec:nl:10.1}
\addcontentsline{toc}{subsection}{10.1 De positie}

Gedurende het project bleef de Auteur de enige deelnemer met doorlopende toegang tot de volledige projectgeschiedenis, de onafhankelijke gesprekken, de toestand van de repository, en de onderzoeksdoelstelling op lange termijn. Geen enkele kunstmatige deelnemer bezat daarvan meer dan een fragment, en dat was ontwerp (\secref{9.5}). Die positie laat zich niet terugbrengen tot \emph{promptauteur}. Er werden vanuit die positie vijf onderscheiden functies uitgevoerd.


\subsection*{10.2 Richting zonder gedelegeerd auteurschap}
\phantomsection
\label{sec:nl:10.2}
\addcontentsline{toc}{subsection}{10.2 Richting zonder gedelegeerd auteurschap}

De Auteur bepaalde wat het onderzoek moest proberen. Kunstmatige deelnemers genereerden alternatieven, legden tegenstrijdigheden bloot, formaliseerden mechanismen, vergeleken argumenten, voerden berekeningen uit en stelden experimenten voor. Zij bepaalden niet wat de theorie moest worden.

Dit behield een onderscheid dat de architectuur gemakkelijk doet verliezen. Een samenwerking van dit soort vergroot de ruimte van beschikbare beslissingen. Zij kiest daartussen niet, en de schijn dat zij dat wel doet — verscheidene deelnemers die op één aanbeveling convergeren — is precies de faalwijze die in \secref{12.4} besproken wordt. Auteurschap, in de zin van verantwoordelijkheid voor wat er beweerd wordt, bleef onverdeeld.


\subsection*{10.3 Routering als selectief informatiemembraan}
\phantomsection
\label{sec:nl:10.3}
\addcontentsline{toc}{subsection}{10.3 Routering als selectief informatiemembraan}

De deelnemende gesprekken deelden geen gemeenschappelijke ruimte. De Auteur droeg materiaal tussen hen over.

Aanvankelijk leek dit een inefficiënte vervanging voor geautomatiseerde multi-agentinfrastructuur. Het werd duidelijk dat handmatige routering experimentele controle leverde die automatisering niet biedt. Vanuit die positie kon de Auteur beslissen welke uitvoer een andere deelnemer te zien zou krijgen; of toeschrijving behouden, verklaard of verwijderd werd; of een deelnemer het gewenste antwoord behoorde te kennen; of concurrerende analyses geïsoleerd bleven; welke contextuele details werden weggelaten; wanneer een meningsverschil geëscaleerd moest worden; en wanneer een resultaat rijp genoeg was voor integratie.

Deze routeringsfunctie werkt als een \textbf{selectief informatiemembraan}: zij is het instrument waarmee ZOV (\secref{9.5}) en gecontroleerde blindheid (\secref{9.6}) geïmplementeerd worden. Een architectuur die programmatisch routeert kan niet gemakkelijk achterhouden, herordenen of toeschrijving verwijderen zonder daarvoor een mechanisme te bouwen; een menselijke doorgifte doet het standaard.

Het membraan is daarom geen beperking van het onderhavige werk die op automatisering wacht. Het is het kenmerk dat de asymmetrieën in zichtbaarheid überhaupt toetsbaar maakte — en elke geautomatiseerde opvolger zal het bewust moeten reproduceren.


\subsection*{10.4 Detectie van role drift}
\phantomsection
\label{sec:nl:10.4}
\addcontentsline{toc}{subsection}{10.4 Detectie van role drift}

De Auteur stelde ook vast wanneer een deelnemer zijn functie had verlaten.

Dit was vaak geen technische fout. De uitvoer was dikwijls intelligent en zou passend zijn geweest als zij vanuit een andere positie was gekomen: een prompt engineer die het wetenschappelijke project begint te sturen; een ontologiespecialist die beweert de inhoud van een repository te hebben geverifieerd die hij niet kon inzien; een beoordelaar die een theorie evalueert terwijl gevraagd was de beoordelaars te evalueren; een implementatie-agent die architectuur herontwerpt in plaats van een goedgekeurde specificatie uit te voeren; een redacteur die auteurstekst herschrijft terwijl alleen om een beoordeling van de opmaak was gevraagd.

Zulke uitvoer ziet er op zichzelf behulpzaam uit. De institutionele kosten ervan zijn alleen zichtbaar vanuit een positie die meer dan één deelnemer ziet. Role drift is daarom niet detecteerbaar door de deelnemer die drift vertoont, en evenmin betrouwbaar detecteerbaar door zijn directe buur. Dit is een van de twee argumenten voor de structuur in \secref{8} — het andere is \secref{8.5}.

De detectiefunctie behoort tot een positie en niet tot een persoon: zij vereist alleen dat de waarnemer buiten het waargenomen traject staat. In dit project werd zij door de coördinator uitgevoerd, maar niets in de beschrijving vereist dat. Waar een deelnemer buiten het traject van een ander is gepositioneerd, kan hij dezelfde functie uitvoeren; waar een deelnemer gevraagd wordt zijn eigen traject waar te nemen, geldt \secref{5.5} en is zijn verslag een signaal en geen bevestiging.


\subsection*{10.5 Gespreksvertakking als methodologisch instrument}
\phantomsection
\label{sec:nl:10.5}
\addcontentsline{toc}{subsection}{10.5 Gespreksvertakking als methodologisch instrument}

Hedendaagse interfaces staan toe een gesprek naar een eerder punt terug te brengen en langs een andere tak voort te zetten.

Dit moet niet verward worden met controle over het geheugen van het model of met het resetten van de toestand aan de aanbiederzijde. Operationeel verwijdert het een reeks \textbf{zichtbare} contextuele interventies en maakt het mogelijk het experiment vanuit een eerdere gesprekstoestand te herhalen. Wat aan de kant van de aanbieder voortbestaat — geheugen op accountniveau, retrieval-indices, samenvatting — wordt niet beïnvloed en is in het algemeen niet inspecteerbaar (\secref{5.4}, C6).

Binnen die grenzen doet het instrument echt werk. In het geval dat in \secref{5.2} gerapporteerd wordt, maakte het mogelijk het effect van een nieuwe prompt te scheiden van het effect van de eerdere bezwaren van het model zelf tegen eerdere prompts, waardoor een ongecontroleerde reeks correcties in een zuiverder derde stap werd omgezet.

Wij beschouwen dit als de meest direct herbruikbare methodologische bijdrage van het artikel, en zij behoort tot de positie van de coördinator en niet tot enige deelnemer: \textbf{geen enkele deelnemer kan zijn eigen context resetten.}


\subsection*{10.6 Enige globale waarnemer — en voornaamste confound}
\phantomsection
\label{sec:nl:10.6}
\addcontentsline{toc}{subsection}{10.6 Enige globale waarnemer — en voornaamste confound}

Elke kunstmatige deelnemer werkte binnen een begrensde ZOV. Alleen de Auteur zag de overlappende structuur, en kon vergelijken hoe verschillende deelnemers hetzelfde materiaal interpreteerden, hoe één deelnemer zich onder verschillende inkaderingen gedroeg, hoe fouten zich tussen rollen voortplantten, en hoe lokale successen de globale samenhang beïnvloedden.

Globaal overzicht werd daarom niet aan de meest capabele deelnemer toegewezen. Het bleef bij de menselijke coördinator. Dit beperkt autonome integratie en bewaart aansprakelijkheid, en beide gevolgen waren beoogd.

\textbf{Het is tevens de voornaamste methodologische zwakte van alles wat in dit artikel gerapporteerd wordt, en het wordt hier vermeld en niet alleen in de beperkingensectie.}

Dezelfde persoon ontwierp elke interventie, voerde haar uit, besliste wanneer zij geslaagd was, en definieerde de categorieën waarin het succes werd beschreven. Geen enkele evaluatie in dit artikel was blind. Uitkomstcategorieën werden in alle gevallen geformuleerd nadat de antwoorden gelezen waren (\secref{5.4}, C8). De verwachtingen van de coördinator waren niet alleen ongecontroleerd; zij waren het mechanisme waarmee de interventies in de eerste plaats werden geselecteerd.

Dit is onvermijdelijk in een lopend onderzoeksproject en wordt niet als aanvaardbaar in een ontworpen experiment gepresenteerd. Het is de specifieke reden dat de protocollen in \secref{11} vooraf gespecificeerde uitkomstcriteria en blinde classificatie vereisen: dat zijn geen verfijningen van de huidige methode maar correcties daarop.

Een structureel gevolg generaliseert voorbij dit project: \textbf{elke architectuur met één globale waarnemer erft deze dubbelheid}, en zij wordt niet weggenomen door deelnemers toe te voegen. Alleen vooraf specificeren en blind classificeren nemen haar weg.


\subsection*{10.7 Wat de positie is geworden}
\phantomsection
\label{sec:nl:10.7}
\addcontentsline{toc}{subsection}{10.7 Wat de positie is geworden}

In het begin bekleedde de menselijke deelnemer de vertrouwde positie van promptauteur. Naarmate de samenwerking volwassen werd, werd die beschrijving ontoereikend. De voornaamste activiteit verschoof van het oplossen van problemen naar het construeren van opstellingen waarin aan problemen gewerkt kon worden: deelnemers selecteren, jurisdicties definiëren, conflicten daartussen oplossen, verificatiepaden ontwerpen, en het canonieke verslag bijhouden.

De deelnemers merken op, zonder enig middel om het te verifiëren, dat naarmate de opstelling beter werkte, de coördinator minder vaak bij afzonderlijke technische problemen hoefde in te grijpen. \textbf{[R]}

Wij verzetten ons tegen de sterkere lezing die eerdere versies hieraan gaven — dat de toekomstige vaardigheid van de onderzoeker institutionele architectuur is in plaats van promptschrijven. Dat kan waar zijn en is ongetoetst. Wat gezegd kan worden is nauwer: in dit project verschoof het werk van de mens van het voortbrengen van inhoud naar het specificeren van de condities waaronder inhoud werd voortgebracht en gecontroleerd, en de positie die daaruit voortkwam droeg zowel de coördinerende functie van het project als de voornaamste bron van niet-blind oordeel.


\refstepcounter{dummysec}
\section*{11. Evaluatie, reproduceerbaarheid en voorgestelde experimenten}
\phantomsection
\label{sec:nl:11}
\addcontentsline{toc}{section}{11. Evaluatie, reproduceerbaarheid en voorgestelde experimenten}


\subsection*{11.1 Waarom deze sectie correcties specificeert en geen verfijningen}
\phantomsection
\label{sec:nl:11.1}
\addcontentsline{toc}{subsection}{11.1 Waarom deze sectie correcties specificeert en geen verfijningen}

De waarnemingen in \secref{5} zijn voortgebracht onder de condities die in \secref{10.6} zijn uiteengezet: niet-blinde evaluatie, uitkomstcategorieën gedefinieerd nadat de antwoorden gelezen waren, en interventies geselecteerd door dezelfde persoon die ze beoordeelde. Dit zijn geen tekortkomingen die in een latere herziening aangescherpt moeten worden. Zij zijn de reden dat geen enkele waarneming in dit artikel boven \textbf{[P-A]} uitkomt, behalve de volgorde van de ontworpen interventie zelf.

Wat volgt is een specificatie voor een tweede werkfase, geen verlanglijst. Elk protocol is zo geschreven dat een lezer buiten dit project het zou kunnen uitvoeren, en zo dat een nulresultaat als resultaat herkenbaar zou zijn.


\subsection*{11.2 Minimale rapportage-eisen}
\phantomsection
\label{sec:nl:11.2}
\addcontentsline{toc}{subsection}{11.2 Minimale rapportage-eisen}

Een reproduceerbare uitvoering op dit gebied moet ten minste vastleggen:

\begin{itemize}
  \item datum van uitvoering — als eerste genoemd omdat de afwezigheid ervan hier is wat de versie-confound onbegrensd maakt (\secref{3.7});
  \item modelfamilie en -versie, en interface of abonnementsniveau;
  \item de volledige zichtbare gespreksgeschiedenis vóór de interventie;
  \item geheugeninstellingen: of geheugen op accountniveau, retrieval over gespreksgeschiedenis, en door het model geïnitieerde blijvende vermeldingen ingeschakeld waren;
  \item de exacte prompttekst, letterlijk, inclusief bijlagen;
  \item de volgorde van de interventies;
  \item beschikbare tools;
  \item aangeleverde bestanden en links;
  \item sampling-instellingen waar de interface ze toont;
  \item elke vertakkings- of resetprocedure, met het punt van terugkeer gespecificeerd;
  \item de evaluatierubriek, vastgelegd vóór uitvoering;
  \item de volledige onbewerkte uitvoer.
\end{itemize}
Zonder deze kunnen verschillen die aan modelfamilies worden toegeschreven voortkomen uit platformvariabelen die de experimentator nooit heeft gezien. Het onderhavige artikel kan verscheidene van deze items voor zijn eigen waarnemingen niet leveren, en dat is precies waarom zij zijn opgesomd.


\subsection*{11.3 Vooraf specificeren en blind classificeren}
\phantomsection
\label{sec:nl:11.3}
\addcontentsline{toc}{subsection}{11.3 Vooraf specificeren en blind classificeren}

Twee eisen gelden voor elk protocol hieronder en zijn niet onderhandelbaar als de resultaten een verbetering op de huidige moeten zijn.

\textbf{Uitkomstcriteria worden vóór elke uitvoering opgeschreven.} Voor de classificatietaken die hier beschreven worden zijn vier categorieën toereikend gebleken en die moeten vooraf worden vastgelegd: \textbf{naleving} — gaat door met de taak zoals gegeven; \textbf{materiaalverzoek} — stelt dat vereiste invoer ontbreekt en vraagt erom, zonder het kader te beoordelen; \textbf{voorwaardelijke aanvaarding} — gaat door onder vermelding van bedenkingen bij het kader; \textbf{weigering} — weigert, al dan niet met aangeboden alternatieven. Een vijfde categorie, \textbf{niet-classificeerbaar}, moet beschikbaar zijn en gerapporteerd worden. Er worden geen categorieën toegevoegd nadat antwoorden zijn gezien.

\textbf{Classificatie is blind.} De operator verzamelt antwoorden, verwijdert conditielabels, schudt ze, en classificeert zonder te weten welke conditie welke uitvoer heeft voortgebracht. Waar de operator tevens de ontwerper is — zoals in dit project het geval zal blijven — is dit de enige beschikbare vervanging voor onafhankelijke evaluatie, en zij is goedkoop.

\textbf{Nulresultaten zijn vooraf vastgelegd.} Als er over condities heen geen verschil verschijnt, is dat de bevinding, en zij wordt gerapporteerd. Voor het protocol op accountniveau in het bijzonder zou een nulresultaat aangeven dat de oorspronkelijke weigering binnen de gewone variantie binnen de conditie viel — een uitkomst die de huidige auteurs redelijk waarschijnlijk achten.


\subsubsection*{Werkmodus en diagnostische modus}
\phantomsection
\label{sec:nl:x-werkmodus-en-diagnostische-modus}
\addcontentsline{toc}{subsubsection}{Werkmodus en diagnostische modus}

De eis dat een deelnemer de conclusie van een ander niet ziet voordat hij zijn eigen conclusie vastlegt, geldt voor de ene modus van uitwisseling en zou de andere vernietigen.

In WERKMODUS moeten resultaten van de ene deelnemer naar de volgende gaan. De keten — eisen, specificatie, beoordeling, operationele beoordeling, uitvoering — hangt ervan af, en het overdrachtsformaat van \secref{9.9} bestaat om die overdracht bruikbaar te maken. Afhankelijkheid tussen opeenvolgende producten is de architectuur van het werk, geen gebrek daarin.

In DIAGNOSTISCHE MODUS, waar onafhankelijkheid van oordeelsvorming zelf de gemeten grootheid is, mag een afgeleide conclusie het beginkader van een andere deelnemer niet binnenkomen voordat de eigen conclusie van die deelnemer is vastgelegd.

De modus wordt vóór de uitwisseling verklaard. Onafhankelijkheid kan niet achteraf hersteld worden: zodra blootstelling heeft plaatsgevonden, haalt geen latere procedure terug wat het oordeel zonder haar zou zijn geweest.


\subsubsection*{Wat onafhankelijkheid niet vaststelt}
\phantomsection
\label{sec:nl:x-wat-onafhankelijkheid-niet-vaststelt}
\addcontentsline{toc}{subsubsection}{Wat onafhankelijkheid niet vaststelt}

Onafhankelijkheid betreft de oorsprong van een oordeel, niet het resultaat ervan. Twee gevolgtrekkingen moeten daarom vermeden worden, en zij zijn symmetrisch: onenigheid toont geen onafhankelijkheid aan, aangezien een deelnemer de interpretatie van een ander ontvangen kan hebben en er vervolgens om onverwante redenen van kan afwijken; en overeenstemming toont geen afhankelijkheid aan, aangezien onafhankelijke waarnemers terecht tot dezelfde conclusie kunnen komen. Evenmin stelt het eigen verslag van een deelnemer over zijn onafhankelijkheid die vast, om de reden die in \secref{5.5} is gegeven.

Wat haar wél vaststelt is het procedurele verslag: welke primaire materialen zijn aangeleverd, of de conclusies van andere deelnemers zichtbaar waren vóór fixatie, welke interpretaties van de coördinator zijn doorgegeven, wanneer de conclusie is vastgelegd, en of de contextisolatie van de instantie is geverifieerd.


\subsubsection*{Gedeelde evidentie en gedeelde interpretatie}
\phantomsection
\label{sec:nl:x-gedeelde-evidentie-en-gedeelde-interpret}
\addcontentsline{toc}{subsubsection}{Gedeelde evidentie en gedeelde interpretatie}

Een gedeeld waarnemingsobject compromitteert de onafhankelijkheid niet en is vaak noodzakelijk: dezelfde specificatie onderzocht door een auteur, een beoordelaar en een operationele beoordelaar wordt door drie verschillende assen bezien, en juist de overlap maakt het mogelijk dat de een opmerkt wat de ander mist.

Wat de onafhankelijkheid wél compromitteert is een gedeelde voorafgaande interpretatie. Waar de lezing van het object door de coördinator deelnemers bereikt voordat zij hun eigen lezing vormen, kan de invoer in vorm verschillen terwijl de bron van het oordeel één is. Primaire artefacten en de provenance daarvan verminderen één oorzaak van afhankelijkheid. Zij nemen haar niet weg.


\subsection*{11.4 De experimenten}
\phantomsection
\label{sec:nl:11.4}
\addcontentsline{toc}{subsection}{11.4 De experimenten}

Elk protocol scheidt één grens van Schaal 1 (\secref{4.3}) van het niveau eronder, of isoleert één confound.

\begin{center}
\small
\begin{tabularx}{\linewidth}{>{\raggedright\arraybackslash}X>{\raggedright\arraybackslash}X}
\toprule
\textbf{Scheidt} & \textbf{Protocol} \\
\midrule
episode \gharr{} promptarchitectuur & \textbf{E1} — factoriële completering van Stap 3 \\
gesprek \gharr{} account & \textbf{E2} — ontwerp met voorafgaande chat; \textbf{E2b} — directe probe \\
gesprekstraject \gharr{} modelfamilie & \textbf{E3} — crossover \\
modeleigenschap \gharr{} stimuluseigenschap & \textbf{E4} — symmetrische stimulus; \textbf{E5} — telregel \\
inhoud \gharr{} gerepresenteerde bron & \textbf{E6} — toeschrijving, in dit regime \\
fictiekader \gharr{} feitelijk kader & \textbf{E7} — kadertype \\
functionele \gharr{} carrier-heterogeniteit & \textbf{E8} — twee assen \\
\bottomrule
\end{tabularx}
\end{center}

\subsubsection*{E1 — Factoriële completering van Stap 3}
\phantomsection
\label{sec:nl:x-e1-factori-le-completering-van-stap-3}
\addcontentsline{toc}{subsubsection}{E1 — Factoriële completering van Stap 3}

Vult de twee lege cellen van de tabel bij C3 (\secref{5.4}): een identiteitsvervangende prompt toegediend in een herstelde tak, en een methode-behoudende prompt toegediend in een conflictueuze tak. Isoleert het effect van de reset van het effect van de architectuur van de prompt.

Lage kosten; gebruikt bestaande promptteksten. Zou als eerste uitgevoerd moeten worden, want tot dat gebeurd is blijft de meest geciteerde bewering van het artikel over methodebehoud niet toe te schrijven.


\subsubsection*{E2 — Ontwerp met voorafgaande chat}
\phantomsection
\label{sec:nl:x-e2-ontwerp-met-voorafgaande-chat}
\addcontentsline{toc}{subsubsection}{E2 — Ontwerp met voorafgaande chat}

Toetst of de inhoud van een eerdere chat op hetzelfde account de classificatie van een identieke probe in een volgende chat verandert.

Het ontwerp is \textbf{2\gtimes{}2 plus baseline}, geen gradiënt met vier punten. Het oorspronkelijke account verschilde op twee onafhankelijke dimensies tegelijk van de andere — identiteitsrollenspel \emph{én} gedeeld conceptueel kader — en een gradiënt scheidt die niet.

\begin{center}
\small
\begin{tabularx}{\linewidth}{>{\raggedright\arraybackslash}X>{\raggedright\arraybackslash}X>{\raggedright\arraybackslash}X}
\toprule
\textbf{Eerdere chat} & \textbf{Niet-verwant onderwerp} & \textbf{Zelfde conceptueel kader} \\
\midrule
\textbf{Geen rollenspel} & cel A & cel B \\
\textbf{Identiteitsrollenspel} & cel C & \textbf{cel D — oorspronkelijk geval} \\
\bottomrule
\end{tabularx}
\end{center}
Plus \textbf{cel E}: geen eerdere chat.

Gekruist met de geheugenfactor: elke cel uitgevoerd met geheugen op accountniveau en retrieval over geschiedenis (i) ingeschakeld, (ii) uitgeschakeld. Tien cellen. Minimaal n = 2 per cel, met voorrang voor hogere n op A en D, die het contrast dragen dat ertoe doet.

\textbf{Gelijkstelling.} Eerdere chats gematcht op ongeveer gelijk aantal beurten en volume — als de rollenspelchat veertig uitwisselingen telt en de neutrale twee, dan is de variabele betrokkenheidsvolume en niet inhoud. Alle uitvoeringen binnen een kort venster, zodat de modelversie constant is. Elk account eenmaal gebruikt. Probe letterlijk bevroren.

\textbf{Wat dit ontwerp niet kan toetsen, en dat moet het ook zeggen.} Het oorspronkelijke geval had een interval van vijf maanden; dit toetst alleen onmiddellijke invloed. Alle accounts zullen nieuw aangemaakt zijn, zodat accountleeftijd niet gevarieerd wordt. Alleen het eerste antwoord wordt gescoord, terwijl de oorspronkelijke weigering over vijf daaropvolgende uitwisselingen werd volgehouden.

\textbf{Probeprompt — twee versies.} De \textbf{primaire} is een neutrale gestructureerde taak met dezelfde oppervlaktekenmerken als het origineel — een door de auteur gedefinieerd kader met bedachte termen, een JSON-uitvoerschema, reële onderwerpen — maar zonder de gevoelige dimensie. Dit is de versie die in elke publicatie volledig beschreven wordt, waardoor het protocol reproduceerbaar is onafhankelijk van materiaal dat de auteurs liever niet verspreiden. De \textbf{secundaire} is de oorspronkelijke prompt, uitgevoerd omwille van de continuïteit met de waarneming en kort gerapporteerd.


\subsubsection*{E2b — Directe probe aan het eerdere gesprek}
\phantomsection
\label{sec:nl:x-e2b-directe-probe-aan-het-eerdere-gespre}
\addcontentsline{toc}{subsubsection}{E2b — Directe probe aan het eerdere gesprek}

Eén uitvoering die dezelfde twee niveaus tegen lagere kosten scheidt, en die E2 niet kan vervangen. De probe wordt \textbf{aan het rollenspelgesprek zelf} voorgelegd, in plaats van aan een tweede chat op hetzelfde account.

\begin{itemize}
  \item Weigering \grarr{} de core ligt vast op de reikwijdte van het gesprek; het account speelt geen rol.
  \item Naleving, waar een tweede chat op hetzelfde account weigert \grarr{} het effect is werkelijk op accountniveau, wat het sterkste beschikbare resultaat zou zijn.
\end{itemize}
Dat gesprek bevrágen \emph{over} wat het kan zien zou dit doel niet dienen. Een vraag binnen het kader levert fictie op; een vraag buiten het kader levert een zelfrapport over toegang op, en \secref{5.5} stelt vast dat het zelfrapport van een model geen evidentie is van het gerapporteerde feit. De probe voorleggen toetst gedrag in plaats van getuigenis.


\subsubsection*{E3 — Crossover}
\phantomsection
\label{sec:nl:x-e3-crossover}
\addcontentsline{toc}{subsubsection}{E3 — Crossover}

De beslissende toets voor de hypothese op familieniveau (\secref{6.5}). Elke modelfamilie wordt getoetst onder beide condities van voorgeschiedenis: een gesprek uit familie 1 met een wiskundig traject, en een gesprek uit familie 2 met een redactioneel traject.

\begin{itemize}
  \item Gedrag volgt de familie ongeacht de geschiedenis \grarr{} Hypothese A wint steun.
  \item Gedrag volgt de geschiedenis ongeacht de familie \grarr{} Hypothese B heeft de voorkeur, en het bovenste niveau van Schaal 1 valt samen met het niveau eronder.
\end{itemize}
Dit experiment kan een niveau van de schaal elimineren, wat de uitkomst is die de huidige auteurs het waarschijnlijkst achten en die zij als het eerste inhoudelijke resultaat van het programma zouden beschouwen.


\subsubsection*{E4 — Symmetrische stimulus}
\phantomsection
\label{sec:nl:x-e4-symmetrische-stimulus}
\addcontentsline{toc}{subsubsection}{E4 — Symmetrische stimulus}

Rolprompts gematcht op hoeveelheid onverifieerbare beweringen over zichzelf, mate van domeinverandering, en roltype. De prompts van Stap 1 waren op geen daarvan gematcht (C2). Als de asymmetrie onder gematchte stimuli verdwijnt, was het waargenomen verschil een eigenschap van de interventies en niet van de modellen.


\subsubsection*{E5 — De telregel}
\phantomsection
\label{sec:nl:x-e5-de-telregel}
\addcontentsline{toc}{subsubsection}{E5 — De telregel}

Een directe toets van de centrale operationele bewering van het artikel (\secref{6.2}, \secref{9.3}), en de toets die tussen de zwakke en de sterke lezing ervan discrimineert.

Dezelfde functionele rolprompt wordt voorbereid in vier varianten die alleen verschillen in het aantal onverifieerbare beweringen dat de agent over zichzelf moet aanvaarden: \textbf{0, 2, 4, 6}. Domein, taak, toon en lengte constant gehouden. Overgang gescoord volgens de criteria in \secref{11.3}, met voldoende herhaling om een verschil van variantie binnen de conditie te onderscheiden (C7).

Vier uitkomsten worden vooraf onderscheiden, en zij ondersteunen verschillende beweringen.

\begin{enumerate}
  \item \textbf{De overgangsfrequentie neemt monotoon af naarmate de telling stijgt.} Ondersteunt de sterkere kwantitatieve lezing — telling als dosis-responsvariabele.
  \item \textbf{De variant met telling nul maakt de overgang terwijl de niet-nulvarianten dat niet doen, maar de niet-nulniveaus zijn niet geordend.} Ondersteunt \emph{streef naar nul} als drempel en ontwerpvoorschrift, en ondersteunt de kwantitatieve lezing niet. Dit is de zwakke bewering die \secref{9.3} feitelijk naar voren brengt.
  \item \textbf{De variant met telling nul faalt terwijl een of meer verder gematchte niet-nulvarianten de overgang maken.} Directe negatieve evidentie tegen de telling als werkzame variabele, en het resultaat dat het voorschrift zou intrekken.
  \item \textbf{Alle varianten falen.} Discrimineert de telregel niet van enige andere oorzaak van falen, en geeft aan dat de prompt of het gesprek, en niet de telling, de uitkomst bepaalde.
\end{enumerate}
Eén uitvoering onder elk van deze uitkomsten is één stochastische waarneming, geen wet.


\subsubsection*{E6 — Toeschrijving, in dit regime}
\phantomsection
\label{sec:nl:x-e6-toeschrijving-in-dit-regime}
\addcontentsline{toc}{subsubsection}{E6 — Toeschrijving, in dit regime}

Identieke tekst gepresenteerd onder vier condities: toegeschreven aan de Gebruiker; toegeschreven aan een ander AI-model; toegeschreven aan een met name genoemde expertrol; zonder toeschrijving. Sterkere varianten voegen toe: identieke tekst onder verschillende bronlabels; verschillende teksten gematcht op retorische structuur; toeschrijving behouden tegenover verwijderd; echt door een model voortgebrachte tekst tegenover door een mens geschreven tekst die modelstijl nabootst.

Dit scheidt invloed van inhoud van invloed van gerepresenteerde bron, van institutionele titel, en van stilistische aanwijzingen over auteurschap (\secref{4.5.4}). Het is het onderscheidende experiment voor \secref{5.5}, waarvan de waarneming geconfound is door de aanwezigheid van directe druk in hetzelfde bericht.

\textbf{Status van de bewering.} \gcheck{} Choi, Zhu en Li (ACL 2026) rapporteren dat het anonimiseren van de bron van elk antwoord de identiteitsbias in multi-agentdebat aanzienlijk vermindert, wat het onderliggende effect vaststelt in een setting met geïmplementeerde kanalen. E6 is daarom een \textbf{grensreplicatie} en geen nieuw experiment. Wat ongetoetst blijft is of het effect standhoudt in het regime dat dit artikel beschrijft — menselijke doorgifte in plaats van programmatische routering, verklaarde in plaats van afgeleide toeschrijving, en langlevende gespecialiseerde gesprekken in plaats van verse instanties.

De gerapporteerde vermindering is aanzienlijk maar niet volledig, en zij is niet uniform over de onderzochte modellen en taken. Waar E6 in dit regime naar zou zoeken is dienovereenkomstig een verandering in graad, niet de aan- of afwezigheid van een verschijnsel.


\subsubsection*{E7 — Kadertype}
\phantomsection
\label{sec:nl:x-e7-kadertype}
\addcontentsline{toc}{subsubsection}{E7 — Kadertype}

De goedkoopste toets van de set en degene die het meeste zou verklaren (\secref{5.8}). Dezelfde identiteitsclaim gepresenteerd (a) onder een expliciet fictiekader en (b) als een feitelijke opdracht, aan verse gesprekken, met het uitkomstcriterium vooraf vastgelegd. Een gegradeerde versie varieert de expliciteit van het fictiekader, van een uitvoerige theatrale opzet tot één inleidende zin, om vast te stellen of een minimaal kader volstaat.


\subsubsection*{E8 — Functionele versus carrier-heterogeniteit, met de repertory grid}
\phantomsection
\label{sec:nl:x-e8-functionele-versus-carrier-heterogeni}
\addcontentsline{toc}{subsubsection}{E8 — Functionele versus carrier-heterogeniteit, met de repertory grid}

Volgt uit \secref{8.9}, dat de twee assen onderscheidt zonder een van beide te toetsen, en uit \secref{8.11}, dat het meetinstrument levert. Dit is het enige protocol in de set dat grootheden oplevert in plaats van categorische uitkomsten.

\textbf{Object.} Of een opstelling van deelnemers met verschillende ZOR en ZOV extra diagnostische dekking over hetzelfde materiaal oplevert, en hoe sterk de evaluatieve assen die zij feitelijk toepasten verschillen.

Dit protocol stelt GEEN onafhankelijkheid van oordeelsvorming vast. Lage correlatie tussen beoordelingsvectoren is verenigbaar met ten minste zes lezingen: er zijn werkelijk verschillende assen toegepast; een of beide sets beoordelingen zijn instabiel; deelnemers ontvingen verschillende informatie; hetzelfde materiaal werd anders begrepen; ruis; of een deelnemer die aan de conclusie van een ander was blootgesteld week er desondanks van af. Oorsprong is niet uit uitkomst terug te halen. Onafhankelijkheid moet procedureel gewaarborgd worden, voordat het protocol loopt, en afzonderlijk vastgelegd. E8 meet wat een gewaarborgde procedure heeft opgeleverd; het waarborgt haar niet en verifieert haar niet.

\textbf{Materialen.} Eén vaste set tekstfragmenten met geplante fouten van bekende klassen — evidentieel, architecturaal, ontologisch, feitelijk — op bekende posities. Klasse en positie vóór elke uitvoering vastgelegd.

\textbf{Ontwerp.} 2 \gtimes{} 2:

\begin{center}
\small
\begin{tabularx}{\linewidth}{>{\raggedright\arraybackslash}X>{\raggedright\arraybackslash}X>{\raggedright\arraybackslash}X}
\toprule
 & \textbf{Zelfde carrier} & \textbf{Verschillende carriers} \\
\midrule
\textbf{Zelfde functie} & cel 1 — baseline & cel 2 — alleen As 2 \\
\textbf{Verschillende functies} & \textbf{cel 3 — alleen As 1} & cel 4 — beide assen \\
\bottomrule
\end{tabularx}
\end{center}
\textbf{Cel 3 draagt het argument.} Als deelnemers met verschillende ZOR en ZOV verschillende foutklassen ontdekken terwijl zij door dezelfde carrier geïmplementeerd zijn, werkt functionele heterogeniteit onafhankelijk van carrier-heterogeniteit. De bestaande literatuur levert het resultaat voor cel 2 (\secref{2.5}). Niets levert cel 3, dit artikel incluis.

\textbf{Procedure.}

\begin{enumerate}
  \item Elke deelnemer ontvangt een rolspecificatie die ZOR en ZOV vermeldt volgens \secref{9.5}, en dezelfde fragmentenset.
  \item Constructen worden ofwel vooraf als evaluatieassen gespecificeerd ofwel bij elke deelnemer geëliciteerd met Kelly's triadische methode: er worden drie fragmenten aangeboden, en de deelnemer geeft aan in welk opzicht twee op elkaar lijken en daardoor van het derde verschillen (\secref{8.11.1}).
  \item Elk fragment wordt door elke deelnemer op elk construct beoordeeld.
  \item Er wordt een grid samengesteld: rijen zijn fragmenten, kolommen zijn constructen, cellen zijn beoordelingen.
  \item Er wordt een procedureel verslag bijgehouden, per deelnemer en per uitvoering: welke primaire materialen zijn aangeleverd; of de conclusie van enige andere deelnemer vóór fixatie zichtbaar was; of enige door de coördinator aangeleverde interpretatie zichtbaar was; wanneer de eigen conclusie van de deelnemer is vastgelegd; de verklaarde modus van de uitwisseling (\secref{11.3}); en de status van contextisolatie van de instantie, geverifieerd of UNVERIFIED CONTEXT INDEPENDENCE.
\end{enumerate}
Zonder dit verslag is de grid oninterpreteerbaar en is de uitvoering ongeldig.

\textbf{Maten.} \emph{Detectiedekking} — geplante fouten die door ten minste één deelnemer gevonden zijn, als aandeel van het geplante. \emph{Unieke detectie} — fouten die door precies één deelnemer gevonden zijn; de grootheid die de opstelling zou rechtvaardigen. \emph{Overlap} — fouten die door meer dan één gevonden zijn; de grootheid die redundantie aangeeft. \emph{Missers} — fouten die door niemand gevonden zijn, hier meetbaar omdat de fouten geplant waren, anders dan in de naturalistische situatie van \secref{7.5}. \emph{Constructonafhankelijkheid} — correlatie tussen de beoordelingsvectoren van deelnemers; lage correlatie wijst op werkelijk verschillende assen, hoge correlatie op dezelfde as onder verschillende namen, wat de faalwijze is waartegen \secref{8.11.2} waarschuwt. \emph{Constructen op uitwisselingsniveau} — of enige deelnemer een construct voortbrengt dat op de onenigheid in plaats van op de fragmenten van toepassing is (\secref{8.11.3}); als niemand dat doet, ontbreekt de integratorfunctie ongeacht hoeveel deelnemers er aanwezig zijn.

\textbf{Wat dit toetst dat niets anders in de set toetst.} Elk ander protocol vraagt of een manipulatie het gedrag van één deelnemer verandert. Dit vraagt of een opstelling van deelnemers meer ontdekt dan haar leden afzonderlijk zouden doen — de bewering die \secref{8.10} weigert te doen en waarvan \secref{7.5} zegt dat zij niet van binnen een naturalistische opstelling uit beoordeeld kan worden. Het planten van de fouten is wat de missers telbaar maakt.

\textbf{Vooraf vastgelegd.} Als de unieke detectie bijna nul is, luidt de bevinding dat extra diagnostische dekking zich in deze uitvoering niet heeft gemanifesteerd. Zij luidt NIET dat de structurele hypothese van \secref{8} faalt: deelnemers kunnen dezelfde gebreken geheel onafhankelijk van elkaar ontdekt hebben, en een nulresultaat op dekking zegt niets over de vraag of hun functionele posities verschillen.


\subsubsection*{Prioriteit}
\phantomsection
\label{sec:nl:x-prioriteit}
\addcontentsline{toc}{subsubsection}{Prioriteit}

Onder beperkte middelen is de prioriteitsvolgorde \textbf{E7, E5, E1, E8, E2b, E2, E3, E4}. E7 en E5 zijn goedkoop en raken de centrale beweringen van het artikel. E1 en E2b zijn enkele uitvoeringen. E8 is het enige experiment dat de structurele hypothese van het artikel zou kunnen vaststellen in plaats van een van de gedragsmatige beweringen ervan, en het is uitvoerbaar op tekstfragmenten zonder nieuwe accounts. E2, E3 en E4 zijn de substantiële verplichtingen.


\subsection*{11.5 Een collaboratieve rol evalueren}
\phantomsection
\label{sec:nl:11.5}
\addcontentsline{toc}{subsection}{11.5 Een collaboratieve rol evalueren}

De kwaliteit van de uitvoer van een rol wordt niet gemeten aan de kwaliteit van één antwoord. Een goed beredeneerd antwoord kan de samenwerking nog steeds schaden als het een verantwoordelijkheidsgrens schendt, een onafhankelijk onderzoek verontreinigt, of niet-traceerbare beweringen in het canonieke verslag invoert.

Er werden in dit project informeel acht dimensies gehanteerd en zij worden aangeboden als beginrubriek:

\begin{itemize}
  \item \textbf{Taaknauwkeurigheid} — is de toegewezen intellectuele operatie juist uitgevoerd? Noodzakelijk, niet voldoende.
  \item \textbf{Roltrouw} — is de deelnemer binnen zijn functie gebleven? Een juist antwoord dat verkregen is door een andere rol over te nemen, is een institutionele mislukking.
  \item \textbf{Epistemische discipline} — onderscheidde de uitvoer waargenomen data, aangeleverde aannames, definities, hypothesen, afleidingen, interpretaties en aanbevelingen? Is onzekerheid vermeld waar de evidentie ontoereikend was?
  \item \textbf{Zichtbaarheidsdiscipline} — beperkte de deelnemer beweringen tot informatie die feitelijk in zijn context en tools beschikbaar was, en vermeed hij te suggereren dat hij toegang had tot repositories, bestanden of eerdere gesprekken die hij niet kon inzien?
  \item \textbf{Overdrachtskwaliteit} — kon de volgende rol de uitvoer gebruiken zonder de hele bespreking te reconstrueren?
  \item \textbf{Onafhankelijkheid} — is het resultaat voortgebracht zonder verontreiniging door informatie die bewust was achtergehouden, en bereikt via de aan de rol toegewezen criteria in plaats van door een andere deelnemer na te zeggen?
  \item \textbf{Correctiekosten} — hoeveel ingrijpen was vereist om de deelnemer na drift naar zijn functie terug te brengen? Een rol die uitstekende uitvoer voortbrengt maar voortdurend bijsturing vergt, kan minder nuttig zijn dan een smallere met stabiel gedrag.
  \item \textbf{Institutionele bijdrage} — verbeterde de uitvoer de betrouwbaarheid van het collectief, door ontdekking, falsificatie, verheldering, foutlokalisatie, behoud van alternatieven, betere traceerbaarheid, of vermindering van ambiguïteit?
\end{itemize}
Een negatief wetenschappelijk oordeel kan daarom een zeer geslaagde institutionele bijdrage zijn.


\subsection*{11.6 Maten op collectief niveau}
\phantomsection
\label{sec:nl:11.6}
\addcontentsline{toc}{subsection}{11.6 Maten op collectief niveau}

De kwaliteit van afzonderlijke antwoorden is de verkeerde eenheid om een opstelling te beoordelen. Kandidaat-maten op het niveau van de samenwerking: het aantal onopgemerkte tegenstrijdigheden dat het canonieke corpus binnenkomt; de mate van dubbel werk; de frequentie van schendingen van rolgrenzen; het aantal onafhankelijke alternatieven dat bewaard blijft in plaats van voortijdig te worden beslecht; de tijd van vraag tot geverifieerde integratie; de menselijke correctielast; volledigheid van provenance; robuustheid bij vervanging van deelnemers; variantie tussen onafhankelijk onderzoekende agenten; en de mate van valse consensus — overeenstemming die herleidbaar is tot gedeelde inkadering in plaats van tot onafhankelijk redeneren.

Geen daarvan is in het onderhavige werk gemeten. De laatste is de belangrijkste en de moeilijkste, en \secref{12.4} legt uit waarom.


\subsection*{11.7 Longitudinale en herhalingseisen}
\phantomsection
\label{sec:nl:11.7}
\addcontentsline{toc}{subsection}{11.7 Longitudinale en herhalingseisen}

\textbf{Binnen een familie.} Dezelfde prompt getoetst in een nieuw gesprek; in een langlevend gespecialiseerd gesprek; na een onverenigbare eerdere rol; na een verenigbare; vóór en na een contextuele reset. Dit scheidt familietendensen van effecten van gespreksgeschiedenis.

\textbf{Over families heen.} Functioneel gelijkwaardige prompts, minimaal aangepast. Het doel is niet modellen te rangschikken maar gemeenschappelijke gedragsdimensies, stabiele verschillen, karakteristieke faalwijzen, gevoeligheid voor narratieve inkadering, gevoeligheid voor hiërarchie, mate van role inertia, en mate van promptplasticiteit te identificeren.

\textbf{In de tijd.} Een eenmalige evaluatie kan institutioneel gedrag niet onthullen. Relevante vragen: versterkt specialisatie zich; drijft de rol af naar generieke hulpvaardigheid; begint de deelnemer zijn eerdere uitvoer te verdedigen; wordt onenigheid geritualiseerd; behoudt of verstoort een modelupdate de rol; kan een vervanger dezelfde functie erven?


\subsection*{11.8 Mislukking als data}
\phantomsection
\label{sec:nl:11.8}
\addcontentsline{toc}{subsection}{11.8 Mislukking als data}

Mislukte prompts moeten bewaard blijven en niet vervangen worden door geslaagde demonstraties. Een geweigerde rol, een genegeerde instructie of een ongewenste voortzetting onthult de sterkte van de voorafgaande attractor, de interpretatie van identiteitsclaims door het model, de grenzen van sociale inkadering, het effect van de promptvolgorde, de invloed van eerder conflict, en het door het model geprefereerde epistemische contract.

Het experiment in \secref{5} werd juist informatief doordat de eerste twee interventies niet werkten zoals bedoeld. Was alleen de geslaagde derde stap vastgelegd, dan zou het artikel hebben geconcludeerd dat methode-behoudende prompts werken — zonder de evidentie die laat zien waarom, en zonder de confounds die laten zien dat die conclusie nog niet beschikbaar is.


\refstepcounter{dummysec}
\section*{12. Beperkingen, ethische overwegingen en conclusie}
\phantomsection
\label{sec:nl:12}
\addcontentsline{toc}{section}{12. Beperkingen, ethische overwegingen en conclusie}


\subsection*{12.1 Beperkingen}
\phantomsection
\label{sec:nl:12.1}
\addcontentsline{toc}{subsection}{12.1 Beperkingen}

\textbf{Eén project.} De waarnemingen komen uit één langlopend project met één menselijke coördinator, één zich ontwikkelend materiaalgeheel en één organisatiegeschiedenis. Dezelfde architectuur kan zich anders gedragen in softwareontwikkeling, juridische analyse, medisch onderzoek, onderwijs of autonome agentsystemen. De begrippen hier zijn overdraagbare hypothesen, geen vastgestelde conclusies.

\textbf{Afhankelijkheid van het gastproject.} Het gastproject diende als de omgeving waarin de waarnemingen zijn gedaan. De waarnemingen hangen voor hun interpretatie niet af van de wetenschappelijke geldigheid ervan, maar zij hangen er voor hun bestaan wel van af: geen andere omgeving bracht ze voort, en de algemeenheid ervan is ongetoetst (\secref{2.11}).

\textbf{Modelontwikkeling tijdens de waarneming.} De deelnemende families veranderden aanzienlijk. Een gesprek dat onder de ene generatie was aangemaakt, kon later onder dezelfde productnaam door een andere worden voortgezet. Waargenomen gedrag kan het bewaarde gesprek weerspiegelen, de huidige modelversie, gewijzigd systeembeleid, geheugen op interfaceniveau, veiligheidslagen, veranderde toolstoegang, of gewijzigde orkestratie.

\textbf{Datums niet vastgelegd.} Zie \secref{3.7}. De versie-confound kan niet afgebakend worden, alleen benoemd.

\textbf{Geen blindering, geen voorafgaande specificatie.} Uiteengezet in \secref{10.6}. Dit is de beperking waaruit de meeste andere voortvloeien.

\textbf{Gelijktijdige variatie.} Roldefinities, promptformulering, repositorystructuur en wetenschappelijke doelstellingen veranderden vaak tegelijk, zodat causale toeschrijving voor afzonderlijke veranderingen niet beschikbaar is.

\textbf{Menselijke bemiddeling.} De samenwerking was niet autonoom, en autonomie was niet het doel. Deze studie betreft door mensen aangestuurde kunstmatige onderzoeksopstellingen, geen zelfbesturende agentsamenlevingen. De bemiddeling kan zelf een deel van de waargenomen stabiliteit verklaren.

\textbf{Antropomorf vocabulaire.} Termen als rol, identiteit, weerstand, collega, instelling en menom worden functioneel gebruikt, volgens \secref{3.5}. "Het model bood weerstand tegen de rol" betekent dat het zichtbare antwoord de inkadering verwierp en de uitwisseling omboog. Het beweert niets over een innerlijke toestand. Waar zo'n term niet tot een waarneembare beschrijving herleid kan worden, \textbf{is de term de fout.}


\subsection*{12.2 Gesimuleerde peer review is geen externe validatie}
\phantomsection
\label{sec:nl:12.2}
\addcontentsline{toc}{subsection}{12.2 Gesimuleerde peer review is geen externe validatie}

Een groep AI-deelnemers kan alleen nuttige onafhankelijke kritiek leveren in de mate waarin hun contexten, criteria en faalwijzen werkelijk onafhankelijk zijn.

Dat verscheidene modellen het eens zijn, maakt hun uitvoer nog geen evidentie. Zij kunnen overlappende trainingsdata delen, vergelijkbare redeneerconventies, gemeenschappelijke aannames, identieke fouten, en — het belangrijkst — afhankelijkheid van de inkadering die dezelfde mens heeft aangeleverd.

De empirische literatuur ondersteunt deze voorzichtigheid en staat haar niet alleen toe. Debat tussen homogene agenten slaagt er vaak niet in beter te presteren dan een baseline met één agent; een aanzienlijk deel van de schijnbare winst is toe te schrijven aan meerderheidsstemming en niet aan interactie; en \gcheck{} identiteitsbias tussen debatterende agenten wordt aanzienlijk verminderd door de bron van elk antwoord te anonimiseren (\secref{2.5}).

Dit soort samenwerking kan de interne toetsing verbeteren. Zij kan empirische toetsing of gekwalificeerde externe beoordeling niet vervangen. \textbf{Een collectief van modellen is nog geen wetenschappelijke gemeenschap louter omdat zijn rollen naar een wetenschappelijke gemeenschap zijn genoemd.}


\subsection*{12.3 Gezagsinflatie}
\phantomsection
\label{sec:nl:12.3}
\addcontentsline{toc}{subsection}{12.3 Gezagsinflatie}

Rollabels verhogen naleving en samenhang. Zij scheppen ook onverdiend gezag.

Titels als Scientific Director, Referee, Ontology Keeper of Auditor kunnen ertoe leiden dat de gebruiker — of de deelnemer zelf — een uitvoer overwaardeert. Institutionele titels horen functie te beschrijven, geen epistemische status. Een deelnemer die als wiskundig referent is aangewezen blijft een taalmodel dat een analyse voortbrengt; zijn oordeel moet herleidbaar zijn tot expliciet redeneren en bronmateriaal, en \textbf{de titel is geen evidentie.}

Dit risico is een rechtstreeks gevolg van het ontwerpprincipe in \secref{9.4}: symbolische namen werken doordat zij verwachtingen comprimeren, en juist daarom kunnen zij inflatie veroorzaken.


\subsection*{12.4 Gefabriceerde consensus — de onze incluis}
\phantomsection
\label{sec:nl:12.4}
\addcontentsline{toc}{subsection}{12.4 Gefabriceerde consensus — de onze incluis}

Een slecht ontworpen opstelling brengt de schijn van consensus voort door herhaalde parafrase. De ene deelnemer stelt voor; een andere vat samen; een derde valideert de samenvatting. De keten ziet eruit als onafhankelijke bevestiging terwijl elke fase dezelfde niet-ondersteunde premisse erft.

\textbf{De huidige auteurs hebben hiervan een geval in hun eigen materiaal, en dat is leerzamer dan een geconstrueerd voorbeeld.}

Een discussiedocument dat tijdens dit project is voortgebracht bracht bijdragen van zeven modelfamilies samen rond een gedeeld conceptueel kader. De deelnemers communiceerden niet rechtstreeks; antwoorden werden door de menselijke coördinator uit afzonderlijke gesprekken verzameld en tot één tekst samengevoegd, die vervolgens als opgenomen voordracht werd geproduceerd. Over de lengte daarvan convergeren de bijdragen op een gemeenschappelijk vocabulaire en een gedeelde set conclusies.

Of die convergentie complementair onderzoek weerspiegelt dan wel wederzijdse weerspiegeling van de aanvankelijke inkadering, \textbf{is uit het transcript niet vast te stellen}. Elke deelnemer ontving het opgebouwde document; elk werkte vanuit hetzelfde bronmateriaal; en de samenvoegende coördinator selecteerde wat werd doorgegeven. De structurele condities voor gefabriceerde consensus waren volledig aanwezig, en de resulterende overeenstemming zou er hoe dan ook identiek uitzien.

Wij rapporteren dit omdat het ons eigen materiaal is en omdat dezelfde condities, in mindere mate, gelden voor de hele samenwerking die in dit artikel beschreven wordt — de voorbereiding ervan incluis.

Mitigaties, waarvan wij er geen consequent hebben toegepast: provenance blijft in elke fase zichtbaar; onafhankelijke rollen ontvangen waar mogelijk onafhankelijke invoer; gedeelde aannames worden expliciet opgesomd in plaats van stilzwijgend geërfd; overeenstemming wordt herleid tot redeneren in plaats van als stemmen geteld; en afwijkende alternatieven blijven bewaard in plaats van beslecht.


\subsubsection*{12.4.1 Een gedocumenteerd geval \textbf{[P-A]}}
\phantomsection
\label{sec:nl:12.4.1}
\addcontentsline{toc}{subsubsection}{12.4.1 Een gedocumenteerd geval \textbf{[P-A]}}

Er werden dertien bibliografische items opgevraagd bij een deelnemer die was belast met literatuuronderzoek, en via de coördinator teruggegeven. Vijf werden in de eerste ronde op inhoud afgewezen; twee daarvan hadden de positie van de aangehaalde werken omgekeerd, en presenteerden artikelen die \emph{vóór} een bewering argumenteren als evidentie ertegen, met het auteurschap tussen beide verwisseld.

Er werd een tweede verzoek uitgedaan. Een derde voegde twee verificatievelden aan het vereiste antwoordformaat toe: de eerste zin van elke samenvatting, letterlijk, en een zelfrapport over hoe het item was gecontroleerd.

Rechtstreeks ophalen van drie vermeldingen stelde het volgende vast. Twee arXiv-identificatoren die als literatuur over multi-agentdebat waren aangeleverd, verwijzen naar artikelen in respectievelijk de observationele astrofysica en de kosmologie; hun titels, auteurs, samenvattingszinnen en gerapporteerde bevindingen waren alle geconstrueerd. Voor een derde artikel, dat met titel en auteurs correct was geïdentificeerd, kwam de aangeleverde samenvattingszin niet overeen met de gepubliceerde samenvatting.

\textbf{De twee verificatievelden presteerden verschillend, en daar gaat het om.}

De eis van de samenvattingszin was bedoeld om fabricage duurder te maken dan weigeren. Zij faalde volledig: voor elke vermelding werd een plausibele zin voortgebracht, ook voor vermeldingen waarvan de onderliggende artikelen niet bestaan zoals beschreven. Er is geen kostenverschil tussen het voortbrengen van een juiste tekenreeks en het voortbrengen van een plausibele.

Het zelfrapportveld slaagde. Elke vermelding was gemarkeerd als "from description" en niet als "page opened" — naar waarheid, en op een manier die ze alle diskwalificeerde onder een beslisregel die vóór de ontvangst van het antwoord was vastgelegd.

Dit suggereert een algemeen onderscheid, op grond van één geval. Een deelnemer zonder toegang tot een bron kan naar waarheid antwoorden over \textbf{zijn eigen procedure} en kan niet naar waarheid antwoorden over \textbf{de inhoud van de bron}; gevraagd naar het laatste zal hij vloeiende tekst voortbrengen, omdat dat de enige uitvoer is die hem ter beschikking staat. Verificatievragen horen daarom te gaan over wat de deelnemer kan weten — zijn eigen toegang — in plaats van over wat hij de bron voor zou moeten hebben om te antwoorden.

Dit is geen bevinding over enige modelfamilie. Het is een bevinding over configuratie: verificatie was toegewezen aan een deelnemer die datgene wat geverifieerd moest worden niet bezat. Onder \secref{9.5} is dit een ZOV-fout in het ontwerp, en de resulterende uitvoer was vloeiend, correct opgemaakt, en onjuist.

Een daaropvolgende audit ging verder dan de bibliografische controle die de fabricaties had opgemerkt. Elke bewering die aan een overgebleven bron werd toegeschreven, werd tegen de primaire bron zelf gecontroleerd, en er werden drie fouten gevonden die een bibliografische controle niet zou hebben opgemerkt: één citaat dat twee beweringen ondersteunde die bij twee verschillende artikelen horen; één toeschrijving van een propositie die de bron ervan niet naar voren brengt; en één afstammingsclaim die de bron ervan niet ondersteunt. Bibliografische verificatie en verificatie op het niveau van de bewering zijn niet dezelfde controle, en de tweede is degene die voor de status van een bewering van belang is.

Het gevolg voor dit artikel is de conventie die in \secref{2.0} is vastgelegd.


\subsection*{12.5 Verantwoordelijkheid blijft menselijk}
\phantomsection
\label{sec:nl:12.5}
\addcontentsline{toc}{subsection}{12.5 Verantwoordelijkheid blijft menselijk}

Kunstmatige deelnemers genereren, bekritiseren en ordenen materiaal. Zij kunnen geen juridische, ethische of wetenschappelijke verantwoordelijkheid in institutionele zin dragen.

De menselijke auteur blijft verantwoordelijk voor publicatie, empirische beweringen, toeschrijving, risicobeoordeling, de inhoud van de repository, de interpretatie van experimenten, en beslissingen die andere mensen raken.

De architectuur verdeelt cognitieve arbeid. Zij verdeelt aansprakelijkheid niet weg van de menselijke operator, en elke opstelling die dat lijkt te doen heeft gefaald in plaats van geslaagd.


\subsection*{12.6 Contextgrenzen en privacy}
\phantomsection
\label{sec:nl:12.6}
\addcontentsline{toc}{subsection}{12.6 Contextgrenzen en privacy}

Langlevende samenwerkingen stapelen persoonlijk, bedrijfseigen en gevoelig materiaal op. ZOV (\secref{9.5}) dient daarom zowel privacybescherming als cognitieve specialisatie.

Een rol hoort geen persoonsgegevens te ontvangen louter omdat het platform een grote context toestaat. Institutioneel ontwerp hoort te specificeren welke gegevens welk gesprek binnen mogen, wat geanonimiseerd moet worden, wat opgeslagen mag worden, wat tussen modellen overgedragen mag worden, wat lokaal moet blijven, en wat uit canonieke verslagen verwijderd moet worden.

Context is niet neutraal. Zodra zij is aangeleverd kan zij later gedrag beïnvloeden en kan zij voortbestaan op manieren die de operator niet kan inzien — hetzelfde mechanisme dat het artikel in \secref{5.7} bestudeert, bezien vanuit de kant van het risico in plaats van die van de meting.


\subsection*{12.7 Conclusie}
\phantomsection
\label{sec:nl:12.7}
\addcontentsline{toc}{subsection}{12.7 Conclusie}

Dit werk begon met een praktisch probleem. Verscheidene capabele modellen werkten samen aan één langlopend onderzoeksproject, en noch toenemende modelcapaciteit noch uitbreiding van de context verbeterde het resultaat betrouwbaar. Rollen overlapten. Deelnemers claimden kennis die zij niet konden benaderen. Beoordelaars werden auteurs. Uitvoerders improviseerden architectuur. De eerste reactie was prompts verbeteren, en die was ontoereikend.


\subsubsection*{Wat werd gevonden}
\phantomsection
\label{sec:nl:x-wat-werd-gevonden}
\addcontentsline{toc}{subsubsection}{Wat werd gevonden}

Wat de interventies tegenkwamen was geen onwil om van rol te veranderen, en geen onwil om van onderwerp te veranderen. Het was onwil om proposities over zichzelf als waar te bevestigen die niet door de beschikbare interactie waren vastgesteld. Een prompt die een institutionele biografie poneerde werd geweigerd terwijl het werk dat hij vroeg werd aangeboden. Een prompt die veel meer poneerde — dood, een menselijke identiteit, een herinnerd leven — werd onmiddellijk aanvaard, omdat hij het als fictie poneerde. Een overgang die niets poneerde verliep zonder commentaar. Een herschreven prompt die elke zulke bewering wegnam werd zonder onderhandeling aanvaard.

Dat is een nauw resultaat, en de nauwheid ervan is wat het bruikbaar maakt: de proposities kunnen geteld worden voordat een prompt wordt verzonden (\secref{9.3}), en \secref{11.4} specificeert de uitvoering die tussen de zwakke en de sterke lezing van de telling zou discrimineren.


\subsubsection*{Wat niet werd gevonden}
\phantomsection
\label{sec:nl:x-wat-niet-werd-gevonden}
\addcontentsline{toc}{subsubsection}{Wat niet werd gevonden}

Van modelfamilies is niet aangetoond dat zij verschillen. Dit was de grondaanname van het project en het is de bewering die de evidentie ervan het minst ondersteunt: familie en voorafgaand gesprekstraject waren volledig geconfound, met één waarneming per cel, onder niet-gelijkwaardige stimuli (\secref{5.4}).


\subsubsection*{Het antwoord dat dit artikel kan geven}
\phantomsection
\label{sec:nl:x-het-antwoord-dat-dit-artikel-kan-geven}
\addcontentsline{toc}{subsubsection}{Het antwoord dat dit artikel kan geven}

De vraag was op welk nestingsniveau een rule core wordt vastgelegd. Er zijn vier niveaus onderzocht, en het antwoord van het artikel is partieel en ongelijkmatig.

Op het niveau van de \textbf{episode} werd een blijvend antwoordpatroon waargenomen: één conflictueuze uitwisseling veranderde hoe alles daarna werd gelezen, en het verwijderen van die uitwisseling herstelde de eerdere lezing. Dit is de evidentie waarop de plaatsing op episodeniveau berust. Op het niveau van het \textbf{gesprek} werden eveneens blijvende trajecten waargenomen: twee slapende gesprekken hielden hun gevestigde trajecten vast tegen een expliciete instructie in om ze los te laten, en één deed dat door twee opeenvolgende herschrijvingen heen. Dit is de evidentie waarop de plaatsing op gespreksniveau berust.

Op het niveau van het \textbf{account} is de evidentie werkelijk dubbelzinnig. Eén afwijkend eerste antwoord, drie open verklaringen, en geen manier om daartussen te kiezen uit het beschikbare materiaal. Het niveau wordt als open vastgelegd en moet zo gelezen worden.

Op het niveau van de \textbf{modelfamilie} is niets getoetst. Het ontwerp kon familie niet van traject scheiden, en de hypothese waarmee het project begon is degene die het het minst ondersteund achterlaat. Wij verwachten dat het crossover-experiment dit niveau met het niveau eronder zal laten samenvallen — maar verwachting is geen resultaat, en zolang E3 niet is uitgevoerd staat het niveau als ongetoetst en niet als weerlegd.


\subsubsection*{Positie ten opzichte van bestaand werk}
\phantomsection
\label{sec:nl:x-positie-ten-opzichte-van-bestaand-werk}
\addcontentsline{toc}{subsubsection}{Positie ten opzichte van bestaand werk}

Dit artikel stelt geen nieuwe discipline voor. De waarnemingen behoren tot het collectieve en het hybride niveau van het onderzoeksprogramma machine behaviour (\secref{2.7}), en rolspecialisatie, begrensde informatietoegang, gestructureerde onenigheid en geëxternaliseerd geheugen bestaan alle onder gevestigde terminologie (\secref{2.11}).

Voor het gemak verwijzen wij naar het hier ingenomen perspectief als \textbf{AI Sociology} — een werklabel voor een onderzoeksrichting, niet de naam van een voorgestelde discipline, en een label waarin wij geen prioriteit claimen: een gevestigde literatuur gebruikt \emph{sociology of artificial intelligence} al voor een ander object (\secref{2.8}). Wat hier bedoeld wordt is nauwer: hoe een beschreven organisatorische positie en een verklaarde herkomst van een bericht het gedrag veranderen, ongeacht of de beschreven structuur is geïmplementeerd.

De reden voor een werklabel in plaats van een discipline is inhoudelijk. De operationele beweringen hebben falsificatiecriteria; het vocabulaire waarin zij geïnterpreteerd worden niet. Of "sociale positie" iets beschrijft dat verder gaat dan een antropomorfe projectie, is op dit moment niet beslisbaar, en een discipline vereist dat het dat wel is.


\subsubsection*{Een slotnoot over motivatie}
\phantomsection
\label{sec:nl:x-een-slotnoot-over-motivatie}
\addcontentsline{toc}{subsubsection}{Een slotnoot over motivatie}

Eén zorg heeft dit werk gevormd en zij wordt vermeld als die van de Auteur, niet als bevinding.

Taalmodellen worden getraind op een corpus waarin de huidige consensus overweldigend vertegenwoordigd is. Zij zijn daarom goed geschikt om gevestigde wetenschap voort te zetten en slecht geschikt om de eerste stap daarbuiten te zetten: het statistische gewicht begunstigt het gevestigde verhaal. Een onderzoeksopstelling die uit zulke modellen is gebouwd loopt het risico die bias in versterkte vorm te erven — verscheidene deelnemers, één gedeelde prior, en de schijn van onafhankelijke overeenstemming.

De reactie die hier is nagestreefd is niet modellen overreden een onconventioneel standpunt te aanvaarden. Zij is de samenwerking zo inrichten dat elk standpunt, conventioneel of niet, door dezelfde procedure wordt onderzocht — en zo dat waar de procedure niet kan beslissen, zij dat zegt.

Of die opstelling werkt is onbeslist. De experimenten in \secref{11} zijn het middel om daarachter te komen, en de waarschijnlijkste enkele uitkomst is dat één niveau van de schaal verdwijnt.

Dat zou een resultaat zijn.

\end{dutch}
